%% 
%% Copyright 2007-2024 Elsevier Ltd
%% 
%% This file is part of the 'Elsarticle Bundle'.
%% ---------------------------------------------
%% 
%% It may be distributed under the conditions of the LaTeX Project Public
%% License, either version 1.3 of this license or (at your option) any
%% later version.  The latest version of this license is in
%%    http://www.latex-project.org/lppl.txt
%% and version 1.3 or later is part of all distributions of LaTeX
%% version 1999/12/01 or later.
%% 
%% The list of all files belonging to the 'Elsarticle Bundle' is
%% given in the file `manifest.txt'.
%% 
%% Template article for Elsevier's document class `elsarticle'
%% with harvard style bibliographic references

% \documentclass[preprint,12pt]{elsarticle}

%% Use the option review to obtain double line spacing
%% \documentclass[preprint,review,12pt]{elsarticle}

%% Use the options 1p,twocolumn; 3p; 3p,twocolumn; 5p; or 5p,twocolumn
%% for a journal layout:
% \documentclass[final,1p,times]{elsarticle}
% \documentclass[final,1p,times,twocolumn]{elsarticle}
\documentclass[final,3p,times]{elsarticle}
% \documentclass[final,3p,times,twocolumn]{elsarticle}
% \documentclass[final,5p,times]{elsarticle}
% \documentclass[final,5p,times,twocolumn]{elsarticle}

%% For including figures, graphicx.sty has been loaded in
%% elsarticle.cls. If you prefer to use the old commands
%% please give \usepackage{epsfig}

%% The amssymb package provides various useful mathematical symbols
\usepackage{amssymb}
%% The amsmath package provides various useful equation environments.
\usepackage{amsmath}
\usepackage{booktabs}
\usepackage{algpseudocode}
\usepackage{longtable}
\usepackage{tabularx}
\usepackage{siunitx}
%% The amsthm package provides extended theorem environments
%% \usepackage{amsthm}
\usepackage{bbm}
\usepackage{graphicx}
\usepackage{subcaption}
\usepackage{caption}
\usepackage[colorinlistoftodos]{todonotes}
% \usepackage[margin=2cm]{geometry}
\usepackage{booktabs}
\usepackage{enumitem}
\usepackage{amsfonts}
\usepackage{amssymb}
\usepackage{amsmath}
\usepackage{cancel}
\usepackage{mathtools}
\usepackage{bm}
\usepackage{esvect}
\usepackage{gensymb}
\usepackage{float}
% \usepackage{abstract}
\usepackage{tabulary}
\usepackage{nomencl}
\usepackage{multicol}
\usepackage[hyphens]{url}
\usepackage[hidelinks]{hyperref}
\usepackage[capitalize]{cleveref}

% \def\appendixname{}
% elsarticle embeds "Appendix" in \thesection, so we output the counter only
\crefname{appendix}{}{}
\Crefname{appendix}{}{}
\crefformat{appendix}{#2#1#3}
\Crefformat{appendix}{#2#1#3}

\journal{Applied Soft Computing}

\begin{document}

\begin{frontmatter}

\title{Multiobjective optimization-based design and dispatch of islanded, hybrid microgrids for remote, off-grid communities in sub-Saharan Africa} %% Article title

\author[1]{Vineet Jagadeesan Nair} %% Author name

%% Author affiliation
\affiliation[1]{
            organization={Department of Mechanical Engineering, Massachusetts Institute of Technology},
            addressline={77 Massachusetts Avenue}, 
            city={Cambridge},
            postcode={02141}, 
            state={MA},
            country={USA}
}

\begin{abstract}
%% Text of abstract
% Reliable, affordable electricity remains inaccessible to over 600 million people in sub-Saharan Africa (SSA), where islanded microgrids combining renewable sources with battery storage and diesel backup represent the most viable electrification pathway. This paper presents a multiobjective, multiperiod global optimization framework for the design, sizing, and dispatch of such islanded hybrid microgrids, with a detailed case study for a remote community in Kenya. System sizing is optimized over a one-year horizon and operational dispatch over a representative day, both using hourly resolution. The formulation simultaneously minimizes lifecycle levelized cost of energy (LCOE), emissions, lost load, and dumped energy, while maximizing renewable energy penetration. Among seven optimization algorithms benchmarked, particle swarm optimization (PSO) achieves the best combined runtime (63\,s) and solution quality (minimum normalized objective of 0.146), and is used for all subsequent sizing results. The optimal technology combination -- solar PV with wind turbines, lithium-ion battery storage (LI), and diesel engine (DE) backup -- achieves a normalized LCOE of 0.46\,\$/kWh with greater than 94\% renewable penetration, outperforming all alternative configurations. Pareto-optimal solution sets reveal key trade-offs among economic, environmental, and reliability objectives, demonstrating that cost-only optimization leads to suboptimal outcomes in emissions and reliability. Sensitivity analyses show that fuel prices and financial discount rates are the most influential parameters for SSA conditions. A break-even distance analysis confirms the microgrid is economically competitive with grid extension at the chosen location. The economic-environmental dispatch model produces day-ahead schedules that are generally robust to short-term uncertainties in load and renewable generation, although prolonged wind lulls significantly increase diesel reliance. This work addresses a critical gap in the literature: despite SSA hosting the largest concentration of people without electricity access, no prior study has provided a comprehensive multi-objective design and dispatch optimization tailored to this region's specific resource, economic, and operational conditions.

Optimal sizing and dispatch of islanded hybrid microgrids constitutes a challenging class of non-convex, non-smooth mixed-integer optimization problems, arising from discrete technology choices, continuous capacity variables, and an embedded 8,760-hour annual simulation. This paper presents a systematic benchmark of seven global metaheuristic and direct-search algorithms — particle swarm optimization (PSO), genetic algorithm (GA), simulated annealing, surrogate optimization, multiple-start search, global search, and direct pattern search — on a five-objective microgrid sizing problem that simultaneously minimizes levelized cost of energy (LCOE), emissions, power supply deficiency, and dumped energy while maximizing renewable penetration. PSO achieves the best combined runtime (63 s) and solution quality (normalized objective 0.146), converging 4.6× faster than GA and 35× faster than surrogate optimization while matching or exceeding global solution quality across all trials. The benchmark demonstrates that population-based swarm and evolutionary methods consistently outperform gradient-based and surrogate approaches on this non-smooth, mixed-integer landscape, providing practical solver selection guidance for this problem class. The PSO-based framework is applied to the design and dispatch of an islanded hybrid microgrid for a remote community in sub-Saharan Africa (SSA) — a region where resource economics, high fuel prices, and elevated discount rates impose conditions not addressed in prior multi-objective work. The optimal system — solar PV, wind turbines, lithium-ion battery storage, and diesel engine backup — achieves 94\% renewable penetration and reduces LCOE by over 50\% relative to diesel-only generation. Pareto-optimal solution sets quantify trade-offs among economic, environmental, and reliability objectives, while parametric sensitivity analyses and robustness testing under renewable uncertainty characterize design behavior under SSA-specific conditions.
\end{abstract}

%%Research highlights
\begin{highlights}
\item Particle swarm global optimization is best suited to multiobjective microgrid design.
\item A hybrid system with solar, wind, Li-ion batteries, and diesel engines is optimal.
\item Levelized cost is the primary driver for allocating capacity between wind and solar.
\item Accurate sizing and sensitivity analyses can help reduce overbuild and overall costs.
\item Battery and backup generator dispatch enable resilience to short-term disturbances.
\end{highlights}

%% Keywords
\begin{keyword}
%% keywords here, in the form: keyword \sep keyword

%% PACS codes here, in the form: \PACS code \sep code

%% MSC codes here, in the form: \MSC code \sep code
%% or \MSC[2008] code \sep code (2000 is the default)

Microgrids, multiobjective optimization, hybrid energy systems, sensitivity analysis, renewables, batteries, energy access

\end{keyword}

\end{frontmatter}

\section*{Nomenclature}

\begin{tabular}{ll}
\multicolumn{2}{l}{\textit{Acronyms}} \\
BS  & Battery storage system \\
CRF & Capital recovery factor \\
DE  & Diesel engine \\
DG  & Distributed generator (backup) \\
DPSP & Deficiency of power supply probability \\
DR  & Demand response \\
DSM & Demand-side management \\
EED & Economic and environmental dispatch \\
GA  & Genetic algorithm \\
LA  & Lead-acid battery \\
LCOE & Levelized cost of electricity (\$/kWh) \\
LI  & Lithium-ion battery \\
MG  & Microgrid \\
MILP & Mixed-integer linear program \\
MPC & Model predictive control \\
MT  & Microturbine \\
PSO & Particle swarm optimization \\
PV  & Solar photovoltaic array \\
REF & Renewable energy fraction (penetration) \\
REPG & Relative excess power generated (dump fraction) \\
RES & Renewable energy sources \\
SA  & Simulated annealing \\
SOC & State of charge \\
SSA & Sub-Saharan Africa \\
TAC & Total annualized cost (\$/yr) \\
WT  & Wind turbine \\[6pt]
\multicolumn{2}{l}{\textit{Key variables}} \\
$E_{b,\text{init}}$ & Initial rated battery energy capacity (kWh) \\
$n_s$ & Number of solar PV modules \\
$n_w$ & Number of wind turbines \\
$P_{BS}(t)$ & Battery power output at time $t$ (kW) \\
$P_{DG}(t)$ & Distributed generator power output at time $t$ (kW) \\
$P_{RES}(t)$ & Total renewable power output at time $t$ (kW) \\
$w_i$ & Weight on objective $i$ \\
\end{tabular}

\section{Introduction and problem background \label{sec:access}}

Over 1 billion people or 13\% of the world's population still lacks access to electricity, with most of these people living in Africa, India and other parts of southeast Asia. In particular, sub-Saharan Africa (SSA) suffers from low electrification rates and hosts over 621 million people (i.e., approximately two in every three persons in Africa) who currently do not have reliable electrical power \cite{SSA_decent}. Up to 85\% of the unserved population lives in remote rural villages \cite{quak_CBA}, located far away from traditional 'macrogrids' that rely on centralized thermal power generation dominated by fossil fuels like coal and natural gas. Of the remaining 15\% that live closer to the grid, many cannot afford to pay the exorbitant connection fees or metering and wiring costs \cite{quak_CBA}. Electricity access is also related to various other socioeconomic issues as well as the health and well-being of these populations. Improving electricity provision is highly correlated with poverty alleviation through improved living standards, increased productivity and better employment opportunities. This will also help achieve several sustainable development goals (SDGs) set by the United Nations (UN) \cite{SSA_decent}.

It is challenging to connect such isolated communities to the conventional grid due to the high costs of constructing transmission and distribution (T\&D) infrastructure such as high voltage AC and DC lines. According to Power for All, grid extension projects and infrastructure upgrades in sub-Saharan Africa can cost up to \$5.5 million per MW/km \cite{grid_gtm}. Long planning, procurement, environmental assessment, and construction phases result in budget overruns and long delays in project completion. This implies that connection costs are extremely high across the continent, averaging around \$136 (for single-phase connections), and can be significantly higher than monthly household incomes. This makes such projects unprofitable, especially for sparsely populated rural areas suffering from high marginal costs and less concentrated electricity demand. As a result, significant segments of the global population either live without any electricity access at all or lack reliable and continuous access to power. Even when they do have access, most communities suffer from 'energy poverty' whereby they spend significant portions of their income on energy-related expenses i.e., electricity and fuel consumption. Furthermore, close to half of the global population still relies on biomass fuels like firewood or charcoal for cooking and heating, and kerosene for lighting which are expensive and have adverse impacts on human health \cite{kammen}. This is especially the case in developing and less developed countries, with almost two-thirds of the world's poorest people residing in rural areas \cite{kammen}. Thus, there is a need to come up with more affordable and reliable solutions for rural electrification that can be implemented rapidly in the near future. 

\subsection{Microgrids and their role in rural electrification}

Microgrids are low-voltage, localized hubs of electricity generation, storage and utilization that remain synchronous with the larger, traditional grid under normal operation, but can disconnect at the common point of coupling and operate independently if needed. Such decentralized systems have become more popular and practical due to the increased penetration of distributed generation and power production from renewable energy sources (RES) like wind, solar, geothermal, biomass and small-scale hydroelectric plants. Microgrids can operate either in:
\begin{enumerate}
    \item \textit{Grid-connected} mode: This allows it to exchange power with the macrogrid in both directions (under certain transmission constraints), in order to balance supply and demand. In this configuration, the microgrid operator can generate additional revenue and profits by selling excess clean power from RES back to the main grid.
    
    \item \textit{Islanded} mode: Here, the microgrid system is completely isolated and not connected to the larger grid network. Since no external power exchange is allowed, all power demand must be met locally and any excess generation must be curtailed by grounding to earth or using dump loads.
\end{enumerate}

Due to challenges with grid extension described in \cref{sec:access}, this work only considers stand-alone, off-grid systems that always operate in islanded mode. This choice introduces some additional special challenges that will be discussed in greater detail in \cref{sec:MG_challenges}. 

\subsection{Advantages of microgrids}

In comparison to large-scale conventional sources such as thermal and nuclear baseload power plants used in legacy grids, microgrids have much smaller capacities but offer greater modularity and flexibility. Thus, they facilitate better integration of renewable, distributed energy resources (DERs) that are highly intermittent, variable in supply and often located far away from load centers. The unpredictability of renewables also leads to challenges and compatibility issues with the conventional grid such as poor power quality, unsteady supply voltages, undesirable harmonics and a high degree of vulnerability to disturbances like natural disasters, extreme weather, terrorism and other cyber-physical attacks. The interconnectedness and sensitivity of such centralized systems can result in cascading failures and domino effects, where a fault in one small section of the grid affects everyone else to some extent. This is evident through blackouts, brownouts and load-shedding that often occur in urban areas, especially in developing nations \cite{macrogrid_prob}. 

By reducing our reliance on fossil fuels, microgrids can help decrease greenhouse gas emissions and mitigate climate change. They will play a critical role in transitioning to a sustainable future based on clean energy sources and achieving ambitious goals such as the SDGs and Paris Climate Agreement. By allowing islanded operation, they can simply disconnect from the macrogrid in case of failures or emergency to still supply local load, thus creating more resilient and secure power networks that are robust to disturbances. Well designed, controllable microgrids can also enable grid optimization by relieving congestion and providing various ancillary services to the main grid as needed, e.g., load shifting, frequency or voltage control and peak shaving. Moving generation sites closer to demand will help reduce transmission and distribution (T\&D) losses while allowing for T\&D investment deferral. Furthermore, microgrids are one of the only realistic options to meet the need for low-cost, locally available energy supply in other remote, resource-constrained locations apart from rural areas as well, such as islands, military bases, harsh environments like the Arctic. Overall, microgrids promise to democratize cost-effective electricity access globally, with the International Energy Agency (IEA) estimating that about 70\% of rural areas that currently lack access will be connected using mini-grid, microgrid or off-grid solutions \cite{WEO_2018,mentis}. This means $\approx$ 140 million rural Africans could get access to electricity through the creation of 100,000 mini-grids by 2040 \cite{wri}.

\subsection{Challenges with microgrids \label{sec:MG_challenges}}

Despite the numerous merits of microgrids, they still face some serious practical challenges that hinder widespread implementation. Owing to their smaller size, they lack the averaging effects that help balance supply and demand across different locations and nodes in larger networks. This can lead to stability issues as well as unmet load in case of excess demand or insufficient RES generation due to unfavorable weather conditions e.g., extended periods of wind lulls, cloudy or rainy days. Load fluctuations occurring in microgrids can be more extreme due to their smaller capacities. Larger, conventional power systems have transmission spanning across multiple regions with distinct weather patterns and load profiles, which helps smooth out variations over time. In addition to highly stochastic demand, microgrids must balance this with the unpredictability and intermittency of renewables. In order to accommodate for the worst-case scenarios, both RES and storage capacity often need to be oversized which raises project costs. 

These issues are further intensified in the case of isolated microgrids since these do not have the main grid (e.g., backup peaker plants and interconnectors) as a safety net to fall back on. They also lack large fossil-fueled synchronous generators that offer high inertia and oppose sudden changes in grid frequency or voltage. Thus, decisions need to be made over faster timescales in order to rectify supply-demand mismatches in both real and reactive power to ensure frequency regulation and voltage stability, respectively \cite{tech_challenges}. This study focuses only on power balancing at slower time-scales on the order of months, days and hours. Frequency and voltage stability are not considered here since these require faster control actions and responsiveness on the order of minutes and seconds. These are primarily achieved through the design of appropriate switches and other power electronics equipment, not through the optimization and control methods which are the topic of the current project.

\section{Literature review}

There is a large body of existing research in the field of microgrid design and operation. Past studies range from optimal selection and sizing of generation and storage technologies to optimizing the energy management and dispatch of these components in real-time. 

\subsection{Sizing optimization}

Once the most suitable technology types are selected based on a combination of economic, technical and socio-political factors, the next step is to decide the optimal sizes of these sub-systems. For RES such as wind turbines (WT) and solar photovoltaic panels (PV), this involves setting the power ratings of each individual unit as well as the total number of units needed (e.g., no. of turbines, total surface area of PV array). For energy storage technologies like batteries and pumped hydro schemes, this would mean fixing their rated powers as well as total energy capacity. Finally, for standby diesel engines or gas turbines, their total rated power is the primary decision variable. Most studies consider \textit{hybrid} microgrid designs for standalone power systems i.e., integrated systems combining two or more renewable sources as primary generation with conventional distributed generators (DG) based on fossil fuels, as secondary, backup systems to meet reliability requirements \cite{kaabeche2014} \cite{bhandari2014}. A techno-economic, multi-objective optimization is performed to minimize pre-determined objectives including economic costs \cite{tazvinga2013}, pollutant emissions (to reduce negative environmental impacts) and other performance metrics for power quality, system reliability and efficiency \cite{kaabeche2011a}. These reflect various goals of the community that the microgrid serves. A more detailed description of individual cost terms will be provided in \cref{chap:methods}.
 
Past studies have shed interesting insights into trade-offs that occur during MG design. For instance, it has been shown that PV/WT systems that use both diesel generators and batteries are more economically viable than those using either only batteries or diesel engines alone \cite{kaabeche2014}. The same study also found that system configurations with much lower number of wind turbines relative to solar PV modules have the lowest energy costs and net equivalent $CO_2$ emissions. This resulted from the strong solar potential in the region considered in Algeria \cite{kaabeche2014} which is also likely to be the case for the location of Timbila, Kenya chosen here. This is an example of how location-specific the planning and design process is owing to the regional characteristics of climatic conditions, availability of wind and solar energy as well as electricity demand profiles. The sizing results also vary depending on the choice of objectives and financial parameter values (like interest, discount and inflation rates) used in the analysis since capital costs tend to be significant for renewables and storage while running fuel costs are the major expense for backup generation. Furthermore, different objectives often compete with one another since more reliable, cleaner systems often imply higher project costs. It has also been found that in order to attain the same reliability as the usual grid, \textit{off-grid} systems can be up to ten times more expensive than \textit{on-grid} solutions using solar PV with batteries \cite{weng2016}.

Overall mathematical models for hybrid microgrids are constructed by combining fundamental models for individual sub-systems. Well-established, proven models already exist for each of the wind, solar, diesel, gas and battery bank systems included in this study \cite{bhandari2014}. In particular, studies have found that the sizing of energy storage and backup generators often plays the crucial role in determining system performance. For example, the efficiency of microgrid operation depends strongly on the battery scheduling due to the key role it plays in peak load shifting and stabilizing fluctuations in both demand and RES supply \cite{gu2014}. Most studies that formulate their own optimization models choose to simulate the problem over a 1-year period (i.e., 8760 h), resulting in excellent convergence to feasible optimum solutions \cite{logenthiran2010}. However, it is possible to simulate over longer periods using existing commercial software as well. There are several such commercial and open-source microgrid design software tools available in the industry; a detailed overview of the main tools is provided in \cref{app:design_tools}

\subsection{Optimal dispatch}

This area of research studies methods for both:
\begin{itemize}

    \item Unit commitment: Decisions regarding the production scheduling of different power generation units to meet forecasted future energy demand at minimum cost. This determines which dispatchable units (including both generators and storage devices) are online and producing energy versus being offline at any given time instant \cite{saravanan2013}.
    
    \item Unit dispatch: These are production decisions that set exactly how much power each online unit is generating in a specific time interval.
    
\end{itemize}

Whereas single-objective optimizations generally minimize purely economic costs, multiobjective problem formulations are becoming increasingly more common. These simultaneously optimize with respect to several objectives such as emissions, cost, power loss and heat \cite{pathak2017}. However, conventional optimization methods are not well-suited to such combined emissions-economic dispatch problem since they are generally nonlinear, non-smooth and non-convex with severe equality and inequality constraints \cite{pathak2017}. There is also the issue of convergence to local minima that may not be globally optimal solutions. Furthermore, these mixed integer programming problems generally involve both discrete as well as continuous, multi-index decisions variables and constraints with high dimensionality. This precludes the use of popular solver packages like CVX which only work for disciplined convex programs, and have led to the development of metaheuristic based optimization approaches. Instead of enforcing hard constraints, some studies use soft constraints with large violation penalties to introduce slackness and ensure feasibility of the optimization problem \cite{zachar2016}. 

Some focus areas of past microgrid operations research include mechanisms for seamless, automatic islanding and reconnection to the grid when necessary to improve reliability and robustness \cite{mg_review}. Passive primary and secondary control strategies have also been explored to regulate both grid supply frequency and voltage, respectively. Furthermore, unlike larger grid networks, microgrids require a third, higher level of control for micro-economic optimization \cite{mg_review}. Modeling and optimization in microgrids can also borrow insights, tools and methods from related applications including electric vehicles as controllable loads, feature-dependent power generation (e.g., solar PV) which depends on environmental factors often outside the designer's control, controllable generation (feature-independent feeds) and existing models for energy storage systems \cite{mg_mod_opt}. Some studies have also used popular programming tools like MATLAB \texttt{Simulink} to create and simulate simplified, dynamic models of small-scale microgrids.

In addition to geographical and spatial factors, temporal variations are critical, including both long (seasonal, monthly etc.) and short-term (daily, hourly etc.) patterns in load and RES supply. In particular, these determine the optimal cycling strategies for the battery and operational time of the backup DG, which in turn affect fuel consumption and operating costs \cite{tazvinga2013}. Optimal configuration and operation of microgrids have been shown to improve power quality, reliability, security, sustainability and decrease electricity costs to consumers \cite{pathak2017}. In case of islanded and isolated microgrids lacking macrogrid support, literature indicates that incentive-based demand response (DR) programs \cite{nwulu2017} and/or model predictive control \cite{zachar2016} are essential to reject forecasting errors and noise in RES generation and load, improve optimality of both supply and demand and thus reduce costs. In terms of DR implementation and demand side management (DSM), it has been shown using game-theoretic approaches that globally optimal aggregate demand profiles can be achieved by incentivizing individual agents to locally minimize costs, through price signals \cite{paola2017}. Finally, several recent simulations have also focused on using \textit{integrated design and operation optimization}, a combined approach where both system sizing and dispatch are optimized simultaneously in a single stage rather than sequentially in two distinct steps \cite{moshi2016}.

\subsection{Survey of optimization algorithms}

Many studies formulate sizing and dispatch as non-linear mixed-integer programs solved with global heuristic methods; others use MILP relaxations for computational tractability \cite{parisio2014, Qu2018}. The seven algorithms benchmarked in this work (\cref{sec:results}) are summarized below; full derivations and parameter details are in \cref{app:opt_alg}. Genetic algorithms (GA) and evolutionary algorithms (EA) are stochastic, population-based methods that search the solution space by mimicking biological evolution through selection, mutation, and crossover of candidate solutions \cite{nemati2015}. They require no gradient information, handle integer and mixed-integer variables natively, and can generate full Pareto-optimal solution sets in a single run \cite{fadaee2012}. Their main drawback is that performance depends heavily on tuning of population size, crossover rate, and mutation rate, and complexity grows with the number of decision variables \cite{fathima2015}. Particle swarm optimization (PSO) simulates the social behavior of a swarm: each particle moves through the search space guided by its own best-known position and the swarm's global best, updating velocity at each iteration \cite{fathima2015, pathak2017}. It is generally easier to implement than GA with fewer hyperparameters, converges quickly on non-convex problems, and has been applied successfully to microgrid sizing and economic dispatch \cite{alrashidi2010}. It can struggle with equality constraints and may stagnate in flat regions of the objective landscape. The velocity-update equations are given in \cref{app:pso}.

Surrogate optimization (SO) constructs a radial-basis-function interpolant of the objective using quasirandom sample points, then minimizes the surrogate rather than the expensive true function \cite{gutmann2001}. This is advantageous for computationally costly evaluations, and the method is proven to converge to the global optimum on bounded domains. In practice, however, it is slow because a large number of evaluations are needed to build an accurate surrogate, and it lacks an explicit stopping criterion based on optimality. Simulated annealing (SA) is a single-solution method that probabilistically accepts worse solutions with a temperature-dependent probability, allowing escape from local optima analogously to the physical annealing of metals \cite{fathima2015}. It is theoretically guaranteed to converge globally under a logarithmic cooling schedule, but this schedule is extremely slow in practice. SA also does not support parallel computation. Multiple start search (MS), Global search (GS), and Direct pattern search (DPS) are gradient-based or direct-search local methods run from multiple starting points to approximate global solutions. MS and GS use local gradient solvers (\texttt{active-set}) launched from many random starts, which makes them faster than population methods on smooth landscapes but less reliable on non-smooth or flat objectives. DPS uses a mesh-based direct search with no gradient information, which is more robust to non-smoothness but considerably slower due to fine mesh resolution requirements.

Recent work extends these methods in three directions relevant here. Improved multiobjective variants of evolutionary and swarm algorithms return Pareto sets rather than single designs, for example an improved NSGA-II for the dispatch of hybrid ship power systems \cite{wang2024nsga} and an improved multiobjective grey wolf optimizer for sizing an island microgrid \cite{zhu2020imogwo}; grey wolf optimization has also been compared with PSO for wind--solar sizing \cite{geleta2022gwo}, and hybrids of PSO with other heuristics aim to balance exploration and exploitation \cite{javidrad2017psosa}. Several studies pair multiobjective search with multi-criteria decision making to choose one design from the Pareto set, such as multiobjective differential evolution followed by AHP--TOPSIS for stand-alone PV sizing \cite{muhsen2019mode} and weighted-sum Pareto sets ranked with ELECTRE for microgrid scheduling \cite{karimi2019madm}. Resource and load uncertainty is also increasingly propagated into the design stage, through Markov chain Monte Carlo scenario generation for stand-alone microgrid planning \cite{bakhtiari2021mcmc} or two-stage stochastic programming for off-grid hybrid systems \cite{calabrese2026tssp}. For stochastic solvers, the computational-intelligence literature recommends comparing result distributions over repeated independent runs with non-parametric tests rather than single runs \cite{derrac2011tutorial,carrasco2020trends}; we follow that practice in \cref{sec:results}.

\subsection{Control strategies and design tools \label{sec:control_survey}\label{app:design_tools}}
Real-time microgrid control (hierarchical primary, secondary and tertiary levels; centralized, decentralized and distributed architectures; model predictive control) lies outside the scope of this paper and is surveyed in \cref{app:control}. Widely used design tools (HOMER Pro, DER-CAM, REopt and the Microgrid Design Toolkit) are compared in \cref{app:tools}.

\subsection{Novelties in current study and contributions}

Three gaps motivate this work. First, most multi-objective microgrid studies treat grid-connected systems that island only under faults; permanently isolated systems, which must meet all load from local resources, receive less attention. Second, comparisons of stochastic solvers for microgrid sizing often rest on one run per solver, so differences between solvers cannot be separated from run-to-run variation. Third, studies that compute Pareto sets often stop short of a rule for choosing one design from them, and robustness is commonly checked with a few deterministic scenarios rather than by propagating resource and load uncertainty.

Sub-Saharan African conditions shape the parameter ranges and the design questions in this study. Diesel prices in remote areas are high and volatile (our default of 3.20 USD per US gallon, about 0.85 USD/L, is conservative for remote sites, and \cref{sec:sizing_sensitivity} spans 0.2--12 USD/gal); nominal interest and inflation rates are high (9\% and 5.7\% for Kenya in the base case), which penalizes capital-intensive storage; measured climate data are scarce, so we draw on the TAHMO station network and check it against reanalysis (\cref{app:climate}); and grid extension is expensive, which motivates the break-even distance analysis in \cref{sec:results}.

This paper makes four contributions. (i) We formulate the sizing of an islanded PV--wind--battery--diesel microgrid as a simulation-based, non-smooth optimization problem in which an 8,760-hour rule-based operating simulation scores each candidate on five normalized objectives, and we state how integrality, normalization and reliability enter the formulation. (ii) We benchmark seven global metaheuristic and direct-search solvers and a PSO--pattern-search hybrid on this problem at an equal budget of 2,000 evaluations over 30 independent runs each, using success rates against the integer optimum found by exhaustive enumeration, expected running time and non-parametric tests. (iii) We pair a-priori weighted-sum sizing with a-posteriori Pareto search, select designs from the Pareto set with knee-point, compromise and TOPSIS rules, and re-rank the candidates by expected cost, conditional value-at-risk of cost and loss-of-load probability over 500 Monte Carlo resource and load scenarios. (iv) We apply this workflow to an off-grid community in Timbila, Kenya, with fuel price, interest, inflation and storage cost ranges drawn from sub-Saharan African conditions, and identify which design rules carry over to other sites once local inputs are recalibrated. \cref{tab:litcomp} places the work relative to representative prior studies.

\begin{table}[htbp]
\centering
\footnotesize
\setlength{\tabcolsep}{3pt}
\caption{\label{tab:litcomp} Comparison of representative prior works on multi-objective hybrid microgrid optimization. \textbf{Obj}: objectives (C = cost, E = emissions, R = reliability, D = dump/curtailment; multi = several, not detailed in the abstract). \textbf{Isl.}: islanded/off-grid. \textbf{Size}/\textbf{Disp.}: sizing/dispatch performed. \textbf{R/S}: robustness or sensitivity analysis. \textbf{Sel.}: explicit rule for selecting a design from a Pareto set. \textbf{UQ}: probabilistic uncertainty propagation. \textbf{Stats}: repeated-run statistical comparison of solvers. \textbf{SSA}: sub-Saharan Africa case.}
\begin{tabular}{>{\raggedright\arraybackslash}p{2.8cm} >{\raggedright\arraybackslash}p{1.1cm} >{\raggedright\arraybackslash}p{2.9cm} cccccccc}
\toprule
\textbf{Study} & \textbf{Obj.} & \textbf{Method} & \textbf{Isl.} & \textbf{Size} & \textbf{Disp.} & \textbf{R/S} & \textbf{Sel.} & \textbf{UQ} & \textbf{Stats} & \textbf{SSA} \\
\midrule
Kaabeche et al.~\cite{kaabeche2014} & C, E & Iterative & \checkmark & \checkmark & & & & & & \\
Bhandari et al.~\cite{bhandari2014} & C & HOMER & & \checkmark & & & & & & \\
Tazvinga et al.~\cite{tazvinga2013} & C, E & GA & \checkmark & \checkmark & \checkmark & & & & & \checkmark \\
Logenthiran et al.~\cite{logenthiran2010} & C & GA & \checkmark & \checkmark & \checkmark & & & & & \\
Moshi et al.~\cite{moshi2016} & C & GA & \checkmark & \checkmark & \checkmark & & & & & \checkmark \\
Zachar \& Daoutidis~\cite{zachar2016} & C & MPC & & & \checkmark & \checkmark & & & & \\
Parisio et al.~\cite{parisio2014} & C & MILP & \checkmark & & \checkmark & & & & & \\
Nwulu \& Xia~\cite{nwulu2017} & C, E & DR/MPC & & & \checkmark & & & & & \\
Weng et al.~\cite{weng2016} & C & DER-CAM & & \checkmark & & & & & & \\
Gu et al.~\cite{gu2014} & C & MPC & \checkmark & & \checkmark & & & & & \\
Karimi \& Jadid~\cite{karimi2019madm} & multi & Weighted sum + ELECTRE & & & \checkmark & & \checkmark & & & \\
Muhsen et al.~\cite{muhsen2019mode} & C, R & MODE + AHP--TOPSIS & \checkmark & \checkmark & & & \checkmark & & & \\
Zhu et al.~\cite{zhu2020imogwo} & C, R & Improved MOGWO & \checkmark & \checkmark & & & & & & \\
Bakhtiari et al.~\cite{bakhtiari2021mcmc} & multi & MCMC scenarios & \checkmark & \checkmark & & & & \checkmark & & \\
Wang et al.~\cite{wang2024nsga} & multi & Improved NSGA-II & \checkmark & & \checkmark & & & & & \\
Calabrese et al.~\cite{calabrese2026tssp} & C, E & Two-stage SP, $\varepsilon$-constraint & \checkmark & \checkmark & \checkmark & \checkmark & & \checkmark & & \\
\textbf{This work} & C, E, R, D & Seven solvers + PSO--PS; \texttt{paretosearch}; TOPSIS & \checkmark & \checkmark & \checkmark & \checkmark & \checkmark & \checkmark & \checkmark & \checkmark \\
\bottomrule
\end{tabular}
\end{table}

\section{Microgrid system description \label{chap:methods}}

\begin{figure}[htbp]
  \centering
  \includegraphics[width=0.7\linewidth]{system_diagram_v2.png}
  \caption{Schematic illustrating key components and power flows in the hybrid microgrid.}
  \label{fig:chap3_MG}
\end{figure}

The microgrid consists of the following key components as shown in \cref{fig:chap3_MG}:

\begin{itemize}
    \item Decentralized renewable power generation sources (RES) that are non-controllable, non-deterministic devices: Solar photovoltaic arrays (PV) and Wind turbine generators (WT)
    \item Standby, fossil fuel distributed generators (DG) for backup which are controllable, deterministic units: Diesel engines (DE) and Microturbines (MT)
    \item Battery energy storage (BS): Lithium-ion ($Li^{+}$) (LI) and Lead-acid (LA) batteries
    \item Charge controllers to regulate the state of charge (SOC) and charging and discharging of batteries
    \item Power conditioners to regulate and improve the quality of power delivered from RES to the DC bus
    \item Bidirectional power converters that can function either as an inverter (DC-to-AC conversion) or rectifier (AC-to-DC conversion) depending on the system requirement
\end{itemize}

This is a hybrid microgrid since it consists of both conventional and renewable generation sources. It is an isolated system in islanded mode serving a remote, rural community implying that there are no connections or power transfers with any larger, external macrogrid. Thus, the MG can neither sell excess power to nor purchase power from, the main grid. Detailed physical models of all elements in the MG can be found in \cref{app:mg_models}, including information about climate data inputs (for a specific location in Kenya) and load characteristics. 

\section{Formulation of optimization problem \label{sec:opt_prob}}

\textbf{Two-stage sequential framework.} The overall optimization is solved in two stages (\cref{fig:workflow}). In Stage~1 (sizing), an annual hourly simulation with a rule-based dispatch inner loop is embedded in an outer global optimization loop that chooses the PV, WT and battery capacities. In Stage~2 (dispatch), the sized system is taken as given and a multi-objective economic-environmental dispatch problem is solved for a representative day. The Stage-1 controller is non-anticipative: it uses only the current hour's load, renewable output and state of charge, and it follows the load-following logic of most deployed off-grid controllers and of HOMER Pro's default strategy. Each candidate design is therefore scored under an operating policy that an operator can implement without forecasts. Embedding a 24-hour dispatch optimization inside every sizing evaluation would multiply the cost of each of the thousands of evaluations per run by 365 daily solves. The decoupling has a price: the design ranked best under rule-based operation need not be best under optimized dispatch. \cref{sec:limitations} discusses this limitation, and \cref{sec:dispatch_opt} reports how optimized dispatch changes daily cost and emissions for the selected design.

A discrete-time, iterative optimization is performed involving two components:

\begin{enumerate}
    \item \textbf{Design and sizing optimization (Stage 1)}: Over a time horizon of 1 year (8760 h) using time steps of 1 h each ($\Delta t = 1\;h$). Three variable inputs in this problem are the number of solar PV modules $n_s$, the number of wind turbines ($n_w$) and the initial rated battery capacity ($E_{b,\;init}$). The sizing is done to minimize several objectives including levelized lifecycle costs, emissions, dumped energy, and lost load. Within the PSO outer loop, the annual hourly simulation uses a priority-based rule dispatch (\cref{sec:disp_strat}) rather than an inner optimization, keeping each function evaluation computationally fast. One evaluation of the 8,760-hour simulation takes about 2\,ms on a current laptop, so a sizing run of 2,000 evaluations completes in seconds (\cref{sec:results}).
    \item \textbf{Day-ahead dispatch optimization (Stage 2)}: This is an economic and environmental/emissions dispatch (EED) over a time horizon of 1 day (24 h) using hourly time steps. This is conducted by taking as given the optimal sizes and capacities outputted by the design/sizing optimization above and then optimizing with respect to $P_{DG}(t)$ and $P_{BS}(t)$ as variable inputs. Thus, it optimizes the scheduling and operation of dispatchable units (DG and BS) while renewable power output $P_{RES}$ is decided by climatic conditions and thus non-dispatchable.
\end{enumerate}

The multiobjective optimization problem considers five distinct objectives: (i) Levelized cost of electricity (LCOE), (ii) Emissions (Em), (iii) Deficiency of power supply probability (DPSP), (iv) Dumped or Relative excess or power generated (REPG), and (v) Renewable energy fraction (REF). For brevity, here we only describe the overall optimization problem and constraints. Further details on all costs and objective functions can be found in \cref{app:costs}. 

\subsection{Dispatch strategy \label{sec:disp_strat}}
For the sizing and design optimization, a priority order for scheduling the generation and storage units was encoded directly into the objective function. The load-following dispatch strategy assumed is described in \cref{tab:dispatch_strat_for_sizing}.
\begin{table}[htbp] 
\centering
\begin{tabular}{p{0.75cm} p{13cm}}
\toprule
\multicolumn{1}{c}{\textbf{Scenario}} & \multicolumn{1}{c}{\textbf{Action}}    \\
\midrule
1 & Use RES power output to meet all load. \\ [1ex]
2 & If generation RES (PV and WT) exceeds power demand, use it to charge the storage devices (BS). \\ [1ex]
3 & If the excess RES supply exceeds the capacity and safe operating limits of the battery (i.e., $SOC = SOC_{max}$), then send the remaining excess power generated to dump loads or ground. \\ [1ex]
4 & If RES generation is insufficient to meet demand, then discharge the BS. \\ [1ex]
5 & If the BS reaches its lower limit or $DoD$ (i.e., $SOC = SOC_{min}$) or its power rating and some load is still unsatisfied, then dispatch the backup DG until either all the demand has been met or it has been ramped up to its maximum (rated) power, whichever occurs first. \\ [1ex]
6 & There is a lower limit on the power output at which the DG can be operated when online. If the remaining excess load in step 5 above is less than this value ($P_{DG,\;min}$), then the extra DG output is rectified from AC to DC and then either used to charge the BS or dumped, depending on the SOC of the BS (i.e., go to step 3). \\ [1ex]
7 & At this point, if there is still some unmet demand, then this is recorded as lost load and contributes to DPSP. An alternative option is to use some form of incentive-based demand response (DR) or demand-side management (DSM) and thus encourage consumers to either shift or voluntarily curtail their loads, instead of resorting to load-shedding. This possibility will be discussed in more detail in \cref{sec:dr}. \\ [1ex]
\bottomrule
\end{tabular}
\caption{Summary of prescribed load-following strategy used for sizing optimization.}
\label{tab:dispatch_strat_for_sizing} 
\end{table}

\subsection{Constraints}

\begin{table}[htbp] 
\centering
\begin{tabular}{p{5cm} p{10cm}}
\toprule
\multicolumn{1}{c}{\textbf{Constraint}} & \multicolumn{1}{c}{\textbf{Description}}    \\
\midrule
Power limits of the DG & Whenever the DG is online (i.e., when $P_{DG}(t) > 0$), its power output must lie between the rated minimum and maximum powers specified. \\ [1ex]
Positivity constraints & The power output of the DG must be non-negative at all times. \\ [1ex]
Ramp rate limits on the DG & This determines how quickly it can be ramped up to, or ramped down from the rated power. However, NERC (North American Electric Reliability Corporation) disturbance control standards specify that regulating and supplemental reserves must be able to reach their full name-plate capacity within 10-20 minutes \cite{nerc}. Thus, it did not make sense to try and enforce this constraint in the current model since an hourly time step is used. \\ [1ex]
Power limits on the BS & Both the charging (-ve) and discharging (+ve) powers of the BS must not exceed its maximum power rating. \\ [1ex]
Safe battery operation & Upper and lower limits on the $SOC$ of the BS. \\ [1ex]
Reliability requirement & Upper limit on the DPSP for the MG to be reliable and ensure power balance. \\ [1ex]
\bottomrule
\end{tabular}
\caption{Summary of MG constraints.}
\label{tab:dispatch_constraints} 
\end{table}

The main constraints in this problem are listed in \cref{tab:dispatch_constraints}. Thus, the multiobjective, weighted, and constrained sizing optimization problem can be written as:
\begin{equation}
    \label{eq:sizing}
    \min_{n_s, \; n_w, \; E_{b,\;init}} w_1 \cdot \frac{LCOE}{LCOE_{base}} \; + \; w_2 \cdot \frac{Em}{Em_{base}} \; + \; w_3 \cdot DPSP \; + w_4 \; \cdot REPG \; + \; w_5 \cdot (1-REF)
\end{equation}
where all the weights sum up to unity i.e., $w_1 + w_2 + w_3 + w_4 + w_5 = 1$ and is subject to the following non-negativity constraints on the inputs $n_s, \; n_w, \; E_{b,\; init} \geq 0$. All the remaining constraints are directly coded into the MATLAB program that calculates the objective function through the prescribed dispatch strategy detailed in \cref{sec:disp_strat}. Thus, these are already enforced implicitly.

Eq.~(\ref{eq:sizing}) is a weighted-sum scalarization of the five objectives. For strictly positive weights its minimizer is Pareto-optimal, but the method can only reach supported Pareto points, those on the convex hull of the Pareto front. We therefore use it as an a-priori method for a decision-maker with stated preferences (equal weights by default; \cref{app:weights} varies them) and complement it with the a-posteriori Pareto analysis of \cref{sec:pareto_opt,sec:pareto_select}. Reliability enters through the DPSP objective. With the default 16~kW DG the rule-based controller meets all load, so DPSP is zero for every candidate design; \cref{sec:pareto_rel} therefore also treats the DG rating as a decision variable to expose the cost--reliability trade-off.

Similarly, the optimization problem for day-ahead dispatch can be written as:
\begin{equation}
    \label{eq:dispatch}
    \min_{P_{DG}, \; P_{BS}} w_1 \cdot \frac{COE}{COE_{base}} \; + \; w_2 \cdot \frac{Em}{Em_{base}} \; + w_3 \; \cdot REPG \; + \; w_4 \cdot (1-REF)
\end{equation}
subject to the following constraints:
\begin{gather*}
\vv{0}_{1\times24} \; \leq \; \vv{\bm{P_{DG}}}(t) \; \leq \; P_{DG,\;r}\cdot \vv{\mathbbm{1}}_{1\times24}, \quad P_{DG,\;min} \cdot \mathbbm{1}\{P_{DG}(t) > 0\} \; \leq \; P_{DG}(t) \; \leq \; P_{DG,\;r} \\
SOC_{min}\cdot \vv{\mathbbm{1}}_{1\times25} \; \leq \;\vv{\bm{SOC}}(t) \;\leq\; SOC_{max}\cdot \vv{\mathbbm{1}}_{1\times25}, \quad -P_{BS,\;r}\cdot \vv{\mathbbm{1}}_{1\times24} \;\leq \;\vv{\bm{P_{BS}}}(t) \; \leq \; P_{BS,\;r}\cdot \vv{\mathbbm{1}}_{1\times24} \\
DPSP \; \leq \;DPSP_{max}, \quad \Sigma_{i=1}^{i=4} \; w_i  = 1, \; \mathbbm{1}\{P_{DG}(t) > 0\} = 
\begin{cases} 
    1 & \text{if } P_{DG}(t) > 0 \; \text{i.e., } DG_{ON}(t) = 1\\
    0 & \text{otherwise}
\end{cases} 
\end{gather*}

where $\vv{\bm{P_{DG}}}$ and $\vv{\bm{P_{BS}}}$ are $1\times24$ row vectors representing hourly DG and BS power flows while $\vv{\bm{SOC}}$ is a $1\times25$ row vector (the SOC vector has an extra element, $SOC(25)$ to ensure that the SOC of the BS remains within limits at the end of the last hour of the day as well). The value of $DPSP_{max}$ is set according to the desired reliability of the system and specifies the maximum proportion of daily capacity shortage allowed. Thus, the value of both the sizing and dispatch objective functions always lie between 0 and 1. Initially both problems are run using equal weights on all objectives. Thus, $w_i = 0.2 \; \forall \; i$ in the sizing optimization and $w_i = 0.25 \; \forall \; i$ in the dispatch optimization. 

\subsection{Mathematical properties of the model}

A convex optimization problem is one with a convex objective as well as all convex constraints. In the case of both the sizing and dispatch optimization problems, although all the constraints are linear and thus also convex, the objective functions themselves are non-convex, nonlinear and non-smooth (discontinuous and non-differentiable at certain points) primarily due to the use of discrete variables like $DG_{ON}$ and $BS_{cycled}$ tracking the operational state of the DG and cycling of the battery in each time step. There are also sub-models included within the objective function program in MATLAB that are non-linear. The numbers of PV modules $n_s$ and wind turbines $n_w$ are integers. Our sizing runs search an integer-coded space: the solvers move in the continuous box, the objective rounds $n_s$ and $n_w$ to the nearest integers before simulating, so every evaluated design can be built, and a final one-dimensional search re-optimizes $E_{b,\;init}$ for the rounded design. \cref{sec:results} also reports the alternative relax-and-round strategy, in which the solvers optimize the continuous relaxation and only the returned point is rounded. Mixed-integer formulations of such problems are NP-hard \cite{parisio2013}. Furthermore, it can be clearly seen from above that the number of constraints rapidly increases if a longer study period ($T$) is chosen or smaller time steps are used since the dimensions of the vector input variables would also increase accordingly. Complexity and runtime scale relatively poorly with problem size and thus, it was decided to perform only day-ahead dispatch with $1 \; h$ steps as a starting point. Future work could look into more efficient ways of formulating this problem. MATLAB's inbuilt solvers from the optimization and global optimization toolboxes were used to solve the problems.

\section{Results and discussion \label{sec:results}}

All these results were obtained using a Dell XPS 15 9560 system with a 2.8GHz Intel Core i7-7700HQ (3.8GHz boost) having 4 cores and 8 threads, alongside an NVIDIA GTX 1050 GPU with 4GB RAM. 

\subsection{Comparison of solution algorithms for MG planning and design optimization}

Due to the problem formulated in \cref{sec:opt_prob} being non-convex and non-smooth, the objective function has several local minima. Thus, it may output different local optima, depending on the initial value chosen by the solver. This necessitates the use of solvers that fully search the space of feasible solutions to find global, rather than local optima.  Various solvers from MATLAB's global optimization toolbox were experimented with to find globally optimal solutions, as summarized in \cref{tab:solver_comp}. This testing was done for a base case LI+DE system i.e., using a DE as the backup DG and LI as the BS, with the DE rated at 16 kW, slightly higher than the peak daily load. All of the solvers were run using parallel processing to speed up computation, except for SA where this is not an option. Most of these global solvers are stochastic in the sense that they use a population of multiple randomized local solvers with different initial values (start points) in order to find a global optimum. Thus, a random number seed was fixed in order to ensure the reproducibility of results. Furthermore, equal weights were placed on all five objectives ($w_i = 0.2 \; \forall \; i=1,2,...5$) indicating that all five goals are prioritized equally by the MG designer.

The maximum number of function evaluations and iterations were both set to 2000 in order to prevent the solver from stopping prematurely due to hitting either of these limits. This ensured that all solvers terminated only upon reaching a global solution indicated by the stopping criteria. For most solvers, this occurs when the relative change in the objective function value during a step is less than the specified tolerance or lower limit. For direct (pattern) search, this happens when the mesh size becomes less than the mesh tolerance. In general, all these tolerances and thresholds were kept at their default values supplied by MATLAB. In addition to the function tolerance, these include lower bounds on the step size, mesh size and first-order optimality measure as well as upper bounds on the magnitude of constraint functions (to detect any possible violations). Finally, for SO, there is no set stop criterion - the solver finishes based on the computational budget assigned by the user, i.e., an upper limit on either function evaluations or iterations. For algorithms that required an initial guess (i.e., GS, MS, and SA), $n_s = 50$,\;$n_w = 3$ and $E_{b,\;init}=50\; kWh$ was used as the start value. The lower and upper bounds were set to [0,\;0,\;0] and [100,\;30,\;200] respectively. The pattern search solver was repeatedly run using 89 different start points randomly initialized within the finite bounds. The global and multiple starting point search solvers were run using the \texttt{active-set} algorithm for minimizing constrained nonlinear multivariable functions since this allows them to take large steps which adds speed. 

\begin{table}[htbp]
\centering
\begin{tabular}{p{5cm} p{1.25cm} p{0.5cm} p{0.5cm} p{2cm}}
\toprule
\multicolumn{1}{l}{\textbf{Solver}} & \multicolumn{1}{c}{\textbf{t [$s$]}} & \multicolumn{1}{c}{\textbf{Min obj.}} & \multicolumn{1}{c}{\textbf{Overall}} & \multicolumn{1}{c}{\textbf{Optimal soln}} \\ \midrule
Particle swarm optimization (PSO) &  63.27  &  0.1460  & 9.24 & [11,    4,   70.39]\\
Genetic algorithm (GA)  &  289.75  &  0.1471  & 42.62 &  [20,    4,   70.15] \\
Multiple start search (MS) & 724.07 & 0.1489 & 107.81 & [25,    4,   61.63] \\
Simulated annealing (SA) & 856.82 & 0.1489 & 127.62 & [5,    4,   83.90] \\
Global search (GS) & 868.86 & 0.1487 & 129.20 & [13,    4,   81.65] \\
Direct pattern search (DPS) & 1875.93 & 0.1460 & 273.89 & [11,    4,   70.29] \\
Surrogate optimization (SO) & 2228.32 & 0.1480 & 329.79 & [21,    4,   66.67] \\
\bottomrule
\end{tabular}
\caption{\label{tab:solver_comp} Optimization results and performance (in terms of normalized objective value) of MATLAB's global solvers.} 
\end{table}

\paragraph{Note on computation times}
The runtimes in \cref{tab:solver_comp} (63\,s for PSO up to 2228\,s for SO) may appear surprisingly low for an 8760-hour annual optimization, but they reflect three key design decisions. First, the dispatch within the sizing loop is rule-based (priority ordering, \cref{sec:disp_strat}), not an inner optimization, so each of the 8760 hourly time steps requires only a simple sequential evaluation rather than solving a subproblem. Second, the sizing decision space is three-dimensional ($n_s, n_w, E_{b,\text{init}}$), enabling fast function evaluations per PSO particle. Third, all solvers were run with parallel processing across 8 threads on a 4-core Intel Core i7 processor. MILP-based approaches, which are well-suited to the day-ahead dispatch subproblem (Stage~2) but would require linearized approximations for the non-smooth, non-convex annual sizing problem (Stage~1), were not included in this comparison; their computational overhead for the full annual problem would be substantially higher. Furthermore, modeling the inner loop dispatch in stage 1 as an MILP (rather than rule-based) would entail a multiperiod optimization to account for the battery's intertemporal constraints --- running such a model at an hourly resolution would be prohibitively expensive and make it more challenging. While our approach does potentially trade off some optimality, it is much more tractable and also allows us to rapidly run more case studies and scenarios, such as different strategies (\cref{sec:strat}), multiple technology combinations (\cref{sec:tech_select}), Pareto frontier analyses (\cref{sec:pareto_opt}), and numerous sensitivity analyses. A discussion comparing our results to MILP dispatch formulations in the literature is provided in \cref{sec:dispatch_opt}.

The optimal solution in \cref{tab:solver_comp} lists the values of [$n_s$,$n_w$ and $E_{b,\;init}$] along with the runtime and minimum value of the objective function attained by the solver. The lower the value of the objective, the greater the optimality of corresponding solutions, and the lower the time taken, the faster the solver's computational performance. Thus, an overall performance metric can be devised as the product of runtime and minimum objective value $Overall = Min \; obj \cdot t$ and the solvers in \cref{tab:solver_comp} are listed in ascending order of this measure, with lower performance metric values indicating a 'better' solver. All of the above solvers converged with successful outcomes, i.e., the solver did not halt by exceeding an iteration or function evaluation limit. The results are also depicted graphically in \cref{fig:solver_comp}.

\begin{figure}[htbp]
  \centering
  \includegraphics[width=0.6\linewidth]{solver_comp.pdf}
  \caption{\label{fig:solver_comp} Plot of minimum objective value vs run time.}
\end{figure}

As can be seen from \cref{tab:solver_comp} and \cref{fig:solver_comp}, GA and PSO give the best overall performance when considering a balance of both speed (efficiency) and accuracy (optimality). However, such genetic or evolutionary programming methods and swarm intelligence algorithms have little supporting theory and do not guarantee globally optimal solutions. Solver parameters (like population size and initial population range in case of GA) need to be carefully tuned to thoroughly search the hyperspace of possible minima. Despite the lack of proven global convergence, these solvers are usually found to work quite efficiently in practice. Vectorization of the objective can make these methods run even faster. 

Pattern search is also able to find the most optimal solution (minimum value of weighted cost function) but takes a much longer time than GA or PSO to converge. Some of the data points like SO, MS, GS and SA in \cref{fig:solver_comp} contradict intuition since one would expect increased running time to result in improved objectives, but this is not always the case. This implies that biologically-inspired methods like GA and PSO are just better suited to the optimization problem studied here. Also, MS and GS are both gradient-based solvers which could explain why they are much faster than DPS and SO. Although surrogate optimization has good supporting theory and has been proven to converge to a global minimum, it was found to be extremely slow in practice when used to solve the sizing optimization problem. SO took the longest to converge while using the same upper limit on the maximum allowed number of function evaluations (= 2000). Even then, the minimum objective obtained by it is sub-optimal and higher than those obtained by other methods like PSO, DPS and GA.

Similarly, SA is guaranteed to converge to global optima but only with a logarithmic cooling schedule which turns out to be extremely slow in practice.  These examples show how theoretical expectations may not always match actual computational performance. Thus, there is a need to test and tune solvers on the specific problem structure being considered and understand their behavior. Since PSO has the lowest performance metric value, it was concluded as the best solver for this optimization problem and was used to produce all subsequent sizing results. A notable pattern in \cref{tab:solver_comp} is that all seven solvers converge to $n_w = 4$ wind turbines, despite arriving at different values of $n_s$ and $E_{b,\text{init}}$. This suggests that the optimal wind turbine count is a robust feature of the problem landscape --- the objective surface has a well-defined, sharp minimum in the $n_w$ direction that all solvers reliably locate. By contrast, $n_s$ ranges from 5 to 25 across solvers, indicating a flatter, more degenerate objective in the solar dimension where multiple PV capacities yield similar cost-reliability trade-offs (particularly when compensated by adjustments to battery size). This also implies that the sizing decision for the PV component is most sensitive to solver choice and initialization, and practitioners should pay particular attention to ensure adequate search coverage in that dimension.

\subsection{Sizing results}

\subsubsection{\label{sec:strat} Strategy selection}

Before optimizing the capacities of the PV, WT and BS components, two strategy decisions had to be made, relating to (1) the replacement of the BS and (2) whether surplus DG output may charge the BS.
\begin{itemize}
    \item The operator can replace the battery at fixed intervals set by its warranty or expected life (here every 15 years for LI), or when it has completed its rated number of charge--discharge cycles (5,475 for LI and 1,400 for LA). The second approach follows actual use and is the more realistic of the two.
    \item The DG has a minimum stable output. When the residual deficit is below this level, the DG still runs at its minimum and produces a surplus. This surplus can either be rectified and stored in the BS or sent to the dump load.
\end{itemize}

\begin{table}[htbp] 
\centering
\begin{tabular}{p{0.75cm} p{10cm}}
\toprule
\multicolumn{1}{c}{\textbf{Strategy}} & \multicolumn{1}{c}{\textbf{Description}}    \\
\midrule
1 & DE minimum-load surplus can charge LI; LI replaced every 15 y \\[1ex]
2 & DE surplus is dumped; LI replaced every 15 y \\[1ex]
3 & DE minimum-load surplus can charge LI; LI replaced every 5475 cycles \\[1ex]
4 & DE surplus is dumped; LI replaced every 5475 cycles \\
\bottomrule
\end{tabular}
\caption{Strategy selection for the LI+DE microgrid.}
\label{tab:strat_select} 
\end{table}

\cref{tab:strat} lists the best design found for each strategy (best of ten integer-coded PSO runs). Strategy 3 gives the lowest weighted objective (0.0863) and the lowest absolute LCOE (0.103~USD/kWh), and we use it in all subsequent results. Strategy 4 reaches lower emissions and a higher renewable fraction, but only with a 144~kWh battery and more dumped energy. Cycle-based replacement lowers costs because the battery cycles only about 63 times a year, so it is replaced far less often than a fixed 15-year schedule assumes.

The DE never runs in order to charge the battery. Under the load-following rule it starts only when renewable output and battery discharge cannot cover the load. When the residual deficit is smaller than the engine's minimum stable output ($0.3P_{DG,\;r}$), the engine still runs at that minimum, and the difference is surplus energy. Strategies 2 and 4 send this surplus to the dump load; Strategies 1 and 3 rectify it ($\eta_{rec} = 0.9$) and store it (round-trip efficiency 0.9). Both options burn the same fuel in that hour, so storage recovers energy that would otherwise be wasted. Battery wear from the extra throughput is costed through cycle-based replacement and capacity fade. To test whether the preference for Strategy 3 depends on cheap fuel or a long battery life, \cref{tab:strategy_stress} repeats the comparison of Strategies 3 and 4 at diesel prices up to 7.57~USD/gal (2.0~USD/L) and with the cycle life cut to 3,000. Strategy 3 keeps the lower weighted objective in all six combinations, although at diesel prices of 7.57~USD/gal and above Strategy 4 reaches a slightly lower absolute LCOE because its larger battery displaces more fuel. The cycle life does not bind in either strategy: the battery completes only about 45--65 cycles a year, so even 3,000 cycles last longer than the 25-year project.

\begin{table}[htbp]
\centering
\small
\begin{tabular}{l l c c c c c c c}
\toprule
\textbf{Strategy} & \textbf{Soln. $[n_s, n_w, E_{b,init}]$} & \textbf{$J$} & \textbf{LCOE} & \textbf{Em.} & \textbf{DPSP} & \textbf{Dump} & \textbf{1-REF} & \textbf{LCOE [USD/kWh]} \\
\midrule
1 & [23, 4, 67.31] & 0.0997 & 0.2650 & 0.0705 & 0.0000 & 0.1094 & 0.0536 & 0.132 \\
2 & [9, 5, 114.32] & 0.1140 & 0.2907 & 0.0515 & 0.0000 & 0.1938 & 0.0341 & 0.144 \\
3 & [32, 4, 107.27] & 0.0863 & 0.2077 & 0.0539 & 0.0000 & 0.1300 & 0.0401 & 0.103 \\
4 & [5, 5, 143.68] & 0.0965 & 0.2180 & 0.0476 & 0.0000 & 0.1851 & 0.0321 & 0.108 \\
\bottomrule
\end{tabular}
\caption{\label{tab:strat} Best design for each strategy (16~kW DE, equal weights). LCOE and emissions are normalized by the diesel-only baseline (\cref{tab:baseline}).}
\end{table}

\begin{table}[htbp]
\centering
\small
\begin{tabular}{c c c c c c c c}
\toprule
\textbf{Diesel price} & \textbf{LI cycle life} & \multicolumn{2}{c}{\textbf{$J$}} & \multicolumn{2}{c}{\textbf{LCOE [USD/kWh]}} & \multicolumn{2}{c}{\textbf{Dump fraction}} \\
\textbf{[USD/gal]} & \textbf{[cycles]} & S3 & S4 & S3 & S4 & S3 & S4 \\
\midrule
3.20 & 5,475 & 0.0863 & 0.0965 & 0.103 & 0.108 & 0.130 & 0.185 \\
3.20 & 3,000 & 0.0863 & 0.0965 & 0.103 & 0.108 & 0.130 & 0.185 \\
5.00 & 5,475 & 0.0781 & 0.0874 & 0.114 & 0.118 & 0.130 & 0.185 \\
5.00 & 3,000 & 0.0781 & 0.0874 & 0.114 & 0.118 & 0.130 & 0.185 \\
7.57 & 5,475 & 0.0720 & 0.0807 & 0.137 & 0.133 & 0.123 & 0.185 \\
7.57 & 3,000 & 0.0720 & 0.0807 & 0.137 & 0.133 & 0.123 & 0.185 \\
\bottomrule
\end{tabular}
\caption{\label{tab:strategy_stress} Strategy 3 (DE surplus stored) vs Strategy 4 (DE surplus dumped) under higher diesel prices and shorter battery cycle life; best of five PSO runs per cell.}
\end{table}

\subsubsection{Technology selection\label{sec:tech_select}}
The MG designer also has to choose the battery and backup generator technologies. Four combinations were sized (\cref{tab:tech_select}). LI+DE gives the lowest weighted objective. LA+DE reaches a lower normalized LCOE but dumps more than twice as much energy, because its small, cycle-limited battery absorbs little of the renewable surplus. Both MT combinations have markedly higher emissions and fossil shares. LCOE and emissions are normalized with respect to a single-DG system of the same technology; the baseline values are listed in \cref{tab:baseline}.

\begin{table}[htbp]
\centering
\small
\begin{tabular}{l l c c c c c c}
\toprule
\textbf{System} & \textbf{Soln. $[n_s, n_w, E_{b,init}]$} & \textbf{$J$} & \textbf{LCOE} & \textbf{Em.} & \textbf{DPSP} & \textbf{Dump} & \textbf{1-REF} \\
\midrule
LI+DE & [32, 4, 107.27] & 0.0863 & 0.2077 & 0.0539 & 0.0000 & 0.1300 & 0.0401 \\
LI+MT & [13, 4, 64.12] & 0.1133 & 0.3300 & 0.0835 & 0.0000 & 0.0880 & 0.0653 \\
LA+DE & [35, 4, 24.72] & 0.1083 & 0.1614 & 0.0534 & 0.0000 & 0.2875 & 0.0393 \\
LA+MT & [12, 4, 24.02] & 0.1317 & 0.2910 & 0.0779 & 0.0000 & 0.2281 & 0.0614 \\
\bottomrule
\end{tabular}
\caption{\label{tab:tech_select} Best design for each technology combination (best of five to ten integer-coded PSO runs).}
\end{table}

\begin{table}[htbp]
\centering
\begin{tabular}{@{}lcc@{}}
\toprule
\textbf{Parameter} & DE system & MT system \\ 
\midrule
LCOE [$\$/kWh$] & 0.4968 & 0.3321 \\
Emissions [$kg$ of pollutants/y] & 50,189 & 48,562 \\
Emissions [$kg$ of pollutants/$kWh$] & 0.6774 & 0.6555 \\
\bottomrule
\end{tabular}
\caption{\label{tab:baseline} Baseline LCOE and emissions for the MG running on a single fossil-fuel 16 kW DG.}
\end{table}

Microturbines running on natural gas have lower specific emissions than diesel engines and, alone, a lower LCOE (\cref{tab:baseline}). Yet the MG with a DE backup has lower emissions and costs than the one with an MT backup, because the backup runs for only about 597~h per year: the MT's higher capital cost is not recovered by its fuel savings over so few operating hours. The same reasoning explains why LI is preferred over LA despite LA's lower capital cost.

\subsubsection{Optimal capacities for baseline LI+DE system}

We sized the LI+DE system with two sets of decision variables: (i) unit counts $[n_s, n_w, E_{b,\;init}]$ with integer $n_s$ and $n_w$, and (ii) total rated powers $[P_{PV,\;r,\;total}, P_{WT,\;r,\;total}, E_{b,\;init}]$ treated as continuous. Design (i) corresponds to [8.16~kW, 14.0~kW, 107.3~kWh] in rated terms and design (ii) to [4.03~kW, 14.97~kW, 103.5~kWh]. The continuous approach (ii) reaches a slightly lower weighted objective (0.0836 vs 0.0863) because it is not restricted to whole PV modules and turbines; a designer limited to specific products should use approach (i).

\begin{table}[htbp]
\centering
\begin{tabular}{p{5cm} p{4.3cm} p{4.3cm}}
\toprule
\multicolumn{1}{l}{\textbf{Parameter}} & \multicolumn{1}{l}{(i) unit counts} & \multicolumn{1}{l}{(ii) rated powers} \\
\midrule
$LCOE_{normalized}$ & 0.2077 & 0.1999 \\
LCOE [$\$/kWh$] & 0.1032 & 0.0993 \\
$Emissions_{normalized}$ & 0.0539 & 0.0542 \\
Annual emissions [$kg/y$] & 2,703 & 2,722 \\
DPSP & 0 & 0 \\
Dump & 0.1300 & 0.1234 \\
1-REF & 0.0401 & 0.0404 \\
No. of online DE hours [$h/y$] & 597 & 594 \\
No. of LI BS cycles [$/y$] & 63 & 70 \\
Optimal solution & $[n_s, n_w, E_{b,\; init}]$ = [32, 4, 107.27] & $[P_{PV}, P_{WT}, E_{b}]$ = [4.03, 14.97, 103.47] \\
Weighted objective $J$ & 0.0863 & 0.0836 \\
\bottomrule
\end{tabular}
\caption{\label{tab:LIDE} Sizing results for the baseline LI+DE system using two sets of decision variables.}
\end{table}

With either approach the hybrid microgrid improves substantially on the diesel-only baseline: LCOE falls by 79\% (from 0.497 to 0.103~USD/kWh), emissions by 95\%, and renewables supply 96.0\% of generation. The battery cycles only about 63--70 times a year, so at 5,475 cycles it would not need replacement within the 25-year project life. The DE runs about 594--597~h a year, so its 15,000-h life implies roughly one replacement over the project.

The renewable capacity of design (i) is 22.2~kW, 162\% above the average load of 8.46~kW and 48\% above the annual peak of 15.0~kW. The battery stores 53\% of the average daily consumption of 203~kWh. \cref{fig:sizing_operation} shows the resulting rule-based operation over a day, a week and a month. Wind dominates the renewable supply (88\% of renewable energy), because the assumed installed cost per kW is lower for wind than for PV and the scaled station data imply a hub-height mean wind speed of 4.7~m/s. \cref{app:climate} and \cref{sec:sizing_sensitivity} examine how sensitive this result is to the wind data.

\begin{figure}[htbp]
  \centering
  \includegraphics[width=\linewidth]{rev_sizing_operation.pdf}
  \caption{\label{fig:sizing_operation} Rule-based operation of the selected LI+DE design (16~kW DE) over 1 January 2018, the first week and the first 30 days: power flows (left) and battery state of charge (right). Battery power is positive when discharging.}
\end{figure}

Global averages of installed costs were used for both PV and WT. Small-scale PV costs have fallen faster than small-scale wind costs over the past decade, so component-specific prices would likely shift the optimum toward PV.

\subsection{Variation of Pareto-optimal combinations of input variables\label{sec:pareto_opt}}

The weighted-sum problem in Eq.~(\ref{eq:sizing}) has a single minimum value, but the five-objective problem behind it has no single optimum: along the Pareto front, any improvement in one objective worsens at least one other. The Pareto set contains the designs that no other design matches or beats in every objective while beating them in at least one. We approximate this set with the genetic (\texttt{gamultiobj}) and pattern-search (\texttt{paretosearch}) multiobjective solvers in MATLAB's Global Optimization Toolbox. \Cref{fig:pareto_inputs} shows how the three decision variables vary across the non-dominated set found by \texttt{gamultiobj}: black dots are the non-dominated designs returned by the solver, the surface interpolates between them, and the colour shows battery capacity. These points are mutually non-dominated: they represent different trade-offs, not equally good designs, and choosing among them requires stated preferences or a selection rule (\cref{sec:pareto_select}). 

\begin{figure}[htbp]
  \centering
  \begin{subfigure}[b]{0.5\linewidth}
    \includegraphics[width=\linewidth]{Sizing_LI-DE_5equal_gamultiobj_Ps_Pw_fraction0-5_mesh_grid.png}
  \end{subfigure}
  \caption{\label{fig:pareto_inputs} Pareto set obtained using \texttt{gamultiobj} and optimizing with respect to $P_{PV,\;r,\;total}$, $P_{WT,\;r,\;total}$ and $E_{b,\;init}$.}
\end{figure}

% complementarity between wind and solar power availability. solar power peaks along with solar radiation during mid-day while wind power and wind speeds are highest at night when solar is unavailable. 
% at low RES penetrations BS capacity is also low since it is not sensible to utilize DE or MT output to charge the batteries since this would be less efficient (due to charge/discharge losses). In such cases, the MG would rely on using the fossil-DG directly.

In addition, in order to more directly show the relationships between pairs of inputs, 2D graphs were also plotted in \cref{fig:pareto_inputs3}, when all three input values belong to the Pareto set. It can be seen from \cref{fig:pareto_inputs3} that, in general, there is a positive relationship between all three pairs of inputs, although the strength of the correlation varies. \cref{fig:WT_PV_pareto} shows that when installed PV capacity increases, the optimal WT capacity also increases weakly. However, when the increase in either quantity is very large, this can lead to a decrease in the other. This results in the relatively low value of $R^2 = 0.4182$, when fitted with a second-order polynomial curve. This makes sense physically since if the system is already heavily reliant on a large installed capacity of either PV or WT, then less additional capacity of WT or PV, respectively, needs to be installed to satisfy the load. Although none of these combinations dominates another, a designer may still prefer a solution offering a more diversified and balanced portfolio of generation sources, as opposed to one that is too heavily skewed towards either PV or WT. 

Similarly, as seen in \cref{fig:BS_PV_pareto} and \cref{fig:BS_WT_pareto}, the installed BS capacity increases both with an increase in PV capacity as well as an increase in WT capacity. This positive relationship is stronger than in \cref{fig:WT_PV_pareto}, with $R^2 = 0.6055$ for the BS vs PV capacity curve (fitted with a power curve) and $R^2 = 0.6152$ for the BS vs WT capacity (fitted with a linear trendline). These results follow from the fact that renewables like wind and solar power are both intermittent and unpredictable and their time of availability does not necessarily match up well with the actual time of use of electricity by customers in the MG. Thus, as RES capacity increases, the BS capacity needed will also likely need to increase on account of this uncertainty and mismatch between supply and demand, in order to ensure reliable MG operation. This is also reflected in \cref{fig:pareto_inputs}. Here, we also observe interestingly that the optimal BS capacity is generally lower when levels of both PV and WT are comparable, as opposed to when their contributions are very different. This is due to the complementarity between these two RES in terms of their availability --- PV output is unimodal and peaks in the middle of the day while WT is more spread out over the day, as seen in \cref{fig:chap3_hourly}. We also note that a system dominated by WT generally requires more BS than a PV-dominant system since the region studied has a higher solar resource than wind power.

\begin{figure}[htbp]
  \centering
  \begin{subfigure}[b]{0.49\linewidth}
    \includegraphics[width=\linewidth]{PV_WT_pareto_set.png}
    \caption{\label{fig:WT_PV_pareto} Variation of installed WT and PV capacity.}
  \end{subfigure}
  \begin{subfigure}[b]{0.49\linewidth}
    \includegraphics[width=\linewidth]{BS_PV_pareto_set.png}
    \caption{\label{fig:BS_PV_pareto} Variation of installed PV and BS capacity.}
  \end{subfigure}
  \begin{subfigure}[b]{0.5\linewidth}
    \includegraphics[width=\linewidth]{BS_WT_pareto_set.png}
    \caption{\label{fig:BS_WT_pareto} Variation of installed WT and BS capacity.}
  \end{subfigure}
  \caption{\label{fig:pareto_inputs3} 2D plots showing relationships between Pareto-optimal values of the inputs $P_{PV,\;r,\;total}$, $P_{WT,\;r,\;total}$ and $E_{b,\;init}$.}
\end{figure}

% If my objective function is also convex, then the Pareto optimal front or single, unique optimal solution is expected to occur at the boundary of the feasible region/set. 

\subsection{Variation of different objectives on the Pareto front}

Similar to the analysis conducted for the inputs, the Pareto front can also be used to visualize how the optimal, minimum attainable values of the objectives vary with respect to each other. This helps identify trade-offs as well as co-benefits that arise while pursuing different goals during MG design and sizing. In order to obtain some non-zero values for DPSP, this Pareto set was produced using a smaller DG rated at 7 kW. It was found that the $REF$ is strongly negatively correlated with the normalized $Emissions$ term in a nearly linear fashion (see \cref{fig:correl}). This makes intuitive sense since raising the penetration of clean renewables leads to lower emissions. Thus, the Pareto-set generation was simplified by excluding the $(1-REF)$ objective for simplicity, since its trends would follow those of emissions.

The Pareto front was then calculated using pattern search on the remaining four objectives with 200 points so as to obtain a dense set. 2D plots in \cref{fig:pareto_obj} show the relationships between different pairs of objectives. Reducing emissions and load deficiency both raise system costs as shown in \cref{fig:LCOE_Emissions} and \cref{fig:LCOE_DPSP}. Minimizing the amount of energy wasted (i.e., grounded or sent to dump loads) requires increasing BS capacity and/or operating the DE more frequently (with a reduction in installed RES capacity), which generally results in higher costs as in \cref{fig:LCOE_Dump}. Thus, some non-zero dump energy ratio is usually necessary for optimal operation, the value being $\approx$ 10-25\% of total generation in this case. However, beyond a certain point, a further increase in dumped energy also increases LCOE since the system is likely oversized implying higher capital and operational costs. 

\begin{figure}[htbp]
\centering
    \includegraphics[width=0.33\linewidth]{Proof_REF_Emissions_highly_correlated_linear.png}
  \caption{\label{fig:correl} Almost linear relationship obtained between renewable energy fraction and normalized emissions.}
\end{figure}

Interestingly, emissions rise with DPSP because the load is lost most often when the backup DE (even when operated at its full rated power) is insufficient to meet demand. Thus, higher DPSP values generally indicate excess demand situations caused by peaky, sub-optimal load profiles that require the DE to run for longer periods and increase diesel consumption. Finally, normalized emissions in \cref{fig:Emissions_Dump} as well as DPSP in \cref{fig:Dump_DPSP} fall with a rise in the dumped energy ratio, both of which make intuitive sense. An increase in dumped energy implies more extensive use of cleaner renewables leading to more curtailment. Keeping the lost load (DPSP) at low levels also often requires meeting excess demand by bringing the backup generator online, which can result in excess power production, particularly if the excess demand remaining to be met is below the DG's rated power.

\begin{figure}[htbp]
  \centering
  \begin{subfigure}[b]{0.33\linewidth}
    \includegraphics[width=\linewidth]{Pareto_123_LCOEvsEmissions.png}
    \caption{\label{fig:LCOE_Emissions}}
  \end{subfigure}
  \begin{subfigure}[b]{0.33\linewidth}
    \includegraphics[width=\linewidth]{Pareto_123_LCOEvsDPSP.png}
    \caption{\label{fig:LCOE_DPSP}}
  \end{subfigure}
  \begin{subfigure}[b]{0.33\linewidth}
    \includegraphics[width=\linewidth]{Pareto_1234_LCOEvsDump.png}
    \caption{\label{fig:LCOE_Dump}}
  \end{subfigure}
  \begin{subfigure}[b]{0.33\linewidth}
    \includegraphics[width=\linewidth]{Pareto_123_EmissionsvsDPSP.png}
    \caption{\label{fig:Emissions_DPSP}}
  \end{subfigure}
  \begin{subfigure}[b]{0.33\linewidth}
    \includegraphics[width=\linewidth]{Pareto_1234_EmissionsvsDump.png}
    \caption{\label{fig:Emissions_Dump}}
  \end{subfigure}
  \begin{subfigure}[b]{0.33\linewidth}
    \includegraphics[width=\linewidth]{Pareto_1234_DumpvsDPSP.png}
    \caption{\label{fig:Dump_DPSP}}
  \end{subfigure}
  \caption{\label{fig:pareto_obj} Two-way correlations between objectives derived from Pareto-set obtained via \texttt{paretosearch}.}
\end{figure}

\subsection{Break-even distance analysis}

The break-even distance (BED) is the distance from the existing grid beyond which the stand-alone MG becomes cheaper than extending the grid \cite{kaabeche2011a}:
\begin{align*}
    BED \; [km] & = \frac{TAC - LCOE_{grid} \cdot \Sigma_{t=1}^{t=8760}\; P_{load}(t)\cdot \Delta t}{C_{ext} \cdot CRF}
\end{align*}
where $LCOE_{grid}$ is the levelized cost of grid electricity and $C_{ext}$ the cost of line extension. For the selected design, $TAC = 7,645$~USD/y and $CRF = 0.0582$. With $LCOE_{grid} = 0.125$~USD/kWh \cite{LCOEgrid} and $C_{ext} = 157,470$~USD/km \cite{Cext} for Kenya, buying the community's 74,086~kWh from the grid would cost 9,261~USD/y before any extension cost, more than the MG's annualized cost. The break-even distance is therefore zero (the formula returns -0.18~km): the MG is cheaper than grid supply at any distance. Timbila lies about 111~km from Voi, the nearest large town, and 305~km from Nairobi.

The MG's LCOE of 0.103~USD/kWh is at the low end of the 0.12--0.20~USD/kWh residential grid tariff in Kenya \cite{LCOEgrid}, and the tariff excludes connection fees of 136--400~USD per household that are prohibitive for rural customers \cite{quak_CBA, golumbeanu2013} as well as frequent outages in many peri-urban areas \cite{iea_africa2019}. Households without electricity pay far more for kerosene lighting (1.00--2.50~USD/kWh equivalent) \cite{tracy2012} or small diesel generators (0.80--1.50~USD/kWh) \cite{quak_CBA}. This comparison rests on the model's cost assumptions (IRENA 2018 installed costs and a real discount rate of 3.1\%); \cref{sec:sizing_sensitivity} shows how LCOE responds to them.

\section{Sensitivity analysis of sizing optimization\label{sec:sizing_sensitivity}}

We assess the sensitivity of the design by varying one parameter at a time and re-optimizing the sizing problem at each value; all other parameters keep the defaults in \cref{tab:default}. Each point is the best of three integer-coded PSO runs. Because LCOE is normalized by a diesel-only baseline that itself depends on fuel price and financial parameters (\cref{app:costs}), every figure reports absolute LCOE next to the normalized objectives.

\begin{table}[htbp]
\centering
\begin{tabular}{@{}lc@{}}
\toprule
\textbf{Parameter}                 & \textbf{Default value}                        \\ \midrule
DG power rating [$kW$]         & 16                                            \\
Diesel fuel price [USD/US gal] & 3.20                                          \\
Nominal interest rate (\%)         & 9                                             \\
Inflation rate (\%)                & 5.70                                           \\
Price of BS ($\$/kWh$)               & 300  \\
{[}$w_1$ $w_2$ $w_3$ $w_4$ $w_5${]}  &  [0.2 0.2 0.2 0.2 0.2] \\ 
\bottomrule
\end{tabular}
\caption{\label{tab:default} Default parameter values used for sensitivity analysis.}
\end{table}

\subsection{Varying DG power rating}
With the default 16~kW DG, DPSP is zero for every design because the DG exceeds the annual peak load. We varied the DG rating from 0 to 20~kW, keeping the normalization baseline at 16~kW (\cref{fig:vary_DG_rating}). DPSP is zero for ratings of 14~kW and above, and stays below 1\% down to 7~kW; an 8~kW DG gives DPSP = 0.0067. Larger generators raise LCOE (from 0.098~USD/kWh at 8~kW to 0.103 at 16~kW and 0.107 at 20~kW) because of their capital and fixed O\&M costs, while the PV, WT and battery capacities change little above about 8~kW. Without any DG the equal-weight objective is lowest ($J$ = 0.0598), even though 11.0\% of the load is lost: with equal weights a 10\% loss of load costs only 0.02 in $J$. Equal weights therefore undervalue reliability for an off-grid community, and \cref{sec:pareto_rel} treats reliability as an explicit constraint instead.

\begin{figure}[htbp]
  \centering
  \includegraphics[width=\linewidth]{rev_vary_DG_rating.pdf}
  \caption{\label{fig:vary_DG_rating} Effect of the DG power rating. (a) normalized objectives and weighted objective $J$; (b) absolute LCOE of the hybrid MG and of the diesel-only baseline; (c) rated PV and WT capacity; (d) battery capacity. Each point is the best of three integer-coded PSO runs.}
\end{figure}

\subsection{Varying fuel price}
The diesel price was varied from 0.2 to 12~USD/gal, the range of current global prices \cite{diesel_price} (\cref{fig:vary_diesel_price}). The absolute LCOE of the hybrid MG rises from 0.074 to 0.163~USD/kWh, while the normalized LCOE falls, because the diesel-only baseline rises much faster (from 0.180 to 1.427~USD/kWh). Higher fuel prices shift the design toward storage (battery capacity from 82 to 133~kWh) and cut DE operation from 570 to 576~h per year. At the 2.0~USD/L (7.57~USD/gal) level typical of remote sites the hybrid system's advantage over diesel-only supply grows further.

\begin{figure}[htbp]
  \centering
  \includegraphics[width=\linewidth]{rev_vary_diesel_price.pdf}
  \caption{\label{fig:vary_diesel_price} Effect of the diesel price. (a) normalized objectives and weighted objective $J$; (b) absolute LCOE of the hybrid MG and of the diesel-only baseline; (c) rated PV and WT capacity; (d) battery capacity. Each point is the best of three integer-coded PSO runs.}
\end{figure}

\subsection{\label{sec:varyInterest} Varying nominal interest or discount rate}
The nominal interest rate sets the opportunity cost of capital and hence the trade-off between capital-intensive renewables and storage and fuel-intensive backup generation. We varied it from 0 to 20\%, the current spread of central bank rates \cite{interest_rates} (\cref{fig:vary_interest}). With inflation fixed at 5.7\%, the real discount rate $r = (i - f)/(1 + f)$ ranges from $-5.4\%$ to $13.5\%$. Higher rates penalize the upfront investment: the battery shrinks from 142 to 67~kWh, the DE runs more (590 to 837~h per year), and LCOE rises from 0.064 to 0.169~USD/kWh. Interest rates in many SSA countries sit in the upper half of this range, which makes concessional finance an important lever for off-grid renewables.

\begin{figure}[htbp]
  \centering
  \includegraphics[width=\linewidth]{rev_vary_interest.pdf}
  \caption{\label{fig:vary_interest} Effect of the nominal interest rate. (a) normalized objectives and weighted objective $J$; (b) absolute LCOE of the hybrid MG and of the diesel-only baseline; (c) rated PV and WT capacity; (d) battery capacity. Each point is the best of three integer-coded PSO runs.}
\end{figure}

\subsection{Varying inflation rate}
Inflation was varied from 0 to 15\%, the range currently observed worldwide \cite{inflation_rate} (\cref{fig:vary_inflation}). At a fixed nominal interest rate of 9\%, higher inflation lowers the real discount rate, so the effect mirrors that of the interest rate: LCOE falls from 0.138 to 0.065~USD/kWh and the battery grows from 68 to 142~kWh as inflation rises from 0 to 15\%.

\begin{figure}[htbp]
  \centering
  \includegraphics[width=\linewidth]{rev_vary_inflation.pdf}
  \caption{\label{fig:vary_inflation} Effect of the inflation rate. (a) normalized objectives and weighted objective $J$; (b) absolute LCOE of the hybrid MG and of the diesel-only baseline; (c) rated PV and WT capacity; (d) battery capacity. Each point is the best of three integer-coded PSO runs.}
\end{figure}

\subsection{Varying price of battery energy storage (BS)}
Li-ion prices have fallen steeply and are projected to keep falling \cite{bnef}, so we varied the battery price from 50 to 300~USD/kWh (\cref{fig:vary_LI_price}). Cheaper storage enlarges the optimal battery (142~kWh at 50~USD/kWh vs 107~kWh at 300~USD/kWh), reduces DE operation (590 vs 597~h per year) and lowers LCOE (0.079 vs 0.103~USD/kWh).

\begin{figure}[htbp]
  \centering
  \includegraphics[width=\linewidth]{rev_vary_LI_price.pdf}
  \caption{\label{fig:vary_LI_price} Effect of the Li-ion battery price. (a) normalized objectives and weighted objective $J$; (b) absolute LCOE of the hybrid MG and of the diesel-only baseline; (c) rated PV and WT capacity; (d) battery capacity. Each point is the best of three integer-coded PSO runs.}
\end{figure}

\subsection{Varying the wind-speed correction factor\label{sec:windsens}}
The wind input rests on station data scaled by a factor of 3.70 (\cref{app:climate}). A comparison with NASA POWER reanalysis suggests that the resulting hub-height wind may be about 25\% too high, so we re-optimized the design with the wind input scaled by 0.75, 1.00 and 1.25 (\cref{tab:wind_sweep}). With 25\% less wind, the optimal design changes from [32, 4, 107] to [28, 5, 77], the DE runs 1081 instead of 597~h per year, LCOE rises from 0.103 to 0.136~USD/kWh and annual emissions from 2,703 to 5,049~kg. With 25\% more wind the PV array drops out ([0, 4, 119]) and LCOE falls to 0.087~USD/kWh. The conclusion that a hybrid system with renewables, Li-ion storage and DE backup beats diesel-only supply holds across this range, but the PV--WT split is sensitive to the wind data.

\begin{table}[htbp]
\centering
\small
\begin{tabular}{c l c c c c c c c}
\toprule
\textbf{Wind scale} & \textbf{$[n_s, n_w, E_b]$} & \textbf{PV [kW]} & \textbf{WT [kW]} & \textbf{$J$} & \textbf{LCOE [USD/kWh]} & \textbf{Em. [kg/y]} & \textbf{REF [\%]} & \textbf{WT share [\%]} \\
\midrule
2.77 & [28, 5, 76.5] & 7.14 & 17.5 & 0.1188 & 0.136 & 5,049 & 92.1 & 88 \\
3.70 & [32, 4, 107.3] & 8.16 & 14.0 & 0.0863 & 0.103 & 2,703 & 96.0 & 88 \\
4.62 & [0, 4, 118.6] & 0.00 & 14.0 & 0.0671 & 0.087 & 1,586 & 97.7 & 100 \\
\bottomrule
\end{tabular}
\caption{\label{tab:wind_sweep} Optimal design when the wind input is scaled by 0.75, 1.00 and 1.25 (equivalent correction factors 2.78, 3.70 and 4.63). WT share: wind share of renewable energy.}
\end{table}

\section{\label{sec:dispatch_opt} Optimal economic and environmental dispatch results}

The sizing optimization (Stage~1) embeds a rule-based, priority-ordered dispatch (\cref{sec:disp_strat}) within the PSO outer loop. This dispatch heuristic --- which prioritizes renewable generation, then battery, then diesel backup in fixed order -- is computationally fast but is not necessarily economically optimal: it does not minimize cost or emissions, does not jointly optimize battery and generator scheduling, and does not use look-ahead information about future load and renewable availability. The day-ahead dispatch optimization (Stage~2) presented in this section directly addresses these limitations by solving a multi-objective economic-environmental dispatch problem over a 24-hour horizon. As shown below, the optimized dispatch differs meaningfully from the rule-based strategy: it allows simultaneous operation of the battery and DG when economically beneficial, and uses the battery more strategically for peak shaving. These differences motivate the separation of Stage~1 (sizing with tractable rule-based dispatch) and Stage~2 (operation with optimized dispatch), and also illustrate the value of day-ahead scheduling relative to purely reactive, rule-based control. Compared to MILP-based dispatch formulations in the literature \cite{parisio2014, dercam}, the present nonlinear multi-objective formulation captures the nonlinear fuel cost of the diesel engine without linearization, but at the cost of not always guaranteeing global optimality; the \texttt{sqp} solver finds a locally optimal solution from the given initial condition.

The dispatch optimization was performed using the \texttt{fmincon} solver in MATLAB using initial guesses of $P_{LI}(t) = 1 \; kW$ and $P_{DE}(t) = 7 \; kW \; \forall \; t = 1:24$. The MG controller optimizes over the entire next day (time horizon = 24 h) using vector-valued variables $\vv{\bm{P_{DE}}}$ and $\vv{\bm{P_{BS}}}$, hourly time steps and assuming perfect foresight of load and weather. Although the \texttt{active-set} algorithm was used with \texttt{fmincon} earlier for sizing optimization, it can violate bound constraints at intermediate iterations. This caused the solver to sometimes converge to infeasible points while solving the dispatch problem. Thus, the sequential quadratic programming (\texttt{sqp}) algorithm was implemented here instead since it satisfies bounds during all iterations and is faster than the \texttt{interior-point} method. The maximum allowable load shortage was set to 1\% i.e., $DPSP_{max} = 0.01$. To start, the simulation was run using optimal PV, WT, and BS capacities found in \cref{tab:LIDE} for the LI+DE system with a 16 kW DG (ensures 100\% reliability).   

\begin{table}[htbp]
\centering
\begin{tabular}{p{5cm} p{0.1cm} p{0.1cm} p{0.1cm} p{0.1cm} p{0.1cm} p{0.1cm}}
\toprule
\multicolumn{1}{l}{\textbf{System}} & \multicolumn{1}{c}{\textbf{Min obj.}} & \multicolumn{1}{c}{\textbf{COE}} & \multicolumn{1}{c}{\textbf{Emissions}} & \multicolumn{1}{c}{\textbf{DPSP}} & \multicolumn{1}{c}{\textbf{Dump}} & \multicolumn{1}{c}{\textbf{REF}} \\
\midrule
(i) 16 kW DE & 0.4717 & 1.0291 & 0.6202  & 0.004 & 0.3559 & 0.6506     \\
(ii) 8 kW DE (old sizing) & 0.2695 & 0.6948 & 0.2469 & 0.007 & 0.2228 & 0.8239 
 \\
(iii) 8 kW DE (new sizing) & 0.2669 & 0.6895 & 0.2350 & 0.010 & 0.2346 & 0.8346   \\
\bottomrule
\end{tabular}
\caption{\label{tab:DG_size} Summary of the dispatch results and performance with a few different possible DG and RES sizes considered. Both COE and emissions are reported as normalized values. The minimum objective value reported here was calculated as an equally weighted linear combination of $COE$, $emissions$, $DPSP$, $Dump$ and $(1-REF)$.}
\end{table}

From \cref{tab:DG_size}, it can be seen that the MG using a 16 kW DE actually costs more daily than the baseline system running only on diesel. The main reason for poor performance is that there is a lower limit of 4.8 kW on the DE's power output and it needs to be run at this level even if the actual demand may be much lower. This unnecessarily increases fuel consumption, emissions and dumped energy. To improve this, a smaller DG size of 8 kW was also tested since this was found in \cref{fig:varyDE1} to be the lowest backup power rating that minimized weighted cost while still satisfying 99\% reliability. The 8 kW DG was first simulated using RES capacities equal to that for the 16 kW case and then with updated values found by rerunning the sizing optimization to account for the smaller DG size. According to \cref{tab:DG_size}, both cases (ii) and (iii) produce quite similar results (as seen in \cref{fig:dispatch_P_base}) and are a major improvement over the 16 kW option. As expected, the updated sizing values in (ii) produce slightly more optimal outcomes. However, the difference is very small and system (iii) was eventually chosen for further analysis since it is more diversified and well-balanced among PV, WT and BS dependence.

Comparing the actual dispatch obtained during sizing optimization in \cref{fig:sizing_dispatch} with \cref{fig:dispatch_P_base}, it is clear that the optimal strategy found here for day-ahead dispatch differs in several ways from the load-following strategy prescribed in \cref{sec:disp_strat} used for sizing purposes. The load-following strategy enforced a clear, fixed priority order in which generation units are dispatched. For instance, the DG is never brought online in \cref{fig:sizing_dispatch} unless both the RES and BS have been exhausted (i.e., either fully discharged or already discharging at maximum power) or are unavailable. However, this is not the case in \cref{fig:dispatch_P_base} as there are several time intervals where both the BS and DG are dispatched simultaneously, sometimes even during periods of high RES output. Thus, the unit commitment decision here is made according to the current (and future) supply-demand situation(s) rather than sticking to a preset scheduling strategy. 

\begin{figure}[htbp]
  \centering
  \begin{subfigure}[b]{0.8\linewidth}
    \includegraphics[width=\linewidth]{Dispatch_DG_16kW_OldSizing_base.png}
    \caption{\label{fig:dispatch_size1} 16 kW backup DG.}
  \end{subfigure}
  \begin{subfigure}[b]{0.8\linewidth}
    \includegraphics[width=\linewidth]{Dispatch_DG_8kW_OldSizing_base.png}
    \caption{\label{fig:dispatch_size2} 8 kW backup DG with old sizing.}
  \end{subfigure}
  \begin{subfigure}[b]{0.8\linewidth}
    \includegraphics[width=\linewidth]{Dispatch_DG_8kW_NewSizing_base.png}
    \caption{\label{fig:dispatch_size3} 8 kW backup DG with new sizing.}
  \end{subfigure}
  \caption{\label{fig:dispatch_P_base} Hourly power dispatch curves over one day (24 h) period on January 1, 2018.}
\end{figure} 

\section{Robustness of dispatch to uncertainty and disturbances}

Here, the robustness of the system's daily operation to uncertainties and disturbances such as sudden demand peaks and shocks to RES generation, is assessed. All the subsequent analysis was done using an 8 kW DG as opposed to the initially set default value of 16 kW, for reasons specified in \cref{sec:dispatch_opt} and also because meeting the DPSP requirement would be too trivial in the latter case.

\subsection{Intermittency and unpredictability of RES generation}

The system was tested under conditions of low solar irradiance as well as low wind speeds, in order to see how it responds to the variability of renewable sources. Both the low irradiance and low wind-speed days were set to be 10\% of the values for the baseline day used for the simulations in \cref{sec:dispatch_opt}, as shown in \cref{fig:irradiance_cases} and \cref{fig:windspeed_cases}. From \cref{tab:dispatch_robust_RES}, it can be seen that the islanded MG is very sensitive to sudden drops in renewable generation owing to unfavorable weather conditions. Since the system is largely dominated by wind power, it is particularly vulnerable to wind speed changes. In cases of extremely low winds, the MG resorts to running the backup DE almost continuously and close to its rated power, leading to sub-optimal performance by raising costs (even above the baseline diesel system) and emissions. Furthermore, the results presented here are only for one day and the performance may be even more negatively affected if such weather persists for longer periods. e.g., extended wind lulls lasting several consecutive days may drain all the BS and force the MG to rely purely on diesel. Surprisingly, the low solar irradiance situation marginally improves MG performance compared to case (ii) in \cref{tab:DG_size}. The main reason for this is the reduction in dump energy ratio since both peak PV and WT outputs tend to coincide near mid-day and thus less solar curtailment is needed with low PV output.

\begin{figure}[htbp]
\centering
\begin{subfigure}[b]{0.49\linewidth}
   \includegraphics[width=\linewidth]{irradiance_cases.png}
    \caption{\label{fig:irradiance_cases} Comparison of days with normal (baseline) and extremely low solar radiation.}
\end{subfigure}
\begin{subfigure}[b]{0.49\linewidth}
   \includegraphics[width=\linewidth]{windspeed_cases.png}
    \caption{\label{fig:windspeed_cases} Comparison of days with normal (baseline) and extremely low wind speeds.}
\end{subfigure}
\end{figure}

\begin{table}[htbp]
\centering
\begin{tabular}{p{5cm} p{0.1cm} p{0.1cm} p{0.1cm} p{0.1cm} p{0.1cm} p{0.1cm}}
\toprule
\multicolumn{1}{l}{\textbf{Scenario}} & \multicolumn{1}{c}{\textbf{Min obj.}} & \multicolumn{1}{c}{\textbf{COE}} & \multicolumn{1}{c}{\textbf{Emissions}} & \multicolumn{1}{c}{\textbf{DPSP}} & \multicolumn{1}{c}{\textbf{Dump}} & \multicolumn{1}{c}{\textbf{REF}} \\
\midrule
Low solar radiation & 0.2507 & 0.6793 & 0.2242 & 0.007  & 0.1699 & 0.8264     \\
Low wind speeds & 0.5119 & 1.047 & 0.6374 & 0.009 & 0 & 0.1335  \\
Low solar + low wind & 0.5331 & 0.9156 & 0.7 & 0.010 & 0.0546 & 0.0151  \\
\bottomrule
\end{tabular}
\caption{\label{tab:dispatch_robust_RES} Summary of dispatch results and performance under several scenarios representing RES uncertainty.}
\end{table}

\begin{figure}[htbp]
  \centering
  \begin{subfigure}[b]{0.8\linewidth}
    \includegraphics[width=\linewidth]{Dispatch_DG_8kW_OldSizing_lowPV.png}
    \caption{\label{fig:dispatch_RES1} Case with low solar irradiance levels.}
  \end{subfigure}
  \begin{subfigure}[b]{0.8\linewidth}
    \includegraphics[width=\linewidth]{Dispatch_DG_8kW_OldSizing_lowWT.png}
    \caption{\label{fig:dispatch_RES2} Case with low wind speeds.}
  \end{subfigure}
  \begin{subfigure}[b]{0.8\linewidth}
    \includegraphics[width=\linewidth]{Dispatch_DG_8kW_OldSizing_lowPVandWT.png}
    \caption{\label{fig:dispatch_RES3} Case with both low solar irradiance as well as wind speeds.}
  \end{subfigure}
  \caption{\label{fig:dispatch_robust_RES} Hourly power dispatch curves over one day (24 h) period under different RES output scenarios.}
\end{figure} 

\subsection{Stochastic demand}

% Also include histograms for load graphs since they're actually not continuous curves in the simulation but hourly steps?

The day-ahead dispatch optimization was repeated using several different characteristic load profiles to understand the MG performance under various demand patterns. Three different load profiles were used and compared to the baseline NREL ReOpt profile \cite{nrel_data} used earlier in \cref{sec:dispatch_opt}, as shown in \cref{fig:load_cases}. 

\begin{figure}[htbp]
\centering
\includegraphics[width=0.6\linewidth]{load_cases.png}
\caption{\label{fig:load_cases} Different load profiles used in the simulation. The peakier profile was artificially generated by allowing hourly loads to vary by $\pm 30\%$ from the baseline value. The flatter profile was created by shifting hourly load values to better align with the $P_{RES}$ curve shown, along with some additional flattening using moving averages and exponential smoothing.}
\end{figure}

\begin{table}[htbp]
\centering
\begin{tabular}{p{7cm} p{0.1cm} p{0.1cm} p{0.1cm} p{0.1cm} p{0.1cm} p{0.1cm}}
\toprule
\multicolumn{1}{l}{\textbf{Scenario}} & \multicolumn{1}{c}{\textbf{Min obj.}} & \multicolumn{1}{c}{\textbf{COE}} & \multicolumn{1}{c}{\textbf{Emissions}} & \multicolumn{1}{c}{\textbf{DPSP}} & \multicolumn{1}{c}{\textbf{Dump}} & \multicolumn{1}{c}{\textbf{REF}} \\
\midrule
Peakier load profile & 0.3529 & 0.8443 & 0.3565 & 0.008  & 0.3196 & 0.7641 \\
Flatter profile: Load shifting only & 0.2055 & 0.6165 & 0.1539 & 0.004 & 0.1355 & 0.8824  \\
Flatter profile: Load shifting and curtailment & 0.1527 & 0.3573 & 0.0524 & 0.010 & 0.3042 & 0.9607  \\
\bottomrule
\end{tabular}
\caption{\label{tab:dispatch_robust_load} Summary of dispatch results and performance under different possible demand profiles of the MG community.}
\end{table}

From \cref{tab:dispatch_robust_load}, it can be seen that the shape of the load curve has a strong influence on the operation of the MG. Peakier profiles with greater temporal variation over the course of the day result in a greater number of demand spikes that degrade the performance of the system and reduces optimality. As seen from \cref{fig:dispatch_load1}, the DE needs to be ramped up quickly to meet sudden increases in load and the uneven pattern also causes more RES output to be curtailed. The increased mismatch between supply and demand worsens all five objectives. 

On the other hand, flatter load profiles are more optimal and result in improved performance across the board, improving all objectives relative to the baseline profile. In \cref{fig:dispatch_load2}, significant improvements are achieved simply by shifting power demands so that their time of use matches up better with the availability of solar PV and wind resources during the day. This enables fuller and more efficient utilization of renewable power while minimizing curtailment and dumped power. In this case, load shifting is done in a way that also results in peak-shaving and a more uniform, averaged load curve. Furthermore, load shifting when combined with even small amounts of curtailment (e.g., 10\%) can improve results further, as seen in \cref{tab:dispatch_robust_load} and \cref{fig:dispatch_load3}, where the use of the DE is minimized. It remains online throughout the day but operates only at its minimum possible output level (2.4 kW), essentially functioning as a baseload while most of the demand is met by renewables. This results in even lower costs, emissions and higher REF than with just load shifting, although the dump energy ratio also rises quite a bit. These results motivate the need for demand response programs to shape load curves into being flatter and more optimal. Here, MG users shift and/or reduce their power consumption at specific times either voluntarily, or are incentivized to do so by the operator. Energy efficiency programs could also produce a similar, desirable effect. Finally, comparing \cref{tab:dispatch_robust_RES} and \cref{tab:dispatch_robust_load}, it may be concluded that the system is more sensitive to supply-side variations than demand changes, since the former causes larger changes in the value of the weighted objective cost function. 

\begin{figure}[htbp]
  \centering
  \begin{subfigure}[b]{0.8\linewidth}
    \includegraphics[width=\linewidth]{Dispatch_DG_8kW_OldSizing_LoadPeaky.png}
    \caption{\label{fig:dispatch_load1} Peakier load profile with higher variability, more peaks and valleys than baseline.}
  \end{subfigure}
  \begin{subfigure}[b]{0.8\linewidth}
    \includegraphics[width=\linewidth]{Dispatch_DG_8kW_OldSizing_LoadShift.png}
    \caption{\label{fig:dispatch_load2} Flatter load profile using only demand shifting i.e., the total amount of daily load served (area under the curve) is conserved.}
  \end{subfigure}
  \begin{subfigure}[b]{0.8\linewidth}
    \includegraphics[width=\linewidth]{Dispatch_DG_8kW_OldSizing_LoadShift_curtail.png}
    \caption{\label{fig:dispatch_load3} Flatter profile with load shifting along with 10\% daily curtailment.}
  \end{subfigure}
  \caption{\label{fig:dispatch_load} Hourly power dispatch curves over one day period under different demand scenarios.}
\end{figure} 

\section{Discussion and limitations\label{sec:discussion}}

\subsection{Site-specific and transferable findings}
Some findings are specific to the Timbila, Kenya case study and reflect the region's conditions. The PV--WT balance of the optimal mix follows from the local resource: 2018 irradiation at the site averaged 4.4~kWh/m$^2$/day in the station record (5.3~kWh/m$^2$/day in NASA POWER), while the wind resource is moderate and uncertain (\cref{app:climate}). In windier locations, such as coastal East Africa or the Sahel, the balance would shift toward wind. The strong sensitivity of cost to diesel price reflects the elevated and volatile fuel costs of remote sub-Saharan sites; where fuel is cheaper or more stable, the economic case for high renewable penetration weakens.

Other findings depend on the structure of the problem rather than on the site. The design workflow (equal-budget multi-run benchmarking, Pareto search, rule-based selection from the front and Monte Carlo re-ranking) applies to any simulation-based sizing problem once the weather, load and cost inputs are replaced. The solver ranking in \cref{sec:results} is an empirical result for this problem class at a budget of 2,000 evaluations; other formulations or budgets may rank solvers differently, which is why we report the full protocol. The trade-off between generator rating and renewable and storage capacity, and the cost of tighter reliability targets traced by the $\varepsilon$-constraint front, hold for any islanded system that relies on a dispatchable backup unit. Practitioners at other off-grid sites can reuse the workflow and these relationships after recalibrating irradiance, wind, load, fuel price and financing inputs.

\subsection{Limitations\label{sec:limitations}}
Five modelling choices limit these conclusions. First, Stage~2 schedules the generator and battery with exact knowledge of the next 24 hours of load and weather. The dispatch costs, emissions and DPSP in \cref{sec:dispatch_opt} are therefore optimistic relative to a controller that acts on forecasts. Stage~1 does not share this bias, because its rule-based controller uses only current-hour information. Second, sizing and dispatch are decoupled (\cref{sec:opt_prob}): a design ranked best under rule-based operation may rank differently under optimised dispatch, although its Stage-1 cost bounds its cost under any better operating policy. Third, sizing uses one meteorological year (2018) and one load profile. The Monte Carlo analysis perturbs both, but it cannot represent multi-year wind droughts or long-term load growth. Fourth, the wind resource rests on station data scaled by a single factor (\cref{app:climate}); \cref{sec:sizing_sensitivity} shows how the design responds when that factor changes by $\pm$25\%, and on-site hub-height measurements would be needed before construction. Fifth, all solvers run with default settings; tuned settings could change the margins between them, although the equal evaluation budget removes the effect of differing stopping rules.

\section{Conclusions}

This study considered the optimal design and energy management of islanded microgrids serving remote rural areas lacking grid connection and external power transfers. The hybrid MG combines renewable sources with battery storage and backup fossil-fuel generators. By simulating a system located in Timbila, Kenya over a one year period using hourly time steps while assuming perfect foresight over climate and loads, optimal capacities for the PV, WT, and BS components were obtained. This was done for a DG size based on the community's peak daily power demand, to ensure reliability. A multiobjective optimization was performed that considered several goals including economic cost and emissions minimization as well as reliability and efficiency maximization. In order to globally solve the sizing and dispatch optimization problems, several different types of solution algorithms were tested. Particle swarm optimization was found to perform best in terms of both runtime and optimality of the resulting solution. More broadly, biologically inspired techniques like GA and PSO seem to be well-suited to this class of problems, in agreement with the literature. 

For the location chosen, the combination using LI as the BS and DE as the backup DG was found to be more optimal than configurations including LA and/or MT. There exists a whole Pareto-set of input variables, all of which achieve the same global optimum. These were used to extract relationships between different input variables as well as output objectives. For instance, PV and WT act as substitutes relative to each other while increases in RES capacity will likely need to be accompanied by increased investment into BS, to account for intermittency and associated uncertainty in renewables. Similarly, there are trade-offs between pursuing different objectives in the MG. For example, more aggressive emissions reductions or reliability improvements raise system costs. The results show that RES can provide reliable and secure power for an islanded microgrid without external grid support but requires either dispatchable, fossil fuel-based DGs as backup or sufficiently large battery storage, preferably both. In order to ensure high reliability, RES systems often need to be oversized in relation to the community's power demand. This can greatly increase the initial capital investment needed and high upfront costs may discourage the community from setting up the MG. However, the break-even distance analysis confirms that the optimized microgrid is still the most well-suited design choice for the chosen location, being much cheaper than grid extension. Decentralized RES is indeed cheaper than grid extension for areas with low demand and high connection costs, a finding that agrees with the literature. Furthermore, the optimal system has high renewable penetration and much lower emissions and costs than a system relying solely on fossil fuels, while still being efficient and meeting load shortage constraints.

According to the sensitivity analysis, MG design is affected significantly by variations in external parameters like fuel prices and financial factors, and the sizing of the RES and BS components depends heavily on the rated power of the corresponding backup DG used. The dispatch optimization produced scheduling strategies that differ slightly from the load-following strategy used for the sizing stage. It was found that the MG is relatively robust to uncertainty and disturbances in load and RES output, but only with a sufficiently large DG capacity. Low RES availability (especially for the wind resource) and/or undesirable demand patterns degrades system performance since it resorts to mainly using fossil fuels via the DG. The performance is worsened further if such conditions persist for longer periods like wind lulls lasting several days, during which BS resources could also be completely drained. These weaknesses of the current model also motivate the need to incorporate demand-side management schemes and more sophisticated, predictive control to improve the MG's robustness.


\subsection{Future work \label{sec:future}}

The limitations in \cref{sec:limitations} point to three extensions: integrated or bilevel sizing and dispatch, in which the design is scored under optimised operation; receding-horizon or stochastic dispatch that acts on forecasts rather than perfect foresight; and multi-year resource records with on-site wind measurements. There are several further model refinements possible that could improve the system's performance under non-ideal or unforeseen situations, including advanced adaptive and predictive control techniques, more accurate forecasting of load and RES availability as well as demand-response programs to shape and/or curtail load. These are discussed further in \cref{app:improvements}. The current work looks at microgrids producing only electric power output. However, the modeling approach could be extended to the design and control of Combined Heat \& Power (CHP) and Combined Cooling, Heating \& Power (CCHP) microgrids as well i.e., those with thermal loads in addition to electric ones. Finally, future research could also focus on modeling other possible MG designs, particularly those that implement energy storage and backup generation options apart from those considered here. Only electrochemical energy storage in the form of batteries was included in this MG. However, there are several other technologies that could be used including electrical, mechanical, thermal, and magnetic storage methods \cite{ESS_control}. Another attractive option for rural Africa is small-scale pumped hydro or micro-hydropower schemes. Similarly, there are also non-fossil fuel-based backup generators like Hydrogen fuel cells and thermal power using biomass, which is the most abundant renewable resource available in many African countries \cite{usea}.

\section{Data and code availability}

All code and data used for this study are in the public GitHub repository \url{https://github.com/vineetjnair9/microgrid_optimization}. Release \texttt{v2.0-asoc} (branch \texttt{asoc-revision}, folder \texttt{revision\_ASOC}) contains the corrected simulation model, the benchmark, Pareto, sensitivity, Monte Carlo and dispatch scripts, the random seeds and all result files used in this revision.

\section{Acknowledgements}
The author would like to acknowledge Prof. Ioannis Lestas at the University of Cambridge for several useful discussions and guidance on the project. The author was supported by a Gates Cambridge Scholarship. There are no conflicts of interest to report.

\newpage

\appendix
\crefalias{section}{appendix}
\crefalias{subsection}{appendix}
\crefalias{subsubsection}{appendix}

\section{Microgrid system models \label{app:mg_models}}

\subsection{Renewable energy sources (RES)}

The renewable power generation sector today is heavily dependent on (1) biomass (such as firewood, charcoal etc.) which often has adverse implications for vegetation and biodiversity despite being carbon-neutral, and (2) large hydroelectric schemes which negatively affect local ecosystems and displace communities. Furthermore, bioenergy is land-intensive, leads to deforestation and often competes with crop cultivation for food purposes, while most of the geographical sites well-suited to hydropower have already been developed. This is especially true in East African countries like Kenya and Tanzania which are the focus of this study. Furthermore, there have recently been recurring droughts in the region which affect the head (level of water in reservoirs), flow rate, and thus the amount of hydropower available \cite{usea}. Thus, these two major sources were excluded from this study. Wind turbines and solar photovoltaics were the only two renewables considered. 

\subsubsection{Solar PV arrays}

The hourly DC power output [$W$] of the PV generator is calculated as $P_{PV} = n_s \; \eta_{PV} \; A_c \; I_{PV} $\cite{tazvinga2013}, where $n_s$ is the number of PV modules in the system, $\eta_{PV}$ is the power conversion efficiency [\%] of each module, $A_c$ is the collector area of each PV array [$m^2$] and $I_{PV}$ [$kW/m^2$] is the hourly solar irradiation incident on the array i.e., the global horizontal irradiance (GHI). PV module efficiency varies as a function of incident radiation $I_{PV}$, ambient outside air temperature $T_a$, and test parameters of the solar cell which are measured at nominal operating cell temperature (NOCT) and standard test conditions \cite{tazvinga2013}:

\begin{equation}
\label{eq:pv2}
    \eta_{PV} = \eta_r \eta_{pc} \cdot \left(1 - 0.9 \beta \left(\frac{I_{PV}}{I_{PV,NOCT}}\right)(T_{c,NOCT} - T_{a,NOCT}) -\beta(T_a - T_r)\right)
\end{equation}

\noindent where $\eta_r$ is the PV efficiency measured at reference cell temperature $T_r$ and standard test conditions. $\beta$ is the temperature coefficient of cell efficiency which tracks degradation in cell performance with rising temperature. $I_{PV,\;NOCT}$ is the average solar radiation incident on the array per hour at NOCT while $T_{c,\;NOCT}$ and $T_{a,\;NOCT}$ are the cell and ambient external temperatures at NOCT test conditions, respectively. Maximum power point tracking (MPPT) via the power conditioner is also used to extract maximum possible power from the system. The power conditioning efficiency $\eta_{pc}= 1$ since it is assumed that a perfect MPPT controller is used \cite{kaabeche2011b}. This system used Mitsubishi Electric PV modules with the parameters shown in \cref{tab:pv_cell}. These are monocrystalline Silicon cells with a power rating of 255 $W$ each. 

\begin{table}[htbp]
\centering
\begin{tabular}{@{}cccccccc@{}}
\toprule
$\eta_r$ [\%]  & $T_r$ [$\degree$ C] & $I_{PV,NOCT}$ [$kW/m^2$]  & $T_{c,NOCT}$ [$\degree$ C]  & $T_{a,NOCT}$ [$\degree$ C] & $\beta$ [\%/$\degree$ C] & $P_{PV,\;r}$ [$kW$] & $A_c$ [$m^2$] \\ \midrule
15.4 & 25 & 0.8 & 45.7 & 20 & 0.0045 & 0.255 & 1.4602
\\ \bottomrule
\end{tabular}
\caption{Specifications of the 255 Wp PV-MLU255HC model \cite{mitsubishi}.}
\label{tab:pv_cell} 
\end{table}

\subsubsection{Wind turbines (WT)}

A small-scale on-shore horizontal-axis wind turbine system from Raum Energy rated at 3.5 kW each was used. This is a 5-blade downwind system. A low power rating was chosen due to the relatively low energy requirements of the rural microgrid community and the lower installation and permitting fees associated with such compact systems. This results in lower upfront investment and initial capital costs. The technical specifications are listed in \cref{tab:wt_spec}. The power output of the WT depends on several factors including the wind speed pattern at the particular location, air density $\rho_{air}$, swept area $A_{swept}$ and the efficiency $\eta_w$ with which the rotor converts the kinetic energy of wind to electrical energy \cite{tazvinga2014}. The power output can be computed using:

\begin{equation}
\label{eq:Pwind}
    P_{WT} = n_w \; \frac{1}{2} \; \eta_w \; \rho_{air} \; C_p \; A_{swept} \; v_{hub}^3, \quad v_{hub}(t) = v_{ref}(t) \; \left(\frac{h_{hub}}{h_{ref}}\right)^\beta
\end{equation}

\noindent where $n_w$ is the total number of wind turbines and $v_{hub}$ is the wind speed at the rotor hub (at height $h_{hub}$). This can be estimated based on the hourly wind speed $v_{ref}$ measured at reference height $h_{ref}$ close to ground level. $\beta$ is the power law exponent or roughness coefficient which ranges from less than 0.1 for flat land, water, or ice to more than 0.25 for heavily forested regions \cite{kaabeche2011a}. A value of 0.14 was used in this study which is appropriate for wind farm installations on flat surfaces like open-terrain grasslands (such as those found in Kenya and Tanzania) away from tall trees and buildings (as is the case with rural areas) \cite{kaviani2009}. For this work, a variable speed permanent magnet generator was used due to its capability to maintain the system at its optimal designed tip-speed ratio ($\lambda$) that maximizes the power coefficient $C_p$, generator conversion efficiency and thus also the power output produced $\eta_w$ \cite{seeling1997}. 

\begin{table}[htbp]
\centering
\begin{tabular}{@{}cccccccc@{}}
\toprule
\textbf{Parameter} & $h_{hub} \; [m]$ & $h_{ref} \; [m]$ & $P_{WT,r} \; [kW]$ & $A_r \; [m^2]$ & $v_c \; [m/s]$ & $v_r \; [m/s]$ & $v_f \; [m/s]$ \\ \midrule
\textbf{Value}    & 14.5              & 1                  & 3.5                    & 12.6                                                & 2.8                & 11                                         & 22 \\ \bottomrule
\end{tabular}
\caption{Specifications of the small-scale WT system \cite{raum}.}
\label{tab:wt_spec}
\end{table}

However, $\eta_w$ and $C_p$ were not readily available from the manufacturer's data sheet and neither was their variation with wind speed given for this specific model. Thus an empirical model was used to calculate power output based on $P_{WT,\;r,\; total}$ (= $n_w \cdot P_{WT,\;r}$) and $v_{hub}$ by dividing the WT power output curve into three characteristic regions: 

\[
 P_{WT}(v_{hub}) = 
  \begin{cases} 
   n_w \cdot P_{WT,\;r} \cdot (A + B\cdot v_{hub} + C \cdot v_{hub}^2) & \text{if } v_c \leq v_{hub} \leq v_{r} \\
   n_w \cdot P_{WT,\;r}      & \text{if } v_{r} \leq v_{hub} \leq v_{f} \\
   0 & \text{otherwise}
  \end{cases}
\]
where $v_c$, $v_r$, and $v_f$ are the cut-in (start-up), rated, and cut-out (failure) wind speeds, respectively. The constants $A$, $B$ and $C$ can be computed using \cite{wen2009}:

\begin{align*}
    A & = \frac{1}{(v_c - v_r)^2} \left(v_c(v_c + v_r) - 4 v_c v_r \left(\frac{(v_c + v_r)}{2v_r}\right)^3 \right), \quad B = \frac{1}{(v_c - v_r)^2} \left(4 (v_c + v_r) \left(\frac{(v_c + v_r)}{2v_r}\right)^3 - (3v_c + v_r)\right) \\
    C & = \frac{1}{(v_c - v_r)^2} \left(2 - 4 \left(\frac{(v_c + v_r)}{2v_r}\right)^3 \right)
\end{align*}

\subsubsection{Climate data \label{app:climate}}

In order to estimate solar and wind power generation, historical temperature, wind speed, and solar irradiance data were obtained at an hourly resolution for an entire year (2018) from ground-based weather stations which tend to be more accurate than satellite measurements, on average. These hourly averaged values were made available by the Trans-African HydroMeteorological Observatory (TAHMO), a public benefit organization that maintains a vast network of weather stations across Africa \cite{tahmo}. The microgrid was designed for a rural village near Timbila in southern Kenya close to the Tanzanian border as shown in \cref{fig:chap3_timbila}. This location was chosen due to the required weather and load data being available more easily. However, several of the wind speed values were missing in the TAHMO datasets. These values were approximated by averaging the wind speed from four other neighboring TAHMO stations located close to Timbila, specified in \cref{tab:tahmo_stations}.

\begin{figure}[htbp]
  \centering
  \includegraphics[width=0.6\linewidth]{timbila.png}
  \caption{Location of proposed microgrid site in Timbila, Kenya.}
  \label{fig:chap3_timbila}
\end{figure}

\begin{table}[htbp]
\centering
\begin{tabular}{p{9cm} p{1cm} p{1cm} p{1cm} p{1cm}}
\toprule
\multicolumn{1}{l}{\textbf{Location}}    & \multicolumn{1}{c}{\textbf{Lat (\degree N)}} & \multicolumn{1}{c}{\textbf{Lon (\degree E)}} & \multicolumn{1}{c}{\textbf{Elevation (m)}} \\ \midrule
Timbila High School, Kenya                                        & -3.39                                     & 37.72                                      & 743                                        \\
Imurtott Primary School, Kenya                                      & -2.88                                      & 37.39                                       & 1853                                       \\
South Eastern Kenya University - Mtitu Adei Campus             & -2.66                                     & 38.11                                      & 854                                        \\
Our Lady of Perpetuah Girls Secondary School, Kenya                   & -3.50                                       & 38.59                                      & 1158                                       \\
Bangalala, Tanzania                                                 & -4.23                                     & 37.85                                      & 877                                        \\ \bottomrule
\end{tabular}
\caption{Geographical details of stations used to generate combined climate data \cite{tahmo}.}
\label{tab:tahmo_stations}
\end{table}

The TAHMO wind speeds at the Timbila station average only 0.86~m/s, far below typical values for the area. The original analysis therefore scaled the series by a factor of about 3.70 (= 3.19/0.86) so that its annual mean matched forecast-website values for Timbila. For the revision we checked the series against an independent source, NASA POWER hourly data for 2018 at the same coordinates \cite{nasapower}, whose wind speeds come from the MERRA-2 reanalysis \cite{gelaro2017merra2}. After a 2-hour time alignment (chosen by maximizing the correlation of irradiance, $r = 0.96$), the two wind series agree in timing ($r = 0.63$ hourly, 0.60 for daily and 0.82 for monthly means; \cref{fig:wind_val}) but not in level: the reanalysis 2~m wind averages 2.23~m/s, 2.59 times the TAHMO mean, and monthly ratios range from 2.07 to 3.33 (\cref{tab:wind_val}). At the 14.5~m hub height, the scaled TAHMO data with the model's 1~m reference height give a mean of 4.68~m/s, while the reanalysis 10~m wind extrapolated with its own 10--50~m shear gives 3.50~m/s, 25\% lower. A 2~m reference height would lower the model's hub-height mean by 9\% (4.24~m/s). Reanalysis data represent a grid cell of about 50~km in complex terrain near Mount Kilimanjaro rather than the site itself, so neither source is definitive; on-site measurement at hub height would be needed before construction. We keep the original scaling for comparability and report the sensitivity of the design to the wind level in \cref{sec:sizing_sensitivity}: scaling the wind input by 0.75, which matches the reanalysis hub-height mean, changes the optimal design from [32, 4, 107] to [28, 5, 77] and raises LCOE from 0.103 to 0.136~USD/kWh. The TAHMO irradiance agrees closely with the CERES irradiance in NASA POWER in timing ($r = 0.96$) and is 16\% lower on average (4.42 vs 5.28~kWh/m$^2$/day), so the PV results are, if anything, conservative.

\begin{figure}[htbp]
  \centering
  \includegraphics[width=\linewidth]{rev_wind_validation.pdf}
  \caption{\label{fig:wind_val} TAHMO station wind (raw and scaled by 3.70) compared with NASA POWER (MERRA-2) wind at 2~m and 10~m for Timbila, 2018: (a) monthly means, (b) mean diurnal profiles after time alignment.}
\end{figure}

\begin{table}[htbp]
\centering
\scriptsize
\setlength{\tabcolsep}{3pt}
\input{rev_wind_monthly.tex}
\caption{\label{tab:wind_val} Monthly mean wind speed of the raw TAHMO series and NASA POWER at 2~m, and their ratio (2018).}
\end{table}  Sample hourly profiles of solar irradiance, wind speed, and ambient air temperature are shown in \cref{fig:chap3_hourly} for TAHMO data recorded on January 1, 2018. Daily averaged profiles are shown in \cref{fig:chap3_daily}.

\begin{figure}[htbp]
  \centering
  \begin{subfigure}[b]{0.32\linewidth}
    \includegraphics[width=\linewidth]{daily_solar_example_tahmo.png}
  \caption{Solar irradiance.}
  \end{subfigure}
  \begin{subfigure}[b]{0.33\linewidth}
    \includegraphics[width=\linewidth]{daily_wind_example_tahmo.png}
    \caption{Wind speed.}
  \end{subfigure}
  \begin{subfigure}[b]{0.33\linewidth}
    \includegraphics[width=\linewidth]{daily_temp_example_tahmo.png}
    \caption{Temperature profile.}
  \end{subfigure}
  \caption{Examples of hourly profiles of climate variables over a characteristic day.}
  \label{fig:chap3_hourly}
\end{figure}

\begin{figure}[htbp]
  \centering
  \begin{subfigure}[b]{0.32\linewidth}
    \includegraphics[width=\linewidth]{annual_irradiance_tahmo.png}
  \caption{Solar irradiance.}
  \end{subfigure}
  \begin{subfigure}[b]{0.33\linewidth}
    \includegraphics[width=\linewidth]{annual_wind_tahmo.png}
    \caption{Wind speed.}
  \end{subfigure}
  \begin{subfigure}[b]{0.33\linewidth}
    \includegraphics[width=\linewidth]{annual_temp_tahmo.png}
    \caption{Ambient external temperature.}
  \end{subfigure}
  \caption{Daily averaged profiles of climate variables over the year 2018.} \label{fig:chap3_daily}
\end{figure}

\subsection{Electricity demand (load)}

Characteristic rural African load profiles were obtained using NREL's ReOpt tool \cite{nrel_data}. The assumptions made regarding the MG community (small, remote village) are listed in \cref{tab:load}. For simplicity, it was assumed that there is no significant variation in load profiles between weekdays and weekends or from season to season. These assumptions are reasonable given the tropical climate in Timbila with similar weather conditions all year round, and the relatively low penetration of commercial entities and businesses in the largely residential village. Household, commercial and combined demand profiles are shown for one characteristic day (24 h period) in \cref{fig:chap3_load}. Since the sizing optimization was conducted over a horizon of one year, the daily profile was used to artificially generate an annual load profile giving hourly power demands for all 365 days of the year. This was done by setting the profile for January 1, 2018 to that shown in \cref{fig:chap3_load} and then allowing hourly demand in subsequent days to randomly vary by $\pm 20\%$. By introducing random variability in day-to-day (instead of hour-to-hour) fluctuations, the load curves in other days still maintain a similar shape or pattern. Key summary statistics for the resulting daily load profile are listed in \cref{tab:load_summary}. 

\begin{table}[htbp] 
\centering
\begin{tabular}{@{}lc@{}}
\toprule
\multicolumn{1}{c}{\textbf{Parameter}} & \textbf{Value}       \\ 
\midrule
Number of households & 100 \\
\% of high income households & 33\% \\
\% of medium income households & 33\% \\
\% of low income households & 33\% \\
Number of water pumping operations & 6 \\
Number of milling operations & 4 \\
Number of small shops & 10 \\
Number of schools & 1 \\
Number of clinics & 3 \\
Number of street lights & 30 \\
\bottomrule
\end{tabular}
\caption{Description of MG community near Timbila, Kenya.}
\label{tab:load}
\end{table}

\begin{figure}[htbp]
  \centering
  \includegraphics[width=0.6\linewidth]{NREL_load.PNG}
  \caption{Rural African microgrid load profiles used in this study \cite{nrel_data}.}
  \label{fig:chap3_load}
\end{figure}

\begin{table}[htbp]
\centering
\begin{tabular}{@{}ccccc@{}}
\toprule
Mean kWh/day & Total kWh/year & Max kW/day & Min kW/day & Mean kW/day \\ \midrule
203 & 74,251 & 12.52 & 3.21 & 8.48 \\ \bottomrule
\end{tabular}
\caption{Summary statistics for annual and daily load patterns.}
\label{tab:load_summary} 
\end{table}

\subsection{Backup distributed generators (DG) \label{sec:DG}}

In order to ensure system reliability, standby distributed generators using conventional fossil fuels must provide backup power. These are only called upon to satisfy excess demand that cannot be met by renewables and battery storage. This is essential in cases of sudden jumps in demand or insufficient supply from renewable sources due to inclement weather, natural disasters etc. Similar to the wind turbines, variable speed control was used for the fossil-fueled distributed generators as well due to its lower fuel consumption compared to constant speed operation. Such control allows the engine or turbine to operate at lower speeds to meet low power demand levels and vice versa, thus maximizing the operational efficiency of the DG \cite{seeling1997}. In order to simplify the problem, it was assumed that the system uses only one DG rather than having to individually dispatch multiple units. The sizing of the diesel generator or micro turbine was not considered as a variable in the optimization problem. Instead, the rated power of the DG was set to meet the microgrid community's peak demand for electric power as $P_{DG,\;r} = P_{load,\;peak} \cdot (1 + SF)$ \cite{ghasemi2013}, where $SF$ is a safety factor that determines the extent to which the DE or MT is oversized. This excess backup power is considered as spinning or operating reserves since it is the extra capacity obtained by increasing the rotational speed of the synchronous generator.

\subsubsection{Diesel engines (DE)}

These are traditional reciprocating, internal combustion engines based on compression ignition that diesel as their fuel. These offer several advantages including their dispatchable nature (unlike RES), quick startup and load-following ability but also have demerits such as high fuel costs, pollutant emissions including greenhouse gases and particulates and noise generation \cite{mg_review}. The load peaks at 12.52~kW on the representative day used for dispatch and at 15.0~kW over the year (26 June, 11:00). The DE is rated at 16~kW, 28\% above the representative-day peak ($SF \approx 0.278$) and 6.6\% above the annual peak. According to \cite{fathima2015}, the lower limit on the DE's power output can be estimated as $P_{DE,\;min} = 0.3 \cdot P_{DE,\; r} = 4.8 \; kW$. In every hour in which the DE is committed ($DE_{ON}(t) = 1$), its volumetric fuel consumption follows the linear law $Fuel_{DE}(t) = \alpha_{DE} P(t)_{DE,\;gen} + \beta_{DE} P(t)_{DE,\;r}$, where $P(t)_{DE,\;gen}$ [in $L/h$] is power output level needed by the load \cite{kaabeche2014}, $\alpha_{DE}$ [$L/kWh$] and $\beta_{DE}$ [$L/kWh$] are fuel consumption curve coefficients supplied by the manufacturer, assumed to be 0.246 and 0.08145 respectively \cite{kaabeche2014}. The hourly fuel costs of the DE can be calculated assuming a unit price of \$ 3.20 per US liquid gallon of diesel, as $C_{fuel,\; DE}(t) = 3.20 \cdot \frac{Fuel_{DE}(t)}{3.78541}$. The DE consumes no fuel in hours when it is off.

\subsubsection{Micro-gas turbines (MT)}

Another option for the backup DG is a microturbine, which can run on a diverse range of gaseous and liquid fuels. These offer similar advantages to DEs with the additional benefits of being compatible with combined heat and power (CHP) applications. For CHP purposes, natural gas is the most common fuel. Furthermore, these have lower greenhouse gas emissions than DEs and much less than larger, centralized natural gas plants, especially for $NO_x$ and $CO_2$ exhausts \cite{wu2014}. However, microturbines are more expensive upfront with higher initial capital costs and often higher fuel costs as well when compared to diesel. They are also not nearly as prevalent across the African continent as diesel gensets currently are. Microturbines tend to be larger on average than diesel engines, generally ranging from 30 to 330 kW. However, for direct comparison with the diesel alternative, the same size was chosen for the MT as well, with $P_{MT,\; r} = P_{MT, \; max} = 16 \; kW$ and $P_{MT,\;min} = 4.8 \; kW$. According to \cite{chp_mt}, the fuel consumption of a 61 kW MT is 0.84 MMBtu per hour of operation at its fully rated power. Extrapolating from this, the hourly fuel consumption of the MT [in $MMBtu/h$] can be approximated using the output level $P_{MT}$, as $Fuel_{MT}(t) = \frac{0.84}{61} \cdot P_{MT}(t)$. The hourly fuel consumption cost of the MT can then be computed assuming a unit price of $C_{NG} = \$ 2.19/MMBtu$ for natural gas \cite{ng_price}, as $C_{fuel,\;MT}(t) = 2.19 \cdot Fuel_{MT}(t)$.

\subsubsection{Converter}

A bidirectional converter was used that is capable of acting as both an inverter and rectifier, with a power rating equal to the peak load (12.52 kW). The efficiencies of both AC-DC and DC-AC conversion were assumed to be equal i.e., $\eta_{inv} = \eta_{rec} = 90\%$ \cite{ogunjuyigbe2016} and the lifetime was assumed to be 20 y \cite{moshi2016}. The capital cost of the converter was taken to be $IC_{conv} = \$2800$ \cite{conv_price}.

\subsection{Battery energy storage (BS) \label{sec:BS}}

Two options for electrochemical energy storage are considered: Lithium-ion ($Li^+$) and Lead-acid batteries. Although these two differ in their parameters, the same discrete time model was used to describe their charging and discharging processes \cite{tazvinga2014}:
\begin{equation}
    SOC(t+1) = (1-\delta) SOC(t) - \frac{1000 \cdot P_{BS} \; \Delta t \; \eta_{BS}(t)}{E_C}, \qquad \eta_{BS}(t) = \begin{cases} \eta_{overall} & P_{BS} < 0 \text{ (charging)} \\ 1 & P_{BS} \geq 0 \text{ (discharging)} \end{cases}
\end{equation}
where $P_{BS}$ [$kW$] is the DC power dispatched either to or from the battery in each time step $\Delta t$ and $\delta$ is the hourly self-discharge rate (\%/h). A positive $P_{BS}$ value indicates that the battery is being discharged while a negative $P_{BS}$ implies charging. $SOC$ is its state of charge which indicates the remaining available battery capacity as a percentage of its current maximum energy capacity $E_C$, i.e., $SOC = \frac{E_{releasable}}{E_C}$. Since it is difficult to measure both charging and discharging efficiencies separately, an overall round-trip Coulombic efficiency $\eta_{overall}$ for the battery is used and applied once, when charging. In each hour the controller limits charging to the headroom below $SOC_{max}$ and discharging to the energy above $SOC_{min}$, in addition to the power rating $P_{BS,\;r}$. The power rating scales with installed capacity, as for stacked modules: 3.68~kW per 13.5~kWh Powerwall module for Li-ion (0.273~kW per kWh) and a C-rate of 0.2 for lead-acid. Key properties of both batteries are listed in \cref{tab:bs_params}:

\begin{table}[htbp]
\centering
\begin{tabular}{@{}lll@{}}
\toprule
\textbf{Parameter}               & \textbf{Li-ion (LI)} & \textbf{Pb-acid (LA)} \\ \midrule
$SOC_{min}$                         & 10\%            & 50\%             \\
$SOC_{max}$                         & 90\%            & 90\%             \\
$\delta$ [\%/month] & 7.5             & 5                \\
$\eta_{overall}$     & 90\%            & 75\%               \\
Lifetime [in y]                  & 15              & 5                \\
Lifetime $t_{BS}$ [in cycles]             & 5475            & 1400             \\
$P_{BS, \; r}$ [kW per kWh installed]  & 0.273           & 0.20             \\
Capacity fading $\Delta_E$ [\%/cycle]       & 0.0055          & 0.0214           \\ \bottomrule
\end{tabular}
\caption{Parameters of the Li-ion and Lead-acid BS systems \cite{bs_source}.}
\label{tab:bs_params} 
\end{table}

The battery's state of charge must remain within the upper and lower limits $SOC_{max}$ and $SOC_{min}$ respectively, throughout its operation. In the simulation, the BS was assumed to be fully charged at the start ($t = 0$) i.e., $SOC = SOC_{max} = 90\%$ for both the LI and LA systems. The lower limit is determined by the maximum level until which the battery can be safely discharged, known as the depth of discharge $DoD$ i.e., $E_{b,\;min} = (1 - DoD)\cdot E_C$. The maximum capacity of the battery $E_C$ is initially equal to the original rated capacity $E_{b,\; init}$ but then gradually degrades linearly with increased cycling in a process known as capacity fading, given by $E_C = E_{b,\; init} \cdot (1 - n_{cycles}\cdot \Delta_E)$. Here, $n_{cycles}$ is the cumulative number of charge-discharge cycles so far, a variable that is updated at each iteration in the dispatch algorithm. Since the system is simulated over a horizon of one year, this gives the total number of cycles per year. Capacity fading values for both Li-ion and Pb-acid were calculated assuming that the battery is replaced once $E_C$ has reduced to 70\% of its original rated value from \cref{tab:bs_params}. %which is also when the sum total of $n_{cycles}$ over all years in the MG lifetime equals $t_{BS}$. 
Some of the Li-ion parameters were updated using the latest values for the Tesla Powerwall \cite{powerwall}, assuming that the battery can be fully cycled daily for five years beyond the specified 10-year warranty period \cite{bs_source}.

\cref{tab:bs_params} clearly shows trade-offs in terms of the relative advantages and disadvantages of the two options. Lead-acid batteries have lower lifetimes (requiring more frequent replacement) and lower round-trip efficiencies than Li-ion. Li-ion also allows deeper discharge enabling the battery capacity to be utilized more fully, while the much higher peak power allows for faster discharge to meet more of the demand. However, it has a slightly higher self-discharge mainly owing to the parasitic protection circuit, and is currently more expensive, requiring greater initial investment than Lead-acid. However, the slower capacity fading means it needs to be replaced less often, which could make it more economical depending on the manner in which the battery is dispatched and how often it is cycled during operation. 

\section{Details of cost functions in the objective \label{app:costs}}

Since this problem involves a multi-period optimization over the entire project lifecycle ($t_{sys} = 25 y$), it is necessary to discount all cash flows to the present. \cref{tab:financial} lists the financial assumptions.

\begin{table}[htbp]
\centering
\begin{tabular}{@{}lcc@{}}
\toprule
\multicolumn{1}{c}{\textbf{Parameter}} & \textbf{Symbol}      & \textbf{Value}       \\ \midrule
Nominal interest or discount rate (\%) & $i$                    & 9 \cite{interest_rate}                   \\
Inflation or escalation rate (\%)                    & $f$                    & 5.7 \cite{inflation_rate}                 \\
System lifetime (y)                    & $t_{sys}$           & 25   \\
PV lifetime & $t_s$ & 25 \\
WT lifetime & $t_w$ & 20 \\
\bottomrule
\end{tabular}
\caption{Key financial parameters used in the model.}
\label{tab:financial} 
\end{table}

\subsection{Sizing optimization problem costs}

The initial capital (IC) investment for the RES sources (PV and WT) and backup DGs are calculated using installed cost values per unit of rated power, while the upfront expenditure for the batteries is calculated using the cost per unit of rated energy capacity. Replacement costs (RC) are estimated to be a smaller fraction of IC since it excludes one-time fees like permitting and installation. RES sources and batteries were assumed to only have fixed annual operation and maintenance (O\&M) costs calculated as a fraction of one-time capital costs, while the converter was assumed to have zero O\&M costs. However, the DGs have running costs in addition to fixed O\&M, for fuel consumption, variable O\&M (depending on its number of operational hours) as well as startup $SUC$ (to bring the DG online) and shutdown $SDC$ (to turn off the DG) costs. These are summarized in \cref{tab:costs1} and \cref{tab:costs2}. 

\begin{table}[htbp]
\centering
\begin{tabular}{@{}lcccc@{}}
\toprule
\textbf{Parameter}        & PV   & WT   & Li-ion & Pb-acid \\ \midrule
$IC_{kW}$ [$\$/kW$]           & 1210 & 1500 & -      & -       \\
$IC_{kWh}$ [$\$/kWH$]           & -    & -    & 300    & 255     \\
$RC$ [\% of $IC$]         & -    & -    & 90     & 90      \\
$OM_{fix}$ [\% of $IC$/y] & 1    & 3    & 1      & 1       \\ \bottomrule
\end{tabular}
\caption{Costs assumed for the RES and BS.}
\label{tab:costs1}
\end{table}

\begin{table}[htbp]
\centering
\begin{tabular}{@{}lcc@{}}
\toprule
\textbf{Parameter} & DE \cite{moshi2016}    & MT \cite{chp_mt}    \\ \midrule
$IC_{kW}$ [$\$/kW$]     & 781.25 & 3320  \\
$RC$                 & 88     & 90    \\
$OM_{fix}$ [\% of $IC$/y]          & 2      & 2     \\
$OM_{var}$            & 0.24 [$\$/h$]   & 0.013 [$\$/kWh$] \\
$SUC$ [\$]                & 0.45   & 0.45  \\
$SDC$ [\$]               & 0.23   & 0.23  \\
Lifetime $t_{DG} \; [h]$ [h]  & 15000  & 40000 \\ 
\bottomrule
\end{tabular}
\caption{Cost assumptions for the DGs.}
\label{tab:costs2}
\end{table}

The capital costs for WT and PV were assumed to be globally weighted averages of installed capital costs \cite{costs_irena} while the fixed O\&M costs were taken from \cite{kaabeche2011a} and the battery values from \cite{hu2017}. Note that the assumed initial investment costs for WT and PV are likely to be slight underestimates since these average values also consider much larger, utility-scale projects (of the order of MW) which tend to have lower installed costs than small-scale installations like those in MG, due to economies of scale. Replacement was considered only for components with the shortest lifetimes i.e., the BS and backup DGs. Since the converter and WT have much longer lifetimes that are close to the overall project lifetime, these were not included in the replacement costs. Discrete vector variables  $\vv{\bm{DG_{ON}}}(t)$ (i.e., $\vv{\bm{DE_{ON}}}(t)$ and $\vv{\bm{MT_{ON}}}(t)$) were created to keep track of the operational state (ON/OFF) of the DG in each time interval $t$, in order to calculate the total startup, shutdown and variable O\&M costs. $DG_{ON}(t)$ = 0 implies that the DG is switched off while $DG_{ON}(t)$ = 1 indicates that it has been brought online. Another discrete variable $BS_{cycled}(t)$ records whether the battery has already been cycled before (= 1) or not (= 0) in the current iteration, to avoid being self-discharged twice within the same time step during dispatch. 
The total initial capital (or installed) costs are: 
\begin{align*}
    IC_{total} = & \; IC_{PV} + IC_{WT} + IC_{BS} + IC_{DG} + IC_{conv} \\
    = & \; IC_{PV,\; kW} \cdot P_{PV,\;r,\;total} + IC_{WT,\;kW}\cdot P_{WT,\;r,\;total} + IC_{BS,\;kWh} \cdot E_{b,\; init} + IC_{DG,\;kW}\cdot P_{DG,\;r,\;total} + IC_{conv} \\
    = & \; IC_{PV,\; kW} \cdot n_s \cdot P_{PV,\;r} + IC_{WT,\;kW}\cdot n_w \cdot P_{WT,\;r} + IC_{BS,\;kWh} \cdot E_{b,\; init} + IC_{DG,\;kW}\cdot P_{DG,\;r,\;total} + IC_{conv}
\end{align*}
The total annual recurring costs are:
\begin{align*}
    C_{rec} & = OM_{fix,\;PV} + OM_{fix\;WT} + OM_{fix,\;BS} + OM_{fix,\;DG}   \\
    & + OM_{var, \; DG,\; annual} + C_{fuel,\; DG,\; annual} + SUC_{DG,\;annual} + SDC_{DG,\;annual}  \\
\end{align*}
where
\begin{align*}
    OM_{var, \; DG,\; annual} & = \Sigma_{t=1}^{t=8760} \; DG_{ON}(t) \cdot OM_{var, \; DG}, \; C_{fuel,\; DG,\; annual} & = \Sigma_{t=1}^{t=8760} \; C_{fuel,\; DG}(t). 
\end{align*}
and $SUC_{DG,\;annual}$ and $SDC_{DG,\;annual}$ are incremented by $SUC_{DG}$ and $SDC_{DG}$ respectively, whenever there is a change in the operational state of the DG (i.e., if $DG_{ON}(t) \neq DG_{ON}(t-1)$). The recurring costs incurred above in the 1st year of operation are assumed to remain the same in all years of the MG operation henceforth. The present worth (PW) of these recurring costs can be calculated using \cite{lcc,ajan2003}:
\begin{equation}
    PW_{C_{rec}} = C_{rec} \frac{\left(\frac{1+f}{1+i}\right)\left(\left(\frac{1+f}{1+i}\right)^{t_{sys}} - 1\right)}{\left(\frac{1+f}{1+i}\right)- 1}
\end{equation}
The total non-recurring costs are due to periodic replacement of the DG and BS: 
\begin{equation}
    C_{non-rec} = RC_{DG} + RC _{BS}
\end{equation}
These can also be converted to PW using \cite{lcc,ajan2003}:
\begin{equation}
    PW_{C_{non-rec}} = C_{non-rec} \frac{\left(\frac{1+f}{1+i_{adj}}\right)\left(\left(\frac{1+f}{1+i_{adj}}\right)^{t_{sys}} - 1\right)}{\left(\frac{1+f}{1+i_{adj}}\right)- 1}
\end{equation}
where $i_{adj}$ is an adjusted interest rate computed as:
\begin{equation}
    i_{adj} = \frac{(1+i)^L}{(1+f)^{L-1}} - 1
\end{equation}
where $L$ is the time period of replacement. Assuming that the system is operated in a similar manner to its first year throughout its subsequent operation, the replacement period of the DG and BS can be determined using the variables $n_{cycles}$ and $DG_{ON}$. The number of years between successive BS replacements is given by $L_{BS} = \frac{t_{BS}}{n_{cycles}}$. Furthermore, if an hourly simulation time-step is used (i.e., $\Delta t = 1 \; h$), then $DG_{ON}$ gives the total number of operational hours of the DG each year and the number of years between its replacement is $L_{DG} = \frac{t_{DG}}{\Sigma_{t=1}^{t=8760} \; DG_{ON}}$
    
The total net present cost (TNPC) can then be obtained by summing up all the lifecycle costs of the system over its lifetime and then discounting back to a common base (the present) \cite{kaabeche2011a}: 
\begin{equation}
    TNPC \; [\$] = IC_{total} + PW_{C_{non-rec}} + PW_{C_{rec}}
\end{equation}
Subsequently, the total annualized cost (TAC) can be calculated as:
\begin{equation}
    TAC = TNPC \cdot CRF 
\end{equation}
where CRF is the capital recovery factor and is used to calculate the present value of an annuity \cite{kaviani2009}:
\begin{equation}
    CRF = \frac{r(1+r)^{t_{sys}}}{(1+r)^{t_{sys}} - 1}
\end{equation}
where r is the real interest rate:
\begin{equation}
\label{eq:real_interest}
    r = \frac{i-f}{1 + f}
\end{equation}

\subsection{Dispatch optimization problem costs}

The cost calculations for day-ahead dispatch are similar to those for sizing, except that it only considers the running operational costs like daily O\&M and fuel costs incurred over a 24 h period, excluding initial capital and replacement costs. Thus, there is no need for discounting in this case and total daily costs can be readily calculated as: 
\begin{align*}
    C_{daily} & = C_{fuel,\; DG,\; daily} + OM_{var, \; DG,\; daily} + SUC_{DG,\;daily} + SDC_{DG,\;daily} \\
    & + \frac{OM_{fix,\;PV} + OM_{fix,\;WT} + OM_{fix,\;BS} + OM_{fix,\;DG}}{365} \\
    OM_{var, \; DG,\; daily} & = \Sigma_{t=1}^{t=24} \; DG_{ON}(t) \cdot OM_{var, \; DG}, \; C_{fuel,\; DG,\; daily} = \Sigma_{t=1}^{t=24} \; C_{fuel,\; DG}(t). 
\end{align*}

\subsection{Objective functions}

While optimizing the design, sizing and dispatch of the MG, there are several objectives that need to be pursued simultaneously and sometimes, these may even compete or be in conflict with one another. Thus, it is necessary to perform a multiobjective, multiperiod, constrained and bounded optimization in order to obtain a holistic and comprehensive solution while also being able to clearly visualize trade-offs between the various goals. Five distinct objectives are considered in the current optimization problem:

\subsection{Levelized cost of electricity (LCOE)}
$LCOE$ [$\$/kWh$] is the ratio of total annual costs to the total annual electrical load served:
\begin{align}
    LCOE & = \frac{TAC}{\Sigma_{t=1}^{t=8760} \; E_{load}(t)} = \frac{TAC}{\Sigma_{t=1}^{t=8760} \; P_{load}(t)\cdot \Delta t} = \frac{TAC}{\Sigma_{t=1}^{t=8760} \; P_{load}(t)}
\end{align}
where $\Sigma_{t=1}^{t=8760} \; E_{load}(t)$ is the total electrical energy [in $kWh$] consumed by the load in a year (8760 h) assuming that an hourly time step ($\Delta t = 1 \; h$) is used. Designers seek to minimize $LCOE$ during sizing optimization in order to ensure affordable electricity access. For dispatch optimization, the daily cost of electricity (COE) [$\$/kWh$] is minimized instead, given by:
\begin{align*}
    COE & = \frac{C_{daily}}{\Sigma_{t=1}^{t=24} \; P_{load} (t)\cdot \Delta t}  = \frac{C_{daily}}{\Sigma_{t=1}^{t=24} \; P_{load}(t)}
\end{align*}

\subsection{Emissions \label{sec:norm}}
In order to mitigate negative environmental impacts of the MG, it is necessary to minimize emissions of harmful airborne pollutants (including particulate matter (PM), volatile organic compounds (VOC) and greenhouse gases (GHG)) from the backup DGs. The emissions factors arising from the use of diesel in the DE and natural gas in the MT are listed in \cref{tab:emissions}. The total emissions $Em$ [in $kg$] can be obtained by summing up these factors and multiplying it by the total energy derived from fossil sources: $\Sigma_{t=1}^{t=T} \; P_{DG}(t) \cdot \Delta t$ where T = either 24 h or 8760 h, for dispatch and sizing optimization respectively.

\begin{table}[htbp]
\centering
\begin{tabular}{@{}lcc@{}}
\toprule
\textbf{Parameter} & DE & MT \\ 
\midrule
$CO_2$ [$kg/kWh$] & 0.649 \cite{wu2016} & 0.631 \cite{chp_mt} \\
$CO$ [$g/kWh$] & 4.063 & 2.851\\
$NO_x$ [$g/kWh$] & 18.857 & 20.884\\
$SO_2$ [$g/kWh$] & 0.074 & 0.003\\
$VOC$ [$g/kWh$] & 1.502 & 0.604\\
$PM$ [$g/kWh$] & 1.338 & 0.047\\
$PM2.5$ [$g/kWh$] & 1.338 & 0.047\\
$PM10$ [$g/kWh$] & 1.338 & 0.047\\
\bottomrule
\end{tabular}
\caption{Emissions factors for both DE and MT \cite{emissions}.}
\label{tab:emissions}
\end{table}

In order to normalize both $LCOE$ and emissions so that they always lie between 0 and 1 (like the other three objectives), these values are divided by the equivalent LCOE and emissions from a hypothetical, base-case scenario where the same MG community is run solely on a single backup DG (either a DE or MT chosen accordingly) of sufficient rated power (i.e., $P_{DG,\;r} = P_{load,\; peak} \cdot (1+SF)$, where $SF \geq 0$). Here the DG is run continuously throughout the study period to meet the entirety of the load. Startup and shutdown costs are incurred only once each at the beginning and end of the time horizon, respectively. There is also no need for a converter since the DG's AC output can directly supply loads. Thus, the costs for the base case (used to calculate $LCOE_{base}$) are comprised of:

\begin{gather*}
    IC_{base} = IC_{DG}, \; C_{rec,\; base} = OM_{fix,\;DG} + OM_{var, \; DG,\; annual} + C_{fuel,\; DG,\; annual} + SUC_{DG} + SDC_{DG}  \\
    C_{non-rec,\; base} = RC_{DG}, \; C_{daily,\; base} = C_{fuel,\; DG,\; daily} + OM_{var, \; DG,\; daily} + SUC_{DG} + SDC_{DG} + \frac{OM_{fix,\;DG}}{365}
\end{gather*}

\noindent where $OM_{var, \; DG,\; annual}$ and $C_{fuel,\; DG,\; annual}$ are calculated assuming the DG is run for 8760 h to meet the entire annual load. Similarly, $C_{fuel,\; DG,\; daily}$ and $OM_{var, \; DG,\; daily}$ are calculated assuming the DG is run for 24 h to meet the entire daily load. The base emissions $Em_{base}$ are also calculated assuming that the DG meets all the power demand, on a daily and annual basis. The baseline DG follows the load in every hour, so its annual fuel use is $8760\,\beta_{DE} P_{DG,\;r} + \alpha_{DE} \sum_t P_{load}(t)$; it does not run at rated power for 8,760~h. In every sensitivity case we recompute $LCOE_{base}$ with the same fuel price, interest and inflation rates as the hybrid system, so the normalized LCOE compares the two systems under identical economic conditions, while the baseline DG rating stays at 16~kW so that the normalization does not move with the hybrid system's DG rating. Because the baseline itself moves with fuel price and financial parameters, the normalized LCOE can fall while the absolute LCOE rises; \cref{sec:sizing_sensitivity} therefore reports both.

\subsection{Deficiency of power supply probability (DPSP)}
$DPSP$ is the ratio between the total unmet or deficit load that is not satisfied by the MG (i.e., consumption that is either curtailed, interrupted, or disrupted) and the total load demand in the same study period \cite{prasad2006}:
\begin{equation}
    DPSP = \frac{\Sigma_{t=1}^{t=T} \; P_{lost}(t)}{\Sigma_{t=1}^{t=T} \; P_{load}(t)}
\end{equation}
where $T = 8760\;h$ for sizing and $T = 24\;h$ for dispatch. Thus, $0 \leq DPSP \leq 1$ and a value of 1 implies that the load will never be satisfied while a value of 0 means that the load will always be satisfied. $DPSP$ must be minimized in order to maximize the system's reliability.  

\subsection{Relative excess or dump power generated (REPG)}

This is the ratio between excess power (DC) generated by the system (over the total demand) that is sent to dump loads or grounded to earth, and the total power (DC) generated \cite{kaabeche2011b}:
\begin{align}
    REPG = Dump & = \frac{\Sigma_{t=1}^{t=T} \; P_{dump}(t)}{\Sigma_{t=1}^{t=T} \; P_{gen}(t)} = \frac{\Sigma_{t=1}^{t=T} \; P_{RES}(t) + P_{DG}(t) \cdot \eta_{rec} + P_{BS}(t) - \frac{P_{load}}{\eta_{inv}}}{\Sigma_{t=1}^{t=T} \; P_{RES}(t) + P_{DG}(t) \cdot \eta_{rec}} \label{eq:Pbal}
\end{align}
where \cref{eq:Pbal} arises from the power balance in the system. $REPG$ is a measure of the system's efficiency. Thus, it is desirable to minimize this ($0 \leq REPG \leq 1$) in order to reduce the wastage of energy and keep the RES curtailment as low as possible. 

\subsection{Renewable energy penetration}
The renewable energy fraction (REF) measures the proportion of total power generation coming from clean, renewable sources:
\begin{align}
    REF & = \frac{\Sigma_{t=1}^{t=T} \; P_{RES}(t)}{\Sigma_{t=1}^{t=T} \; P_{gen}(t)} = \frac{\Sigma_{t=1}^{t=T} \; P_{PV}(t) + P_{WT}(t) \cdot \eta_{rec}}{\Sigma_{t=1}^{t=T} \; P_{RES}(t) + P_{DG}(t) \cdot \eta_{rec}}
\end{align}
where $T$ is either 8760 h or 24 h for sizing and dispatch, respectively. In order to improve the sustainability of the MG and reduce dependence on fossil fuels, it is essential to maximize $REF$ ($0 \leq REF \leq 1$) or equivalently minimize $1-REF$. 

\section{Effects of varying relative objective weights on the sizing results \label{app:weights}}

Here, the weights $w_i$ placed by the decision maker (i.e., the MG designer) on each of the five objectives in the multiobjective function were individually varied while keeping the values of the other four equal to each other i.e., $\frac{1-w_i}{4}$. This allows the designer to understand the impacts of prioritizing different objectives on the three decision variables and also how these objectives might be correlated with one another. 

% Tradeoffs - increased cycling of battery or DG implies more frequent replacement needed which raises costs. Increasing the number of online operational hours of the DG may also lead to higher startup and shutdown costs.

\subsection{Varying weight $w_1$ on normalized LCOE}

Increasing $w_1$ prioritizes the minimization of levelized electricity costs in the optimization. From \cref{fig:varyw11} there seems to be a minimum threshold to which MG costs can be reduced (relative to the baseline system running solely on the DE). The LCOE only falls initially as $w_1$ rises above zero but then remains more or less stagnant. In other words, the cheapest MG configuration possible costs slightly less than half of the baseline alternative, and further cost-minimization beyond this point may worsen other objectives.

\begin{figure}[htbp]
  \centering
  \begin{subfigure}[b]{0.33\linewidth}
    \includegraphics[width=\linewidth]{Sizing_vary_w1_1.png}
    \caption{\label{fig:varyw11}}
  \end{subfigure}
  \begin{subfigure}[b]{0.33\linewidth}
    \includegraphics[width=\linewidth]{Sizing_vary_w1_2.png}
    \caption{\label{fig:varyw12}}
  \end{subfigure}
  \begin{subfigure}[b]{0.32\linewidth}
    \includegraphics[width=\linewidth]{Sizing_vary_w1_3.png}
    \label{fig:varyw13}
  \end{subfigure}
  \caption{Effects of varying the weight placed on the LCOE term in the multiobjective optimization.}
  \label{fig:varyw1}
\end{figure}

Thus, prioritizing LCOE minimization actually leads to a linear increase in the overall weighted cost function. As $w_1$ increases, emissions gradually fall signaling increased use of RES sources over fossil-fueled backup generators. This also causes the dumped energy ratio to rise, indicating that pursuing cost-effectiveness may lead to solutions that are not as efficient. To cut down on costs, the BS capacity installed is reduced and to account for this, the remaining battery storage needs to be cycled more often as \cref{fig:varyw12} shows. Similarly in \cref{fig:varyw13}, there is a shift in RES capacity away from solar PV towards WT power, which is relatively cheaper on a levelized cost basis, primarily due to higher capacity factors. 

\subsection{Varying weight $w_2$ on emissions}

\begin{figure}[htbp]
  \centering
  \begin{subfigure}[b]{0.33\linewidth}
    \includegraphics[width=\linewidth]{Sizing_vary_w2_1.png}
    \caption{\label{fig:varyw21}}
  \end{subfigure}
  \begin{subfigure}[b]{0.33\linewidth}
    \includegraphics[width=\linewidth]{Sizing_vary_w2_2.png}
    \caption{\label{fig:varyw22}}
  \end{subfigure}
  \begin{subfigure}[b]{0.32\linewidth}
    \includegraphics[width=\linewidth]{Sizing_vary_w2_3.png}
    \caption{\label{fig:varyw23}}
  \end{subfigure}
  \caption{Effects of varying the weight placed on the emissions term in the multiobjective optimization.}
  \label{fig:varyw2}
\end{figure}

Unlike the case with LCOE, prioritizing emissions reduction results in co-benefits by increasing the renewable energy penetration to nearly 100\%. This brings both annual emissions and weighted cost to almost zero in \cref{fig:varyw21}, with an increase of both PV and WT capacity in \cref{fig:varyw23}. However, this causes LCOE and dump energy to rise especially as $w_2$ gets close to 1, implying that the last few $kg$ of emissions are the most expensive to eradicate and usually entail high RES curtailment. Higher reliance on RES reduces DE online time and requires larger BS to compensate, as in \cref{fig:varyw22}.

\subsection{Varying weight $w_3$ on DPSP} 

Since the DPSP is always zero for the default DG power rating of 16 kW, the sensitivity analysis on $w_3$ was performed on a system with a smaller DG rated at 7 kW. 

\begin{figure}[htbp]
  \centering
  \begin{subfigure}[b]{0.33\linewidth}
    \includegraphics[width=\linewidth]{Sizing_vary_w3_1.png}
    \caption{\label{fig:varyw31}}
  \end{subfigure}
  \begin{subfigure}[b]{0.33\linewidth}
    \includegraphics[width=\linewidth]{Sizing_vary_w3_2.png}
    \caption{\label{fig:varyw32}}
  \end{subfigure}
  \begin{subfigure}[b]{0.32\linewidth}
    \includegraphics[width=\linewidth]{Sizing_vary_w3_3.png}
    \caption{\label{fig:varyw33}}
  \end{subfigure}
  \caption{Effects of varying the weight placed on the DPSP term in the multiobjective optimization.}
  \label{fig:varyw3} 
\end{figure} 
Placing greater emphasis on minimizing load deficiency also reduces emissions in \cref{fig:varyw31} and increases REF, but leads to an increase in LCOE and excess energy ratio in \cref{fig:varyw32}. Surprisingly, it also causes a decline in DE operational hours in \cref{fig:varyw33} which suggests that oversizing RES coupled with larger BS may actually be the more optimal design choice to meet a greater portion of load, rather than depending on the backup DE more heavily during peak load periods.

\subsection{Varying weight $w_4$ on dump energy ratio}
 
\begin{figure}[htbp]
  \centering
  \begin{subfigure}[b]{0.33\linewidth}
    \includegraphics[width=\linewidth]{Sizing_vary_w4_1.png}
    \caption{\label{fig:varyw41}}
  \end{subfigure}
  \begin{subfigure}[b]{0.33\linewidth}
    \includegraphics[width=\linewidth]{Sizing_vary_w4_2.png}
    \caption{\label{fig:varyw42}}
  \end{subfigure}
  \begin{subfigure}[b]{0.32\linewidth}
    \includegraphics[width=\linewidth]{Sizing_vary_w4_3.png}
    \caption{\label{fig:varyw43}}
  \end{subfigure}
  \caption{\label{fig:varyw4} Effects of varying the weight placed on the dump or excess energy ratio in the multiobjective optimization.}
\end{figure} 

Excess energy production in the MG occurs mainly due to curtailment of wind and solar output. In order to mitigate this wasted energy that is earthed or sent to dump loads, RES capacity needs to be reduced as seen in \cref{fig:varyw43}, which also leads to contraction of BS size in \cref{fig:varyw42}. Thus, the DE needs to be used to meet a greater extent of the load and is online for longer in \cref{fig:varyw42}. This increased fossil-fuel dependence raises emissions and LCOE making the system less optimal, which is why the weighted cost initially increases slightly with $w_4$ in \cref{fig:varyw41}, before decreasing to zero as $w_4$ approaches 1.

\subsection{Varying weight $w_5$ on (1-REF)}
 
\begin{figure}[htbp]
  \centering
  \begin{subfigure}[b]{0.33\linewidth}
    \includegraphics[width=\linewidth]{Sizing_vary_w5_1.png}
    \caption{\label{fig:varyw51}}
  \end{subfigure}
  \begin{subfigure}[b]{0.33\linewidth}
    \includegraphics[width=\linewidth]{Sizing_vary_w5_2.png}
    \caption{\label{fig:varyw52}}
  \end{subfigure}
  \caption{\label{fig:varyw5} Effects of varying the weight placed on the (1-REF) term in the multiobjective optimization.}
\end{figure}  

The results in \cref{fig:varyw5} which focus on maximizing renewables penetration show similar trends and patterns to those for emissions reduction in \cref{fig:varyw2}. This follows from the strong correlation between these two objectives found earlier in \cref{fig:correl}.

\section{\label{app:improv} Model caveats, potential refinements and extensions}

\subsection{Alternative dispatch strategies}

The current strategy assumed for sizing optimization results in DGs being used during peak demand conditions. However, it may be more optimal to instead use DGs under off-peak situations and use BS or RES during higher demand. This would entail a different priority order for unit commitment and dispatch since all the RES would not just be used to meet as much of the load as possible at the instant when solar or wind power is available. Some portion of this would also be directed towards charging the BS during off-peak times. Similarly, there may be other innovative methods that could improve objectives even more. Thus, it is necessary to explore the full space of possible strategies more comprehensively even at the design and planning stage.

\subsection{Improve computational efficiency}

Non-convex formulations generally entail greater computational complexity and are more challenging to solve, requiring further processing time. Future work could look into relaxations, variable transformations and/or linearizations that would make the problem easier to solve. For e.g., trying to remove the currently used discrete decision variables and making all variables continuous could be one way to convexify the problem and increase the speed of optimization. This could potentially also allow the incorporation of more complex, realistic, non-linear models describing MG components, which improve accuracy but would not be tractable with problems that are not convex. Another potential method to increase speed would be the vectorization of the objective functions.

\subsection{More comprehensive multiyear modeling}

The multi-period optimization could also be refined to simulate the system more accurately over its entire lifetime. Currently, the sizing optimization only simulates the MG for one year and then assumes that all remaining years in its lifecycle are similar to the first, accounting for discounting and interest. While this is a reasonable approach, it has some deficiencies due to which it does not reflect the actual operation of the MG perfectly. For instance, this does not account for the average annual increase in electricity load of 10-15\% expected in Kenya and Tanzania due to population growth, economic development, rising incomes, and living standards \cite{usea}. Similarly, variations are possible in several other model parameters as well (e.g., volatile fuel prices) and the current model does not fully account for this, even though a sensitivity analysis was performed. Finally, the dispatch optimization was only performed for a time horizon of one day due to computational constraints. Optimizing the dispatch over longer periods could help explore the system's performance and response to disturbances in the longer term. 

\subsection{Simulation using smaller time steps}

Discretizing into finer time steps in the order of minutes or even seconds (e.g., $\Delta t = 5\; min$) instead of the currently used hourly time step ($\Delta t = 1\; h$) is likely to produce more realistic results since it simulates the time-marching battery dynamics and evolution of its SOC more accurately. The modeling of the DG dispatch would also be enhanced since ramp rate constraints can now be enforced, which limits how fast the unit can either be brought online (startup) or taken offline (shutdown). Furthermore, it offers a more detailed resolution for all variables across the board. Currently, there are discontinuities in the dispatch and $SOC$ plots produced by the model which are not physically feasible and are likely caused by the time step being too large. Reducing $\Delta t$ will help obtain more continuous, smooth dispatch curves that are closer to those expected during actual operation. However, it is challenging to achieve this in reality due to the additional complexity, extra computational burden, and longer runtimes that result. It is also extremely difficult to obtain reliable and accurate sub-hourly climate and load data, particularly for locations and settlements in Africa. 

\subsection{\label{sec:dr} Demand response (DR) and Demand-side management (DSM)}

The system's robustness to the stochasticity of demand and load fluctuations can be improved by implementing a demand-side management program in the MG. Several possible approaches have been suggested in the past to mathematically model DR programs and interpret their results using concepts like game theory, mechanism design, and auction theory. One such intuitive model is described in \cite{nwulu2017} where the net expected benefit to the MG operator is maximized:  
\begin{equation}
    \max_{x,y} \; V_2 (\theta, \lambda)
\end{equation}
where $V_2$ is the benefit function of the operator:
\begin{equation}
    V_2 (\theta, \lambda) = \lambda x - y
\end{equation}
where $x$ is the amount of power consumption that is shifted or curbed by a customer of type $ 0 \leq \theta \leq 1$. The customer type indicates their readiness to shift or curb their power demand, with $\theta = 1$ for the most willing and for the least willing. $\lambda$ is the 'value of power interruptibility, a measure of how beneficial it is for the grid operator to modify a specific load, and is determined from optimal power flow (OPF). $y$ is the amount of monetary compensation given to a customer for a particular amount of load shifted or curbed. The customer's benefit function is then given by:
\begin{equation}
    V_1(\theta,x,y) = y - c(\theta,x)
\end{equation}
where $c(\theta,x)$ is the cost incurred by a customer of type $\theta$ for shifting or reducing their load by $x$ kW, which is increasing in $x$ and decreasing in $\theta$. Finally, there are constraints on the operator's total budget (for compensating customers) and daily limits on each customer's interruptible power.

Such agent-based formulations can be useful for simulating behavior in small networks like microgrids with relatively few market participants. However, their computational complexity scales poorly with the number of agents, resulting in larger networks, more nodes, and increasingly complex interlinks between these. Another technical barrier is that such programs would require two-way communication channels for all agents involving bidirectional flows. This would entail extensive upgrades to existing network infrastructure. Furthermore, smart meters and home energy monitors integrated with the Internet of Things (IoT) also bring up concerns regarding the privacy of individual users. However, there are techniques available that can help preserve privacy while still providing enough useful information and resolution to perform load management, such as non-intrusive load monitoring (NILM) and disaggregation. 

The results of such models are also non-deterministic since they involve decision-making by consumers, producers, and prosumers, who can often be irrational as well. There is a need to determine response functions for each agent involved and approximate their individual preferences and private costs. For instance, in the example model given above, it is challenging to actually determine the types $\theta_i$ and cost functions $c_i(\theta_i,x)$ of all the consumers and prosumers in the MG. Finally, these results can often only be presented in a probabilistic sense, which presents another barrier to practical implementation.

\subsection{Improved load and RES forecasting}

Another approach to mitigate the MG's sensitivity and effects of uncertainty is to build forecasting models that generate more accurate predictions for the load as well as solar and wind output. These could also be combined with predictive control and online optimization methods to dynamically update forecasts at each time step, rather than using static forecasts or relying on historical weather (temperature, wind speed, irradiance) and consumption data. It would also be worthwhile to test both the sizing and dispatch optimization with community load profiles calculated using methods and sources other than NREL ReOpt. For instance, past studies have determined empirical relationships to approximate the power demand of typical isolated, off-grid rural African villages as a function of time, such as \cite{load_profile}:
\begin{equation}
    P \; [kW] = e^{[sin(0.3409 - sin(0.68039t) - 0.16801t)]}
\end{equation}

% \subsection{Model predictive control (MPC) formulation}

\subsection{Improved control techniques \label{app:improvements}}

The current model assumes perfect foresight and no errors in forecasts. The scheduler knows all the actual values of both demand and RES output over the full prediction or planning horizon, by using precalculated annual (or daily) load profiles and known historical climate data. This method is useful for initially sizing the system and devising nominal dispatch strategies, as well as an idealized benchmark to evaluate other strategies against \cite{parisio2014}. However, it is often very challenging to achieve this in reality due to the lack of availability of such data with high reliability. As described in more detail in \cref{sec:control_survey}, the MG's performance can be improved by applying control strategies like model predictive or receding horizon control. Since this problem involves a non-convex, non-smooth, 'black box' objective function (i.e., having no closed-form solution or gradient), it may prove to be more challenging to reformulate it into standard MPC form. Nevertheless, commercial solvers like GUROBI have been shown to efficiently solve such mixed-integer linear and quadratic programs without requiring complicated heuristics or decomposition methods \cite{parisio2014}. 

\section{Detailed survey of optimization algorithms \label{app:opt_alg}}

In addition to using available tools like HOMER Pro and DER-CAM, many studies have built their own optimization models by formulating sizing and dispatch problems usually as non-linear mixed integer minimization programs. Relaxations and variable transformations could also be used to convert these to mixed integer linear programs while certain assumptions and functions could be modified to potentially also convexify the problem. Linearized formulations are more computationally tractable and efficient than quadratic programs. Thus, quadratic functions like the fuel costs of DGs are often approximated to linear expressions like the maximum over affine functions as done in \cite{parisio2014}. Some of the most common global optimization techniques used to solve such problems are described in more detail in the following subsections. These artificial intelligent search techniques usually prove to be more effective than classical methods like linear programming when the dispatch problem becomes more complex \cite{Qu2018}. In addition, other methods like Lyapunov optimization \cite{pathak2017} and discrete Fourier transforms \cite{dft} have also been used in the literature.

\subsection{Genetic (GA) and evolutionary algorithms (EA)}

These are a class of optimization algorithms capable of solving smooth or non-smooth optimization problems with any type of constraints including integer constraints. It is a stochastic, population-based method which randomly searches the space of candidate solutions by \textit{selection}, \textit{mutation} and \textit{crossover} among population members. It imitates the biological process of evolution by natural selection and represents individual members as potential solutions to a complex optimization problem \cite{nemati2015}. The solution population of chromosomes is first initialized as a vector or array of bits or character strings. The algorithm seeks to minimizes the fitness value since it measures deviations from predefined objectives and constraints. It can be formulated as a weighted sum of the cost (objective function value), penalties on constraint violations and the parity degree (a measure of similarity among individuals or possible solutions) as per \cref{eq:GA_app}, where $W_{Pen}$ and $W_{Par}$ are variable weighting factors while $R_{Cost/Pen}$ and $R_{Cost/Par}$ are ratios of mean values used to normalize the three criteria to be comparable. The parity term shows diversity and improves the survival probability of individuals that differ from the group, thus decreasing the likelihood of quickly converging to, and getting stuck at local optima \cite{nemati2015}.

\begin{equation}
    F = Cost + W_{Pen} \cdot R_{Cost/Pen} \cdot Penalty + W_{Par} \cdot R_{Cost/Par} \cdot Parity
    \label{eq:GA_app}
\end{equation}

In accordance with Charles Darwin's theory of `survival of the fittest', the fitness of each chromosome is assessed at each iteration and the best or minimizing solution is selected \cite{pathak2017}. These solutions then repeatedly undergo processes of slight modification and recombination with one another to retain properties that better adapt them to their environment than competitors, for several generations until a preset stop criterion (generation number) is met \cite{nemati2015}. Similar to genetic algorithms, evolutionary programming also relies on Darwinian principles of evolution and adaptation. Whereas GA attempts to simulate micro-level processes like chromosomes, evolutionary algorithms seek to mimic macro-processes such as phenotype variations and inheritance \cite{pathak2017}. Since these methods use a population of solutions in their search, multiple well-distributed, non-dominated solutions can be found in a single run while requiring little information about the domain or search space for a given problem \cite{Qu2018}. \textit{Differential evolution} is a variant of evolutionary algorithms that differs in the solution selection step due to their use of non-domination sorting and ranking selection procedures \cite{Qu2018}. These have proven to be effective in solving environmental/economic dispatch (EED) as well as unit commitment (UC) problems \cite{lu2011}.

One approach for multiobjective optimization is to combine the different individual objective functions into a single composite using weights that sum up to unity. Another more general method is to identify a whole solution set that is Pareto optimal. This is the set of all non-dominated feasible solutions that cannot be improved further relative to any objective without worsening performance on at least one other objective. Genetic and evolutionary methods prove to be useful to generate such sets and help designers clearly visualize trade-offs amongst pursuing various objectives like cost-effectiveness, efficiency, and reliability \cite{fadaee2012}.

The main disadvantage with these methods is that complexity rises significantly with the number of problem parameters which also increases the programming and solving difficulty \cite{fathima2015}. Exploration of new solution spaces is limited with traditional GA and adaptability is low due to the crossover and mutation operations being preset according to experience. In order to resolve such issues, some studies have developed modified versions such as improved adaptive genetic algorithms which avoid phenomena like premature, local convergence and arrive at global optima in a faster, more stable manner \cite{wu2016}.

\subsection{Particle swarm optimization (PSO) \label{app:pso}}

Particle swarm optimization is another biologically inspired algorithm that attempts to mimic dynamics of group behavior and swarm intelligence in order to converge to the optimal solution. It is an iterative, agent-based, distributed problem-solving method that simulates the social interactions of a swarm of individuals moving together in search of food inside a specific region \cite{fathima2015}, as seen in animal societies and insect colonies \cite{pathak2017}. It employs a population of particles representing potential solutions to thoroughly explore the hyperspace of feasible solution \cite{alrashidi2010}. In comparison to others, the PSO algorithm is often easier to implement owing to its relative simplicity, fewer parameter settings and strong optimization ability. However, it may get stuck at local optima and has trouble satisfying equality constraints in the initialization and updating processes of the particle swarm \cite{wu2014}. Additionally, GA is more efficient in cases where there are more than three components to the solution. Improved and enhanced multiobjective PSO algorithms have been successfully used to accurately size battery systems for microgrids and solve a variety of EED problems efficiently. This includes highly complex and even non-convex large-scale examples such as optimal power flow (OPF) \cite{alrashidi2010}, resulting in quick convergence to high quality solutions with relatively low computational cost \cite{pathak2017}.

The conventional PSO algorithm randomly initializes a particle community with positions $[x_{i1}, x_{i2}, ... \; , x_{id}]$ and speeds $[v_{i1}, v_{i2}, ... \;, v_{id}]$ where $i = 1,2,... \; n$ and $n$ is the population size while $d$ is the dimension of the search space describing each parameter or coordinate of the solution. Each particle's position is determined by two best values, (1) the best fitness value (minimal objective function) that the individual particle itself has achieved so far ($pbest$) and (2) the global best value obtained by the optimizer for the entire population so far ($gbest$). In addition to positions tracking the values of input variables, each particle also has a velocity moving it towards the individual and global bests in terms of fitness \cite{borhanazad2014}, to find the best among all feasible solutions within the predefined space (set by the initial population). The procedure can be summarized as:

\begin{enumerate}
    \item Evaluate fitness of each particle or potential solution
    \item Update both $pbest$ and $gbest$
    \item Update velocity and position of each particle in the population using \cite{wu2014}:
    \begin{align}
    \label{eq:pso}
        x_{i,d}^{k+1} & = x_{i,d}^{k} + v_{i,d}^{k} \\
        v_{i,d}^{k} & = \omega \cdot v_{i,d}^{k} + c_1 \cdot rand_1^k \cdot (pbest_{i,d}^k - x_{i,d}^{k}) + c_2 \cdot rand_2^k \cdot (gbest_{i,d}^k - x_{i,d}^k)
    \end{align}
\end{enumerate}

The velocity is calculated as a random weighted integer as in \cref{eq:pso} where $v_{i,d}^k$ is the speed and $x_{i,d}^k$ is the position along dimension $d$ of particle $i$ in iteration $k$, $\omega$ is the inertia weight factor, $c_1$ and $c_2$ are learning factors, $pbest_{i,d}^k$ and $gbest_{i,d}^k$ are the local and global best values of coordinate $d$ in iteration $k$, and $rand_1^k$ and $rand_2^k$ are random numbers uniformly distributed in $(0,1)$ \cite{wu2014}. The algorithm keeps updating each particle's location and speed until some predefined termination condition such as the maximum number of iterations or a target fitness value is met, at which point the optimal group value $gbest$ provides the desired solution. Similar to PSO, there are several other swarm-intelligence based approaches such as ant or bee colony optimization, bacterial foraging algorithms and group search optimizers \cite{pathak2017}. These are just a small subset of the larger, more general group of bio-inspired algorithms to solve complex, multiobjective optimization problems.

\subsection{Surrogate optimization}

This algorithm is useful while working with expensive, time-consuming objective functions that entail a great deal of computational complexity and effort. It uses a surrogate function that approximates the actual objective function but takes less time to evaluate. The algorithm chooses quasirandom points within the bounds on input variables and then constructs the surrogate by interpolating upon the true objective using a radial basis function interpolator \cite{gutmann2001}. Thus, the best point minimizing the surrogate can be taken as a reasonable approximation for the minimizer of the objective. The method tries to balance two competing goals: (1) Complete exploration of search space for a global minimum and (2) Increase speed by arriving at a good solution using as few objective function evaluations as possible. This works best when the objective is smooth and continuous, and has been proven to arrive at globally optimal solutions on bounded domains even though the convergence is not fast \cite{gutmann2001}. However, smoothness is not strictly required to use this technique. Such an approach can greatly speed up computational times for complex component sizing and economic dispatch problems with non-linear cost functions and a large number of variables and constraints. There is usually no stopping criterion that can accurately recognize when the solver is near a global optimum. Instead the program terminates upon reaching a prescribed number of iterations, function evaluations or time duration and outputs the `best' solution found within that preset computational budget.

\subsection{Simulated annealing (SA)}

This is a stochastic search technique modelled after the annealing process in metals that are molten at high temperatures and then allowed to cool and freeze into crystalline structures with minimum energy, with temperature control being the crucial element. A randomly chosen current state is perturbed to a new state at a fixed temperature $T$ with $\Delta E$ being the energy gap between the two states. If $\Delta E < 0$, the new state is accepted as the updated current state and if $\Delta E \geq 0$, the new state is accepted with probability:
\begin{equation}
\label{eq:sa}
    Pr\;(acceptance) = exp\left(\frac{-\Delta E}{k_BT}\right)
\end{equation}
where $k_B$ is the Boltzmann constant. These perturbations are repeated until the system reaches thermal equilibrium with its environment at temperature $T$, following which the sequence can be repeated by slightly decreasing the temperature set point. Such an approach can be extended to hybrid microgrid optimization problems where system energy $E$ is replaced by the objective cost function, temperature $T$ corresponds to the control parameter and \cref{eq:sa} describes the acceptance probability of a particular solution (i.e., system configuration or dispatch strategy) \cite{fathima2015}. However, SA is less commonly used for sizing and dispatch of hybrid energy systems compared to popular tools like GA, EA and PSO. Certain past studies have concluded that SA can be faster to converge, but are less efficient than other metaheuristic methods like tabu-search algorithms for optimal sizing problems aimed at minimizing LCOE \cite{sa_tabu}.

% ============================================================

\section{Detailed survey of microgrid control strategies \label{app:control}}

After completing (1) the initial design to choose and size the generation and storage technologies and (2) devising the optimal strategy to schedule and dispatch these components, the next major area of research focuses on the real-time, active control of such microgrid systems. Closed-loop systems are likely to be superior to open-loop model-based approaches in practical implementation, due to their ability to better predict future system behavior, and compute control signals and corrective actions needed to deal with both demand and supply variations \cite{tazvinga2014}. In particular, control is essential to accommodate the high degree of uncertainty in RES power generation due to their intermittency and unpredictable fluctuations. This becomes even more critical with increased penetration of localized, distributed, and decentralized power generation, especially in the case of islanded and off-grid systems, where there are special concerns regarding low power quality, safety, and overload issues \cite{pourbabak2017}. Furthermore, there is added complexity due to bidirectional flows of both power and information. Thus, some form of feedback control is necessary to improve the robustness of the system and minimize the impact of forecasting errors (in load and RES availability), uncertainties, and external disturbances. Examples of disturbances include extreme weather (like wind lulls) affecting wind or solar power output, natural disasters, sudden spikes in demand (particularly due to heating or cooling loads), etc. Conventional microgrid control structures can be hierarchically differentiated into three levels. Primary control is responsible for regulating voltage, frequency and power of individual converters. Secondary control regulates power quality and mitigates voltage or frequency deviations, imbalances and undesirable harmonics. In practice, this translates to maintaining grid frequency around a certain target value and stabilizing supply frequency by regulating real and reactive power flows in real-time, respectively. Finally, the tertiary control handles power exchanges with the external macrogrid and/or other adjacent microgrids \cite{dragicevic_2017}. It could also include higher-level tasks like system-wide techno-economic optimization during the design stage and efficiency improvements while making day-to-day dispatch and operational decisions \cite{mg_review}. There are broadly three types of control strategies possible in large scale, cyber-physical systems applications like microgrids: Centralized, decentralized and distributed control.

\subsection{Centralized control}

This is the most mature and well-researched of the three methods and has been commonly used to control and economically dispatch generators in electrical power systems \cite{su2015}. Data is collected from all over the system and processed together in a single controller in order to make decisions which will be directly sent as control signals to each individual agent through supervisory, control, and data acquisition (SCADA). SCADA is an advanced automation control system that centrally manages the control, gathering, and monitoring of a power system's operation \cite{scada}. Such schemes have a simple conceptual and theoretical framework and have been studied for decades, resulting in the development of a large number of controllers particularly for linear systems \cite{farina2012}. However, these are challenging to implement in reality due to the high level of connectivity and two-way communication links needed, especially as the number of agents increases. Real-time transmission of measurements and computation of control variables may not be feasible if the network is deployed over a large physical area. Thus, these systems may not be cost-effective or computationally tractable, and complexity scales poorly with network size. Model-based, offline synthesis may also be difficult to achieve for larger systems with many inputs, states, and outputs. Finally, such an approach may be unappealing for socioeconomic and political reasons as well, especially since all agents are required to share their private information with the central controller \cite{farina2012}. Nevertheless, centralized controllers are a good starting point for smaller microgrids with a limited number of agents (i.e., producers, consumers, and prosumers) and relatively simple network topologies. After designing an effective centralized scheme, the microgrid operator could explore ways to transition this to a distributed scheme to account for increasing size and complexity.

\subsection{Decentralized control}

This is a schematic where each agent, node, or subsystem (e.g. DG, consumer, or BS unit) has its own controller, allowing them to make decisions and apply control actions based only on local measurements like voltage or frequency \cite{pourbabak2017}. This avoids the need for global information over the entire system and allows individual agents to protect their privacy since they no longer exchange their information with a central unit or with other nodes \cite{gan2012}. This can also speed up processes by computing control variables in parallel and allows for easier real-time communication since sensors, controllers and actuators within a given subsystem will be located closer to one another \cite{farina2012}. However, the lack of communication between controllers at different nodes limits attainable performance and cannot guarantee \textit{global} optimality or stability \cite{pourbabak2017}. Nevertheless, both centralized and decentralized methods have been demonstrated to improve transient stability in small-scale power systems like microgrids \cite{senjyu2005}. Decentralized systems can also be more stable than centralized alternatives having similar configurations, if some nodes get disconnected due to faults or failure, other agents in the system can still keep functioning using local information.

\subsection{Distributed control}

This is a consensus-based approach where individual agents or nodes use both locally measured parameters at their node as well as information shared by their nearest neighbors. In contrast to decentralized control with no information sharing, adjacent agents here share their private information with each other through two-way communication channels. This enables global optimization while still allowing efficient, parallel computation. Furthermore, it allows the designer to tune the trade-off between system performance versus communication burden among controllers and modify network design accordingly \cite{farina2012}, while still allowing agents to retain some control over their privacy. Thus, distributed control is a compromise that offers some of the best properties of both centralized and decentralized control, making it an increasingly popular tool for microgrids. In particular, this approach is well-suited to handle the varying nature of the microgrid and is not greatly impacted by structural changes expected in smart grids in the near future, due to local information sharing. This enables us to satisfy both the \textit{peer-to-peer} and \textit{plug-and-play} requirements. The 1st condition ensures reliability, implying that there is no single component that is absolutely critical to microgrid operation and the system can function even in case of failure of individual generators or storage devices. The 2nd, plug-and-play characteristic means that an additional, new unit can be installed at any point in the electrical system without having to reconfigure the network or controller \cite{mg_prop}. This allows the microgrid (or any decentralized power grid, in general) to easily extend and grow as more producers and customers get connected \cite{pourbabak2017}, making it an appropriate solution for remote rural electrification. Distributed approaches also intuitively lend themselves well to real-time DR and demand-side management (DSM) algorithms that engage both users and generators \cite{deng2014}, as well as systems involving dynamic, flexible DGs and smart, responsive loads \cite{dist_1, dist_2}, owing to their scalable nature, fast response and model independence. However, they still face challenges in designing the right communication and control network topology, upgrading existing infrastructure to facilitate bidirectional flows, and achieving convergence through consensus among local agents quickly \cite{pourbabak2017}.

\subsection{Model predictive control \label{app:mpc}}

Model predictive control (MPC) is by far the most common method used in the literature for optimal microgrid operation. This makes the originally static optimization problem dynamic by repeatedly optimizing over a receding, finite time horizon and solving an open-loop optimal control problem at each sampling instant using the current state as the initial state. Only the first control action of the resulting sequence is applied to the plant, after which the process is repeated using the new state \cite{mpc}. The controller therefore uses the current state as a basis for estimating future state variables, by coupling inputs and outputs to be dependent on one another in the observer. Linear MPC generally solves either a quadratic (QP) or linear program (LP) at each iteration or control interval i.e., either a quadratic or linear cost function (objective) subject to linear constraints. Thus, MPC differs from traditional controllers like PI and PID since it solves the optimal control problem `online' for the \textit{current} state of the system instead of offline control law synthesis (i.e., a uniform precomputed feedback policy that applies to all states) \cite{mpc}. This gives it a unique, predictive ability valuable for MG control as it can anticipate future events like RES supply fluctuations or demand changes and take actions needed to damp their effects, while satisfying time-varying request and operation constraints \cite{parisio2014}. This allows it to achieve better performance over relatively long periods during which disturbances could occur.

Variants of predictive control implemented in the literature include formulations using mixed-integer linear as well as nonlinear programs since MG modeling requires both discrete (e.g. ON/OFF states of DGs) and continuous (e.g., battery SOC and charge/discharge rates) decision variables \cite{parisio2013, parisio2014}. Such feedback mechanisms have been successful in compensating for uncertainty in RES outputs and time-varying loads. Robust linear and nonlinear economic MPC (E-MPC) methods have been developed to optimize setpoints for microgrid dispatch that minimize operational cost \cite{zachar2016} and/or maximize profits from selling power back to the main grid \cite{pereira2017}. Finite horizon objectives in MPC comprise stage, terminal, and penalty costs, with E-MPC using economic costs as the stage cost function. Closed loop discrete linear MPC has been applied to a photovoltaic--wind--diesel--battery system similar to that considered here \cite{tazvinga2014}. In addition to load and RES uncertainty, predictive control has also been used to deal with uncertainty arising from variations in model parameters of generators \cite{prodan2014}. Another interesting example is the use of adaptive, switched MPC where switching constraints are used instead of switched state-space models to model shifting between modes like battery charging and discharging \cite{zhu2015}.

% ============================================================

\section{Detailed survey of microgrid design tools \label{app:tools}}

\subsection{HOMER Pro}

HOMER (Hybrid Optimization of Multiple Energy Resources), also known as Hybrid Optimization Model for Electric Renewables Pro, is a tool originally developed by the National Renewable Energy Laboratory (NREL) and now distributed by the commercial entity HOMER Energy. Based on literature surveyed, this is the most widely used commercial tool for microgrid design optimization in both academia and industry today. It is currently the global standard for optimizing microgrid design in all sectors, from village power and island utilities to grid-connected campuses and military bases. HOMER is a simulation model that simulates a viable system for \textit{all} possible combinations of technologies being considered. It uses a proprietary, derivative-free optimization algorithm that uses brute force permutations and exhaustive search \cite{weng2016} to minimize the total net present costs (TNPC) of the system (defined as present value of sum total of costs minus revenues). It requires six types of inputs - meteorological or climate data, load profiles, equipment characteristics, search space specifications, economic and technical data \cite{bahramara2016}. Optimal sizes of components are then determined in three stages of simulation, optimization, and sensitivity analysis, using time steps ranging from 1 minute to 1 hour which capture intra-day variability. The simulation and optimization occur simultaneously while the sensitivity analysis determines the effects of uncertainty in non-deterministic parameters (such as fuel cost, wind speed, solar radiation, and component costs) on the results in terms of the best (lowest NPC), feasible system plans.

\subsection{DER-CAM}

The Distributed Energy Resources Customer Adoption Model (DER-CAM) is a decision-support tool developed by Lawrence Berkeley National Laboratory (LBNL) that is capable of optimizing \textit{both} design \textit{and} dispatch for DER investments used in either buildings or hybrid, multi-energy microgrids. Compared to HOMER, it offers more flexibility in terms of modifying the objective function and allowing for multi-objective optimizations. Unlike HOMER which is a simulation model based on heuristics and non linear formulations, it uses a mixed integer linear program (MILP) to find global optima quickly \cite{dercam} with key inputs required being similar to HOMER. Rather than using a full time series, it models loads and solar irradiance as characteristic day types to capture seasonal variations \cite{weng2016}. Its outputs include optimal selection and capacity of DERs, their relative placement at nodes and dispatch schedules to meet economic, resiliency and reliability targets. In addition to monetary costs, it also considers carbon emissions and environmental criteria in its optimization - providing a breakdown of both costs and emissions associated with various end-use loads among its outputs. It assumes no deterioration in output or efficiency over each component's lifetime and constraints on start-up, shut-down and ramp rates are excluded by default \cite{dercam}.

\subsection{REopt}

REopt is an MILP-based techno-economic decision support model from NREL to optimize decentralized energy systems for a variety of applications including both grid-connected and off-grid microgrids, using hourly or sub-hourly simulation time steps. It recommends the optimal mix, sizes and operating strategies of renewables, conventional thermal generation and storage to maximize economic performance and cost savings while connected to the grid, and increase survivability and days of autonomy during outages in islanded mode \cite{tools_summary}. The full software is not publicly available and the accessible, rudimentary version is not very comprehensive since it offers only a small subset of free features. However, supporting NREL datasets obtained from REopt cost models were used here to estimate load profiles and explore variations in Levelized costs of electricity (LCOE) specific to rural Africa \cite{nrel_data}. From the sample results obtained by NREL comparing LCOE for multiple system configurations, it can be concluded that diesel-only systems are the cheapest option but may not meet other design objectives. Although the PV+battery system is more expensive, the inclusion of a diesel backup makes the combined system even cheaper than the diesel-only option. This indicates that hybrid designs incorporating RES and storage along with fossil-fuels are likely to be optimal and most cost-effective for isolated microgrids.

\subsection{MDT}

The Microgrid Design Toolkit (MDT) is a software developed by Sandia National Laboratories that uses both MILP and genetic algorithms. Rather than using time steps, it is a discrete event simulation (examples of events could include changes in load, fuel outage, equipment failures etc.) \cite{tools_summary}. In addition to its sizing capability, MDT also helps to identify and characterize the trade-space of alternative microgrid design decisions in the early, preliminary stages. It optimizes both technology and performance management according to user-defined criteria to output a set of efficient trade-off design options as a Pareto optimal frontier.

\bibliographystyle{unsrt} % Use for unsorted references  
% \bibliographystyle{plainnat} % use this to have URLs listed in References
% \cleardoublepage

\begin{thebibliography}{9}
% \addcontentsline{toc}{section}{title=References}

\bibitem{SSA_decent}
Dagnachew, A. G., Lucas, P. L., Hof, A. F., Gernaat, D. E., de Boer, H. S., \& van Vuuren, D. P. (2017). The role of decentralized systems in providing universal electricity access in Sub-Saharan Africa–A model-based approach. \textit{Energy}, 139, 184-195.

\bibitem{WEO_2018}
World Energy Outlook (2018). \textit{International Energy Agency (IEA)}
\url{https://www.iea.org/weo2018/}. Accessed 17 June, 2019.  

\bibitem{kammen}
Casillas, C. E., \& Kammen, D. M. (2010). The energy-poverty-climate nexus. \textit{Science}, 330(6008), 1181-1182.

\bibitem{wri}
Odarno, L., Sawe, E., Swai, M., Katyega, M. J., \& Lee, A. (2017). Accelerating Mini-Grid Deployment in Sub-Saharan Africa. \textit{World Resources Institute and TaTEDO}. 

\bibitem{mentis}
Mentis, D., Howells, M., Rogner, H., Korkovelos, A., Arderne, C., Zepeda, E., Siyal, S., Taliotis, C., Bazilian, M., De Roo, A., Tanvez, Y., Oudalov, A., \& Scholtz, E. (2017). Lighting the World: the first application of an open source, spatial electrification tool (OnSSET) on Sub-Saharan Africa. \textit{Environmental Research Letters}, 12(8), 085003. doi:10.1088/1748-9326/aa7b29

\bibitem{quak_CBA}
Quak, E. J. (2018). The costs and benefits of lighting and electricity services for off-grid populations in sub-Sahara Africa. \textit{Institute of Development Studies (IDS)}.

\bibitem{grid_gtm}
Attia, B., Shirley, R. (2018). Distributed Models for Grid Extension Could Save African Utilities Billions of Dollars. \textit{Greentech Media}. \url{https://www.greentechmedia.com/articles/read/grid-extension-done-right-for-sub-saharan-africas-utilities#gs.q0ytps}. Accessed 16 July, 2019. 

\bibitem{macrogrid_prob}
Amin, M. (2008, July). Challenges in reliability, security, efficiency, and resilience of energy infrastructure: Toward smart self-healing electric power grid. In \textit{2008 IEEE Power and energy society general meeting-conversion and delivery of electrical energy in the 21st century (pp. 1-5)}. IEEE.

\bibitem{tech_challenges}
Salam, A. A., Mohamed, A., \& Hannan, M. A. (2008). Technical challenges on microgrids. \textit{ARPN Journal of engineering and applied sciences}, 3(6), 64-69.

\bibitem{mg_review}
Hirsch, A., Parag, Y., \& Guerrero, J. (2018). Microgrids: A review of technologies, key drivers, and outstanding issues. \textit{Renewable and Sustainable Energy Reviews}, 90, 402-411.

\bibitem{mg_mod_opt}
Porsinger, T., Janik, P., Leonowicz, Z., \& Gono, R. (2017). Modeling and optimization in microgrids. \textit{Energies}, 10(4), 523.

\bibitem{kaabeche2014}
Kaabeche, A., \& Ibtiouen, R. (2014). Techno-economic optimization of hybrid photovoltaic/wind/diesel/battery generation in a stand-alone power system. \textit{Solar Energy}, 103, 171-182.

\bibitem{tazvinga2013}
Tazvinga, H., Xia, X., \& Zhang, J. (2013). Minimum cost solution of photovoltaic–diesel–battery hybrid power systems for remote consumers. \textit{Solar Energy}, 96, 292-299.

\bibitem{kaabeche2011a}
Kaabeche, A., Belhamel, M., \& Ibtiouen, R. (2011). Techno-economic valuation and optimization of integrated photovoltaic/wind energy conversion system. \textit{Solar energy}, 85(10), 2407-2420.

\bibitem{bhandari2014}
Bhandari, B., Poudel, S. R., Lee, K. T., \& Ahn, S. H. (2014). Mathematical modeling of hybrid renewable energy system: A review on small hydro-solar-wind power generation. \textit{International journal of precision engineering and manufacturing-green technology}, 1(2), 157-173.

\bibitem{gu2014}
Gu, W., Wu, Z., Bo, R., Liu, W., Zhou, G., Chen, W., \& Wu, Z. (2014). Modeling, planning and optimal energy management of combined cooling, heating and power microgrid: A review. \textit{International Journal of Electrical Power \& Energy Systems}, 54, 26-37.

\bibitem{logenthiran2010}
Logenthiran, T., Srinivasan, D., Khambadkone, A. M., \& Raj, T. S. (2010, June). Optimal sizing of an islanded microgrid using evolutionary strategy. In \textit{2010 IEEE 11th International Conference on Probabilistic Methods Applied to Power Systems} (pp. 12-17). IEEE.

\bibitem{bahramara2016}
Bahramara, S., Moghaddam, M. P., \& Haghifam, M. R. (2016). Optimal planning of hybrid renewable energy systems using HOMER: A review. \textit{Renewable and Sustainable Energy Reviews}, 62, 609-620.

\bibitem{weng2016}
Weng, D. (2016). \textit{Microgrid Techno-Economic Assessment Presentation}. Electric Power Research Institute, Inc.

\bibitem{dercam}
Lawrence Berkeley National Lab. \textit{Distributed Energy Resources - Customer Adoption Model (DER-CAM)}. \url{https://building-microgrid.lbl.gov/projects/der-cam}. Accessed July 18, 2019. 

\bibitem{tools_summary}
Jimenez, Antonio, Haase, Scott G., \& Mathur, Shivani. \textit{Microgrid Analysis Tools Summary}. United States. doi:10.2172/1425364.

\bibitem{nrel_data}
Li, Xiangkun et al. (2018): \textit{Microgrid Load and LCOE modeling Results}. National Renewable Energy Laboratory. \url{https://data.nrel.gov/submissions/79}. Accessed June 12, 2019. 

\bibitem{saravanan2013}
Saravanan, B., Das, S., Sikri, S., \& Kothari, D. P. (2013). A solution to the unit commitment problem—a review. \textit{Frontiers in Energy}, 7(2), 223-236.

\bibitem{pathak2017}
Pathak, A., Pandit, M., \& Kumar, Y. (2017). Combined Economic Emission Dispatch in a Micro Grid: A Comprehensive Survey. \textit{International Journal of Applied Engineering Research}, 12(24), 15876-15901.

\bibitem{zachar2016}
Zachar, M., Daoutidis, P. (2016). Economic dispatch for microgrids with constrained external power exchange. \textit{IFAC-PapersOnLine}, 49(7), 833-838.

\bibitem{parisio2014}
Parisio, A., Rikos, E., \& Glielmo, L. (2014). A model predictive control approach to microgrid operation optimization. \textit{IEEE Transactions on Control Systems Technology}, 22(5), 1813-1827.

\bibitem{parisio2013}
Parisio, A., Rikos, E., Tzamalis, G., \& Glielmo, L. (2014). Use of model predictive control for experimental microgrid optimization. \textit{Applied Energy}, 115, 37-46.

\bibitem{nemati2015}
Nemati, M., Bennimar, K., Tenbohlen, S., Tao, L., Mueller, H., \& Braun, M. (2015, June). Optimization of microgrids short term operation based on an enhanced genetic algorithm. In \textit{2015 IEEE Eindhoven PowerTech} (pp. 1-6). IEEE.

\bibitem{nwulu2017}
Nwulu, N. I., \& Xia, X. (2017). Optimal dispatch for a microgrid incorporating renewables and demand response. \textit{Renewable Energy}, 101, 16-28.

\bibitem{paola2017}
De Paola, A., Angeli, D., \& Strbac, G. (2017). Price-based schemes for distributed coordination of flexible demand in the electricity market. \textit{IEEE Transactions on Smart Grid}, 8(6), 3104-3116.

\bibitem{fadaee2012}
Fadaee, M., \& Radzi, M. A. M. (2012). Multi-objective optimization of a stand-alone hybrid renewable energy system by using evolutionary algorithms: A review. \textit{Renewable and sustainable energy reviews}, 16(5), 3364-3369.

\bibitem{fathima2015}
Fathima, A. H., \& Palanisamy, K. (2015). Optimization in microgrids with hybrid energy systems–A review. \textit{Renewable and Sustainable Energy Reviews}, 45, 431-446.

\bibitem{Qu2018}
Qu, B. Y., Zhu, Y. S., Jiao, Y. C., Wu, M. Y., Suganthan, P. N., \& Liang, J. J. (2018). A survey on multi-objective evolutionary algorithms for the solution of the environmental/economic dispatch problems. \textit{Swarm and Evolutionary Computation}, 38, 1-11.

\bibitem{dft}
Xiao, J., Bai, L., Li, F., Liang, H., \& Wang, C. (2014). Sizing of energy storage and diesel generators in an isolated microgrid using discrete Fourier transform (DFT). \textit{IEEE Transactions on Sustainable Energy}, 5(3), 907-916.

\bibitem{wu2016}
Wu, H., Zhuang, H., Zhang, W., \& Ding, M. (2016). Optimal allocation of microgrid considering economic dispatch based on hybrid weighted bilevel planning method and algorithm improvement. \textit{International Journal of Electrical Power \& Energy Systems}, 75, 28-37.

\bibitem{lu2011}
Lu, Y., Zhou, J., Qin, H., Wang, Y., \& Zhang, Y. (2011). Environmental/economic dispatch problem of power system by using an enhanced multi-objective differential evolution algorithm. \textit{Energy Conversion and Management}, 52(2), 1175-1183.

\bibitem{wu2014}
Wu, H., Liu, X., \& Ding, M. (2014). Dynamic economic dispatch of a microgrid: Mathematical models and solution algorithm. \textit{International Journal of Electrical Power \& Energy Systems}, 63, 336-346.

\bibitem{borhanazad2014}
Borhanazad, H., Mekhilef, S., Ganapathy, V. G., Modiri-Delshad, M., \& Mirtaheri, A. (2014). Optimization of micro-grid system using MOPSO. \textit{Renewable Energy}, 71, 295-306. 

\bibitem{moshi2016}
Moshi, G. G., Bovo, C., Berizzi, A., \& Taccari, L. (2016, June). Optimization of integrated design and operation of microgrids under uncertainty. In \textit{2016 Power Systems Computation Conference (PSCC)} (pp. 1-7). IEEE.

\bibitem{alrashidi2010}
AlRashidi, M. R., AlHajri, M. F., \& El-Hawary, M. E. (2010). Enhanced particle swarm optimization approach for solving the non-convex optimal power flow. \textit{World Academy of Science, Engineering and Technology}, 62(2), 651-655.

\bibitem{gutmann2001}
Gutmann, H. M. (2001). A radial basis function method for global optimization. \textit{Journal of global optimization}, 19(3), 201-227.

\bibitem{tazvinga2014}
Tazvinga, H., Zhu, B., \& Xia, X. (2014). Energy dispatch strategy for a photovoltaic–wind–diesel–battery hybrid power system. \textit{Solar Energy}, 108, 412-420.

\bibitem{pourbabak2017}
Pourbabak, H., Chen, T., Zhang, B., \& Su, W. (2017). Control and energy management system in microgrids. In S. Obara \& J. Morel (Eds.), \textit{Clean Energy Microgrids} (Chapter 3, pp. 109--133). Institution of Engineering and Technology (IET).

\bibitem{sa_tabu}
Katsigiannis, Y. A., Georgilakis, P. S., \& Karapidakis, E. S. (2012). Hybrid simulated annealing–tabu search method for optimal sizing of autonomous power systems with renewables. \textit{IEEE Transactions on Sustainable Energy}, 3(3), 330-338.

\bibitem{su2015}
Su, W., Wang, J., \& Ton, D. (2015). Smart grid impact on operation and planning of electric energy systems. \textit{Handbook of Clean Energy Systems}, 1-13.

\bibitem{dragicevic_2017}
Dragicevic, T., Wu, D., Shafiee, Q., \& Meng, L. (2017). Distributed and decentralized control architectures for converter-interfaced microgrids. \textit{Chinese Journal of Electrical Engineering}, 3(2), 41-52.

\bibitem{scada}
Higgs, M. A. (2000). Electrical SCADA systems from the operators perspective. \textit{Power Engineering Journal}, 14(3), 133-137.

\bibitem{farina2012}
Farina, M., \& Trecate, G. F. (2012). Decentralized and distributed control. In \textit{EECI-HYCON2 Graduate School on Control 2012}.

\bibitem{gan2012}
Gan, L., Topcu, U., \& Low, S. H. (2012). Optimal decentralized protocol for electric vehicle charging. \textit{IEEE Transactions on Power Systems}, 28(2), 940-951.

\bibitem{senjyu2005}
Senjyu, T., Kuninaka, R., Urasaki, N., Fujita, H., \& Funabashi, T. (2005, December). Power system stabilization based on robust centralized and decentralized controllers. In \textit{2005 International Power Engineering Conference} (pp. 905-910). IEEE.

\bibitem{mg_prop}
Nikkhajoei, H., \& Lasseter, R. H. (2007, June). Microgrid protection. In \textit{2007 IEEE Power Engineering Society General Meeting} (pp. 1-6). IEEE.

\bibitem{deng2014}
Deng, R., Yang, Z., Hou, F., Chow, M. Y., \& Chen, J. (2015). Distributed real-time demand response in multiseller–multibuyer smart distribution grid. \textit{IEEE Transactions on Power Systems}, 30(5), 2364-2374.

\bibitem{dist_1}
Rahbari-Asr, N., Ojha, U., Zhang, Z., \& Chow, M. Y. (2014). Incremental welfare consensus algorithm for cooperative distributed generation/demand response in smart grid. \textit{IEEE Transactions on Smart Grid}, 5(6), 2836-2845.

\bibitem{dist_2}
Mudumbai, R., Dasgupta, S., \& Cho, B. B. (2012). Distributed control for optimal economic dispatch of a network of heterogeneous power generators. \textit{IEEE Transactions on Power Systems}, 27(4), 1750-1760.

\bibitem{mitsubishi}
MLU Series Technical Specifications Sheet. \textit{Mitsubishi Electric Photovoltaic Modules}. \url{https://www.mitsubishielectricsolar.com/images/uploads/documents/specs/MLU_spec_sheet_250W_255W.pdf}. Accessed June 23, 2019.

\bibitem{kaabeche2011b}
Kaabeche, A., Belhamel, M., \& Ibtiouen, R. (2011). Sizing optimization of grid-independent hybrid photovoltaic/wind power generation system. \textit{Energy}, 36(2), 1214-1222.

\bibitem{raum}
3.5 kW Wind Turbine System Specification Sheet. \textit{Raum Energy}. \url{https://farm-energy.extension.org/wp-content/uploads/2019/04/3.5kW_Spec_Sheet.pdf}. Accessed June 20, 2019.

\bibitem{seeling1997}
Seeling-Hochmuth, G. C. (1997). A combined optimization concept for the design and operation strategy of hybrid-PV energy systems. \textit{Solar Energy}, 61(2), 77-87.

\bibitem{kaviani2009}
Kaviani, A. K., Riahy, G. H., \& Kouhsari, S. M. (2009). Optimal design of a reliable hydrogen-based stand-alone wind/PV generating system, considering component outages. \textit{Renewable energy}, 34(11), 2380-2390.

\bibitem{wen2009}
Wen, J., Zheng, Y., \& Donghan, F. (2009). A review on reliability assessment for wind power. \textit{Renewable and sustainable energy reviews}, 13(9), 2485-2494.

\bibitem{tahmo}
African Climate Data. \textit{Trans-African HydroMeteorological Observatory}.
\url{https://tahmo.org/__african-climate-data/}. Accessed June 30, 2019.

\bibitem{accuweather}
Weather forecast for Timbila, Kenya. \textit{AccuWeather}.
\url{https://www.accuweather.com/en/ke/timbila/827492/air-travel-hourly-forecast/827492?hour=34}. Accessed July 1, 2019.

\bibitem{ghasemi2013}
Ghasemi, A., Asrari, A., Zarif, M., \& Abdelwahed, S. (2013). Techno-economic analysis of stand-alone hybrid photovoltaic–diesel–battery systems for rural electrification in eastern part of Iran—A step toward sustainable rural development. \textit{Renewable and Sustainable Energy Reviews}, 28, 456-462.

\bibitem{chp_mt}
Combined Heat and Power Technology Fact Sheet Series: Microturbines. \textit{U.S. Department of Energy}. \url{https://www.energy.gov/sites/prod/files/2016/09/f33/CHP-Microturbines_0.pdf}. Accessed July 3, 2019. 

\bibitem{bs_source}
DiOrio, N., Dobos, A., \& Janzou, S. (2015). Economic analysis case studies of battery energy storage with SAM (No. NREL/TP-6A20-64987). \textit{National Renewable Energy Lab.(NREL)}, Golden, CO (United States).

\bibitem{powerwall}
Powerwall technical specifications. \textit{Tesla}. \url{https://www.tesla.com/en_GB/powerwall}. Accessed June 27, 2019. 

\bibitem{ogunjuyigbe2016}
Ogunjuyigbe, A. S. O., Ayodele, T. R., \& Akinola, O. A. (2016). Optimal allocation and sizing of PV/Wind/Split-diesel/Battery hybrid energy system for minimizing life cycle cost, carbon emission and dump energy of remote residential building. Applied Energy, 171, 153-171.

\bibitem{costs_irena}
Renewable Power Generation Costs in 2018. \textit{International Renewable Energy Agency (IRENA)}. \url{https://www.irena.org/publications/2019/May/Renewable-power-generation-costs-in-2018}. Accessed June 16, 2019.

\bibitem{hu2017}
Hu, X., Zou, C., Zhang, C., \& Li, Y. (2017). Technological developments in batteries: a survey of principal roles, types, and management needs. \textit{IEEE Power and Energy Magazine}, 15(5), 20-31.

\bibitem{lcc}
Brown, R. J., \& Yanuck, R. R. (1985). Introduction to life-cycle costing. \textit{Fairmont Press}, Englewood Cliffs, NJ.

\bibitem{ajan2003}
Ajan, C. W., Ahmed, S. S., Ahmad, H. B., Taha, F., \& Zin, A. A. B. M. (2003). On the policy of photovoltaic and diesel generation mix for an off-grid site: East Malaysian perspectives. \textit{Solar Energy}, 74(6), 453-467.

\bibitem{prasad2006}
Prasad, A. R., \& Natarajan, E. (2006). Optimization of integrated photovoltaic–wind power generation systems with battery storage. \textit{Energy}, 31(12), 1943-1954.

\bibitem{emissions}
Potential to Emit Calculator for Stationary Compression Ignition Engines (2016). \textit{U.S. Environmental Protection Agency (EPA)}.
\url{https://www.epa.gov/sites/production/files/2016-07/engines_ci_pte_calculator_version_1.0_0.xlsx}. Accessed June 16, 2019. 

\bibitem{ng_price}
Natural gas price. \textit{Business Insider}. \url{https://markets.businessinsider.com/commodities/natural-gas-price}. Accessed July 1, 2019. 

\bibitem{nerc}
Standard INT-003-0 - Interchange Transaction Implementation. \textit{North American Electric Reliability Corporation}.
\url{https://www.nerc.com/files/INT-003-0.pdf}

\bibitem{conv_price}
SolaX Power (2017). \textit{X-Hybrid Single/Three Phase Inverter Datasheet}. SolaX Power Co., Ltd. \url{https://www.solaxpower.com/wp-content/uploads/2017/02/X-Hybrid-SingleThree-Phase.pdf}. Accessed July 2019.

\bibitem{usea}
Msyani, C. M. (2013). Current status of energy sector in Tanzania. \textit{Presentation at the United States Energy Association Executive Exchange on Developing an Ancillary Service Market, Washington, DC}.

\bibitem{bnef}
Goldie-Scot, L. (2019). A Behind the scenes take on lithium-ion battery prices. \textit{Bloomberg New Energy Finance}.

\bibitem{diesel_price}
Diesel prices. \url{https://www.globalpetrolprices.com/diesel_prices/}. Accessed July 20, 2019. 

\bibitem{interest_rates}
Central banks - summary of current interest rates. \textit{Global-rates.com}. \url{https://www.global-rates.com/interest-rates/central-banks/central-banks.aspx}. Accessed July 20, 2019.

\bibitem{inflation_rate}
List of inflation rates by country. \textit{Trading Economics}.
\url{https://tradingeconomics.com/country-list/inflation-rate}. Accessed July 20, 2019.

\bibitem{interest_rate}
Kenya interest rate. \textit{Trading Economics}. \url{https://tradingeconomics.com/kenya/interest-rate}. Accessed June 30, 2019.

\bibitem{LCOEgrid}
Moksnes, N., Korkovelos, A., Mentis, D., \& Howells, M. (2017). Electrification pathways for Kenya–linking spatial electrification analysis and medium to long term energy planning. \textit{Environmental Research Letters}, 12(9), 095008.

\bibitem{Cext}
Zeyringer, M., Pachauri, S., Schmid, E., Schmidt, J., Worrell, E., \& Morawetz, U. B. (2015). Analyzing grid extension and stand-alone photovoltaic systems for the cost-effective electrification of Kenya. \textit{Energy for Sustainable Development}, 25, 75-86.

\bibitem{ESS_control}
Arani, A. K., Gharehpetian, G. B., \& Abedi, M. (2019). Review on energy storage systems control methods in microgrids. \textit{International Journal of Electrical Power \& Energy Systems}, 107, 745-757.

\bibitem{mpc}
Mayne, D. Q., Rawlings, J. B., Rao, C. V., \& Scokaert, P. O. (2000). Constrained model predictive control: Stability and optimality. \textit{Automatica}, 36(6), 789-814.

\bibitem{pereira2017}
Pereira, M., de la Peña, D. M., \& Limón, D. (2017). Robust economic model predictive control of a community micro-grid. \textit{Renewable Energy}, 100, 3-17.

\bibitem{prodan2014}
Prodan, I., \& Zio, E. (2014). A model predictive control framework for reliable microgrid energy management. \textit{International Journal of Electrical Power \& Energy Systems}, 61, 399-409.

\bibitem{zhu2015}
Zhu, B., Tazvinga, H., \& Xia, X. (2014). Switched model predictive control for energy dispatching of a photovoltaic-diesel-battery hybrid power system. IEEE Transactions on Control Systems Technology, 23(3), 1229-1236.

\bibitem{load_profile}
Prinsloo, G., Dobson, R., \& Brent, A. (2016). Scoping exercise to determine load profile archetype reference shapes for solar co-generation models in isolated off-grid rural African villages. \textit{Journal of Energy in Southern Africa}, 27(3), 11-27.

\bibitem{golumbeanu2013}
Golumbeanu, R., \& Barnes, D. (2013). Connection charges and electricity access
in Sub-Saharan Africa. \textit{World Bank Policy Research Working Paper}, 6511.

\bibitem{iea_africa2019}
International Energy Agency (2019). \textit{Africa Energy Outlook 2019: World
Energy Outlook Special Report}. IEA, Paris.
\url{https://www.iea.org/reports/africa-energy-outlook-2019}

\bibitem{tracy2012}
Tracy, J., \& Jacobson, A. (2012). \textit{The True Cost of Kerosene in Rural
Africa}. Lighting Africa Program, IFC/World Bank Group.
\url{https://www.lightingglobal.org/wp-content/uploads/2012/04/40_kerosene_pricing_Lighting_Africa_Report.pdf}


\bibitem{wang2024nsga}
Wang, X., Zhu, H., Luo, X., Chang, S., \& Guan, X. (2024). A novel optimal dispatch strategy for hybrid energy ship power system based on the improved NSGA-II algorithm. \textit{Electric Power Systems Research}, 232, 110385. \url{https://doi.org/10.1016/j.epsr.2024.110385}

\bibitem{zhu2020imogwo}
Zhu, W., Guo, J., Zhao, G., \& Zeng, B. (2020). Optimal sizing of an island hybrid microgrid based on improved multi-objective grey wolf optimizer. \textit{Processes}, 8(12), 1581. \url{https://doi.org/10.3390/pr8121581}

\bibitem{geleta2022gwo}
Geleta, D. K., Manshahia, M. S., Vasant, P., \& Banik, A. (2022). Grey wolf optimizer for optimal sizing of hybrid wind and solar renewable energy system. \textit{Computational Intelligence}, 38(3), 1133--1162. \url{https://doi.org/10.1111/coin.12349}

\bibitem{javidrad2017psosa}
Javidrad, F., \& Nazari, M. (2017). A new hybrid particle swarm and simulated annealing stochastic optimization method. \textit{Applied Soft Computing}, 60, 634--654. \url{https://doi.org/10.1016/j.asoc.2017.07.023}

\bibitem{muhsen2019mode}
Muhsen, D. H., Nabil, M., Haider, H. T., \& Khatib, T. (2019). A novel method for sizing of standalone photovoltaic system using multi-objective differential evolution algorithm and hybrid multi-criteria decision making methods. \textit{Energy}, 174, 1158--1175. \url{https://doi.org/10.1016/j.energy.2019.03.046}

\bibitem{karimi2019madm}
Karimi, H., \& Jadid, S. (2019). Optimal microgrid operation scheduling by a novel hybrid multi-objective and multi-attribute decision-making framework. \textit{Energy}, 186, 115912. \url{https://doi.org/10.1016/j.energy.2019.115912}

\bibitem{bakhtiari2021mcmc}
Bakhtiari, H., Zhong, J., \& Alvarez, M. (2021). Predicting the stochastic behavior of uncertainty sources in planning a stand-alone renewable energy-based microgrid using Metropolis--coupled Markov chain Monte Carlo simulation. \textit{Applied Energy}, 290, 116719. \url{https://doi.org/10.1016/j.apenergy.2021.116719}

\bibitem{calabrese2026tssp}
Calabrese, M., Ademollo, A., \& Carcasci, C. (2026). Designing off-grid hybrid renewable energy systems under uncertainty: A two-stage stochastic programming approach. \textit{Renewable Energy}, 256, 124193. \url{https://doi.org/10.1016/j.renene.2025.124193}

\bibitem{derrac2011tutorial}
Derrac, J., Garc\'{i}a, S., Molina, D., \& Herrera, F. (2011). A practical tutorial on the use of nonparametric statistical tests as a methodology for comparing evolutionary and swarm intelligence algorithms. \textit{Swarm and Evolutionary Computation}, 1(1), 3--18. \url{https://doi.org/10.1016/j.swevo.2011.02.002}

\bibitem{carrasco2020trends}
Carrasco, J., Garc\'{i}a, S., Rueda, M. M., Das, S., \& Herrera, F. (2020). Recent trends in the use of statistical tests for comparing swarm and evolutionary computing algorithms: Practical guidelines and a critical review. \textit{Swarm and Evolutionary Computation}, 54, 100665. \url{https://doi.org/10.1016/j.swevo.2020.100665}

\bibitem{hansen2021coco}
Hansen, N., Auger, A., Ros, R., Mersmann, O., Tu\v{s}ar, T., \& Brockhoff, D. (2021). COCO: a platform for comparing continuous optimizers in a black-box setting. \textit{Optimization Methods and Software}, 36(1), 114--144. \url{https://doi.org/10.1080/10556788.2020.1808977}

\bibitem{vargha2000a12}
Vargha, A., \& Delaney, H. D. (2000). A critique and improvement of the CL common language effect size statistics of McGraw and Wong. \textit{Journal of Educational and Behavioral Statistics}, 25(2), 101--132. \url{https://doi.org/10.3102/10769986025002101}

\bibitem{deb2011knee}
Deb, K., \& Gupta, S. (2011). Understanding knee points in bicriteria problems and their implications as preferred solution principles. \textit{Engineering Optimization}, 43(11), 1175--1204. \url{https://doi.org/10.1080/0305215X.2010.548863}

\bibitem{hwang1981topsis}
Hwang, C.-L., \& Yoon, K. (1981). \textit{Multiple Attribute Decision Making: Methods and Applications}. Lecture Notes in Economics and Mathematical Systems, Springer, Berlin. \url{https://doi.org/10.1007/978-3-642-48318-9}

\bibitem{gelaro2017merra2}
Gelaro, R., McCarty, W., Su\'{a}rez, M. J., Todling, R., Molod, A., Takacs, L., et al. (2017). The Modern-Era Retrospective Analysis for Research and Applications, Version 2 (MERRA-2). \textit{Journal of Climate}, 30(14), 5419--5454. \url{https://doi.org/10.1175/JCLI-D-16-0758.1}

\bibitem{nasapower}
NASA Langley Research Center. Prediction Of Worldwide Energy Resources (POWER) Project, hourly API, data for 2018 at 3.39\degree S, 37.72\degree E. \url{https://power.larc.nasa.gov/}. Accessed 19 October 2026.

\end{thebibliography}

\end{document}
